%---------------------------------------------------------
\section{Overview}
\label{sec:theory_overview}

This chapter develops, with full mathematical rigor, the theoretical
machinery underlying the functional principal component analysis (FPCA)
pipeline implemented in Chapter~\ref{chap:methods}: smoothing of discretely
observed trajectories, estimation of the mean function and covariance
operator, spectral decomposition of the covariance operator, and asymptotic
validation of the resulting estimators. The exposition follows
\cite{hsingeubank2015}, whose treatment of functional data analysis through
the lens of Hilbert space and linear operator theory provides, at present,
the most complete rigorous account of this pipeline. Definitions, theorem
statements, and proofs throughout this chapter are our own exposition of
that material, adapted to the notation and scope of this thesis, with
\cite{hsingeubank2015} cited at the relevant points.

The chapter is organized to mirror the logical dependency of the FPCA
pipeline rather than the chronological order in which a reader might first
encounter the topics. Section~\ref{sec:preliminaries} fixes notation and
recalls, without proof, the elementary Hilbert space and bounded-operator
concepts (inner products, orthonormal bases, the projection theorem,
reproducing kernel Hilbert spaces, Sobolev spaces, adjoint operators) that
everything else is built on. Section~\ref{sec:spectral_theory} develops the
spectral theory of compact operators, culminating in Mercer's Theorem, which
identifies the covariance kernel of a stochastic process with the spectral
decomposition of its covariance operator. Section~\ref{sec:perturbation}
develops perturbation bounds for the eigenvalues and eigenfunctions of a
compact self-adjoint operator, which are the technical tool later used to
transfer consistency of an estimated covariance operator into consistency
of its estimated eigenelements. Section~\ref{sec:smoothing} treats
smoothing and regularization, in particular smoothing splines, which
underlie the adaptive B-spline basis used in Chapter~\ref{chap:methods}.
Section~\ref{sec:random_elements} develops the probabilistic foundations of
FDA: the mean element and covariance operator of a Hilbert-space-valued
random element, and the Karhunen--Lo\`eve Theorem, which is the population
version of FPCA. Section~\ref{sec:mean_cov_estimation} studies the large
sample behavior of estimators of the mean function and covariance operator,
and Section~\ref{sec:fpca_theory} closes the chapter with the asymptotic
theory of the FPCA estimators themselves --- consistency and asymptotic
normality of the estimated eigenvalues and eigenfunctions --- which is the
theoretical justification for the FPCA results reported in
Chapter~\ref{chap:results}.

A scope note is in order. The source \cite{hsingeubank2015} develops this
entire pipeline for general basis systems and, concretely, for smoothing
splines (equivalently, B-spline bases with a roughness penalty), which
matches the adaptive B-spline smoothing used as the baseline in
Section~\ref{subsec:bspline}. It does not, however, treat the Fourier or
wavelet basis systems used as the comparison bases in this thesis. The
theoretical treatment of those two basis systems --- in particular, the
approximation-theoretic and minimax properties that motivate the use of
wavelets for spatially inhomogeneous (Besov-space) signals such as the ones
studied in Chapter~\ref{chap:results} --- will be added from a second
source in a subsequent revision of this chapter.

Throughout, $\mathbb{H}$, $\mathbb{H}_1$, $\mathbb{H}_2$ denote real Hilbert
spaces with inner products $\langle\cdot,\cdot\rangle$ (or
$\langle\cdot,\cdot\rangle_1$, $\langle\cdot,\cdot\rangle_2$) and induced
norms $\|\cdot\|$. To avoid a notational clash with the number of retained
basis coefficients $K$ used in Chapter~\ref{chap:methods}, we denote the
covariance operator of a random element by $\Gamma$ and its kernel by
$\gamma(s,t)$, rather than the letter $K$ used for this purpose in
\cite{hsingeubank2015}.

\section{Mathematical Preliminaries}
\label{sec:preliminaries}

This section collects, without proof, the elementary facts about metric,
normed, Hilbert, and reproducing kernel Hilbert spaces, and about bounded
linear operators, that are used throughout the rest of the chapter. All
results in this section are standard; proofs may be found in
\cite{hsingeubank2015}, Chapters~2--3, and the classical references cited
therein (Luenberger, 1969; Rynne and Youngson, 2001).

\subsection{Metric, Normed, and Banach Spaces}

\begin{definition}
A \emph{metric} on a set $\mathbb{M}$ is a function $d:\mathbb{M}\times
\mathbb{M}\to\mathbb{R}$ satisfying, for all $x_1,x_2,x_3\in\mathbb{M}$: (i)
$d(x_1,x_2)\geq 0$ with equality iff $x_1=x_2$; (ii) $d(x_1,x_2) =
d(x_2,x_1)$; (iii) $d(x_1,x_3)\leq d(x_1,x_2)+d(x_2,x_3)$ (triangle
inequality). The pair $(\mathbb{M},d)$ is a \emph{metric space}. A sequence
$\{x_n\}\subset\mathbb{M}$ is \emph{Cauchy} if
$\sup_{m,n\geq N}d(x_m,x_n)\to 0$ as $N\to\infty$; $(\mathbb{M},d)$ is
\emph{complete} if every Cauchy sequence converges, and \emph{separable} if
it has a countable dense subset.
\end{definition}

\begin{definition}
A \emph{vector space} $\mathbb{V}$ over $\mathbb{R}$ is a set closed under
addition and scalar multiplication satisfying the usual axioms of
commutativity, associativity, distributivity, and the existence of a zero
element and additive inverses. A \emph{norm} on $\mathbb{V}$ is a function
$\|\cdot\|:\mathbb{V}\to\mathbb{R}$ such that $\|v\|\geq 0$ with equality
iff $v=0$, $\|av\| = |a|\|v\|$, and $\|v_1+v_2\|\leq\|v_1\|+\|v_2\|$, for
all $v,v_1,v_2\in\mathbb{V}$, $a\in\mathbb{R}$. Every norm induces a metric
$d(x,y):=\|x-y\|$. A \emph{Banach space} is a normed vector space that is
complete under this metric.
\end{definition}

\begin{definition}
\label{def:Lp}
Let $(E,\mathcal{B},\mu)$ be a measure space and $p\in[1,\infty)$. The space
$\mathbb{L}^p(E,\mathcal{B},\mu)$ consists of (equivalence classes, under
equality $\mu$-a.e., of) measurable functions $f$ on $E$ with $\int_E
|f|^p\,d\mu<\infty$, equipped with the norm $\|f\|_p =
\big(\int_E|f|^p\,d\mu\big)^{1/p}$. When $E=[0,1]$, $\mathcal{B}$ is the
Borel $\sigma$-field, and $\mu$ is Lebesgue measure, we write
$\mathbb{L}^p[0,1]$. Each $\mathbb{L}^p(E,\mathcal{B},\mu)$, $p\geq 1$, is a
Banach space (the Riesz--Fischer Theorem).
\end{definition}

\subsection{Hilbert Spaces and Orthonormal Bases}

\begin{definition}
An \emph{inner product} on a vector space $\mathbb{V}$ is a function
$\langle\cdot,\cdot\rangle:\mathbb{V}\times\mathbb{V}\to\mathbb{R}$ that is
symmetric, bilinear, and satisfies $\langle v,v\rangle\geq 0$ with equality
iff $v=0$. It induces the norm $\|v\| = \langle v,v\rangle^{1/2}$, for which
the Cauchy--Schwarz inequality $|\langle v_1,v_2\rangle|\leq
\|v_1\|\|v_2\|$ holds. A \emph{Hilbert space} is a complete inner-product
space. The space $\mathbb{L}^2(E,\mathcal{B},\mu)$ of Definition
\ref{def:Lp} is a Hilbert space with $\langle f_1,f_2\rangle =
\int_E f_1f_2\,d\mu$; this is the setting most relevant for FDA, with $E =
[0,1]$ after rescaling of the observation window $\mathcal{T}$.
\end{definition}

\begin{definition}
Elements $x_1,x_2$ of a Hilbert space are \emph{orthogonal} if
$\langle x_1,x_2\rangle = 0$. A countable collection $\{e_j\}$ is an
\emph{orthonormal sequence} if $\|e_j\|=1$ for all $j$ and the $e_j$ are
pairwise orthogonal. An orthonormal sequence $\{e_j\}$ in a Hilbert space
$\mathbb{H}$ is an \emph{orthonormal basis}, or complete orthonormal system
(CONS), if $\overline{\mathrm{span}}\{e_j\} = \mathbb{H}$, where
$\overline{\mathrm{span}}$ denotes closed linear span; equivalently, if
$\langle x,e_j\rangle = 0$ for all $j$ forces $x=0$.
\end{definition}

\begin{theorem}[Fourier expansion / Parseval's relation]
\label{thm:parseval}
If $\{e_j\}$ is a CONS for a Hilbert space $\mathbb{H}$, every
$x\in\mathbb{H}$ satisfies $x = \sum_{j=1}^\infty \langle x,e_j\rangle e_j$
and $\|x\|^2 = \sum_{j=1}^\infty \langle x,e_j\rangle^2$. A Hilbert space is
separable if and only if it admits a countable orthonormal basis; every
infinite-dimensional separable Hilbert space is congruent (isometrically
isomorphic) to $\ell^2$, the space of square-summable real sequences.
\end{theorem}

\subsection{The Projection Theorem}

\begin{theorem}[Projection Theorem]
\label{thm:projection}
Let $\mathbb{M}$ be a closed linear subspace of a Hilbert space
$\mathbb{H}$. For every $x\in\mathbb{H}$ there is a unique element
$\hat{x}\in\mathbb{M}$, called the \emph{projection} of $x$ onto
$\mathbb{M}$, minimizing $\|x-y\|$ over $y\in\mathbb{M}$; it is
characterized by $\langle x-\hat{x},y\rangle = 0$ for all $y\in\mathbb{M}$.
Consequently, writing $\mathbb{M}^\perp := \{x\in\mathbb{H}:
\langle x,y\rangle=0 \text{ for all } y\in\mathbb{M}\}$ for the orthogonal
complement of $\mathbb{M}$, every $x\in\mathbb{H}$ decomposes uniquely as
$x = x_1+x_2$ with $x_1\in\mathbb{M}$, $x_2\in\mathbb{M}^\perp$, written
$\mathbb{H} = \mathbb{M}\oplus\mathbb{M}^\perp$.
\end{theorem}

\subsection{Reproducing Kernel Hilbert Spaces}
\label{subsec:rkhs_prelim}

\begin{definition}
\label{def:rkhs}
Let $\mathbb{H}$ be a Hilbert space of real-valued functions on a set $E$.
A bivariate function $K$ on $E\times E$ is a \emph{reproducing kernel} for
$\mathbb{H}$ if (i) $K(\cdot,t)\in\mathbb{H}$ for every $t\in E$, and (ii)
$f(t) = \langle f,K(\cdot,t)\rangle$ for every $f\in\mathbb{H}$, $t\in E$
(the reproducing property). A Hilbert space possessing a reproducing kernel
is a \emph{reproducing kernel Hilbert space} (RKHS), denoted
$\mathbb{H}(K)$. The evaluation functional $\mathcal{T}_t:f\mapsto f(t)$
has representer $K(\cdot,t)$, and $\mathbb{H}$ is an RKHS if and only if
every evaluation functional $\mathcal{T}_t$, $t\in E$, is bounded.
\end{definition}

A reproducing kernel is necessarily unique, symmetric, and nonnegative
definite, i.e.\ $\sum_{i=1}^n\sum_{j=1}^n a_ia_jK(t_i,t_j)\geq 0$ for any
$n$, $\{a_i\}\subset\mathbb{R}$, $\{t_i\}\subset E$. Conversely, the
Moore--Aronszajn Theorem guarantees that every symmetric, nonnegative
definite $K$ on $E\times E$ is the reproducing kernel of a unique RKHS
$\mathbb{H}(K)$. If $K_1,K_2$ are reproducing kernels for RKHSs
$\mathbb{H}(K_1),\mathbb{H}(K_2)$ of functions on $E$ with
$\mathbb{H}(K_1)\cap\mathbb{H}(K_2)=\{0\}$, then $K=K_1+K_2$ is the
reproducing kernel of $\mathbb{H}(K)=\mathbb{H}(K_1)\oplus\mathbb{H}(K_2)$
(the two summand spaces are then orthogonal subspaces of $\mathbb{H}(K)$).
RKHS norm convergence implies pointwise (and, under a
uniform bound on $K$, uniform) convergence, which is what distinguishes
RKHSs from general $\mathbb{L}^2$-type function spaces and makes them
convenient for smoothing problems (Section~\ref{sec:smoothing}) and for
representing covariance kernels (Section~\ref{sec:random_elements}).

\subsection{Sobolev Spaces}

\begin{definition}
\label{def:sobolev}
For an integer $q\geq 1$, the \emph{Sobolev space} $\mathbb{W}^q[0,1]$
consists of functions $f$ on $[0,1]$ that are $(q-1)$-times differentiable
with $f^{(q-1)}$ absolutely continuous and $f^{(q)}\in\mathbb{L}^2[0,1]$.
Equipped with the inner product
\begin{equation}
\langle f,g\rangle'_{\mathbb{W}^q[0,1]} := \langle f,g\rangle_2 +
\langle f^{(q)},g^{(q)}\rangle_2,
\label{eq:sobolev_ip}
\end{equation}
where $\langle\cdot,\cdot\rangle_2$ is the $\mathbb{L}^2[0,1]$ inner
product, $\mathbb{W}^q[0,1]$ is a reproducing kernel Hilbert space.
\end{definition}

The Sobolev space $\mathbb{W}^q[0,1]$ formalizes the notion of a function
with bounded, uniformly distributed roughness of order $q$: it is precisely
the class of functions whose smoothing-spline estimator is studied in
Section~\ref{sec:smoothing}, and the assumption that GDP growth
trajectories $X_i\in\mathbb{W}^q[0,1]$ underlies the smoothing-spline
justification for the B-spline baseline of Section~\ref{subsec:bspline}.

\subsection{Bounded Linear Operators and Adjoints}

\begin{definition}
Let $\mathbb{X}_1,\mathbb{X}_2$ be normed vector spaces. A linear
transformation $\mathcal{T}:\mathbb{X}_1\to\mathbb{X}_2$ is \emph{bounded}
if there is a finite $C>0$ with $\|\mathcal{T}x\|_2\leq C\|x\|_1$ for all
$x\in\mathbb{X}_1$; boundedness is equivalent to (uniform) continuity of
$\mathcal{T}$. The set $\mathcal{B}(\mathbb{X}_1,\mathbb{X}_2)$ of bounded
linear operators from $\mathbb{X}_1$ to $\mathbb{X}_2$ is itself a normed
space under the \emph{operator norm}
$\|\mathcal{T}\| := \sup_{\|x\|_1=1}\|\mathcal{T}x\|_2$, and is a Banach
space whenever $\mathbb{X}_2$ is. We write $\mathcal{B}(\mathbb{X})$ for
$\mathcal{B}(\mathbb{X},\mathbb{X})$.
\end{definition}

\begin{theorem}[Adjoint operator]
\label{thm:adjoint}
Let $\mathbb{H}_1,\mathbb{H}_2$ be Hilbert spaces. For every
$\mathcal{T}\in\mathcal{B}(\mathbb{H}_1,\mathbb{H}_2)$ there is a unique
$\mathcal{T}^*\in\mathcal{B}(\mathbb{H}_2,\mathbb{H}_1)$, the
\emph{adjoint} of $\mathcal{T}$, satisfying
\begin{equation}
\langle\mathcal{T}x_1,x_2\rangle_2 = \langle x_1,\mathcal{T}^*x_2\rangle_1,
\qquad x_1\in\mathbb{H}_1,\ x_2\in\mathbb{H}_2.
\label{eq:adjoint}
\end{equation}
If $\mathbb{H}_1=\mathbb{H}_2=\mathbb{H}$, $\mathcal{T}$ is
\emph{self-adjoint} if $\mathcal{T}^*=\mathcal{T}$. The adjoint satisfies
$(\mathcal{T}^*)^*=\mathcal{T}$, $\|\mathcal{T}^*\|=\|\mathcal{T}\|$, and
$\|\mathcal{T}^*\mathcal{T}\|=\|\mathcal{T}\|^2$.
\end{theorem}

\begin{definition}
\label{def:nonneg_operator}
A self-adjoint operator $\mathcal{T}\in\mathcal{B}(\mathbb{H})$ is
\emph{nonnegative definite} if $\langle\mathcal{T}x,x\rangle\geq 0$ for all
$x\in\mathbb{H}$. For any $\mathcal{T}\in\mathcal{B}(\mathbb{H}_1,
\mathbb{H}_2)$, the operator $\mathcal{T}^*\mathcal{T}$ is nonnegative
definite on $\mathbb{H}_1$, since it is self-adjoint and
$\langle\mathcal{T}^*\mathcal{T}x,x\rangle = \|\mathcal{T}x\|^2\geq 0$. This
fact is the operator-theoretic reason that covariance operators
(Section~\ref{sec:random_elements}), which have precisely this
$\mathcal{T}^*\mathcal{T}$-type structure, are always nonnegative definite.
\end{definition}

\section{Spectral Theory of Compact Operators}
\label{sec:spectral_theory}

This section develops the eigenvalue/eigenvector decomposition of a compact
self-adjoint operator, its non-self-adjoint (singular value) counterpart,
and Mercer's Theorem, which identifies the spectral decomposition of an
integral operator with a series expansion of its kernel. This machinery is
what will let us pass, in Section~\ref{sec:random_elements}, from a
covariance kernel $\gamma(s,t)$ to the eigenvalues/eigenfunctions that
define FPCA. Throughout, all Hilbert spaces are assumed separable, which
entails no loss of generality for compact operators (Theorem
\ref{thm:compact_finite_approx} below). Proofs follow
\cite{hsingeubank2015}, Chapter~4.

\subsection{Compact Operators}

\begin{definition}
\label{def:compact}
A linear transformation $\mathcal{T}$ from a normed space $\mathbb{X}_1$
into a normed space $\mathbb{X}_2$ is \emph{compact} if for every bounded
sequence $\{x_n\}\subset\mathbb{X}_1$, $\{\mathcal{T}x_n\}$ contains a
subsequence that converges in $\mathbb{X}_2$.
\end{definition}

A compact linear transformation is necessarily bounded: if $\mathcal{T}$
were unbounded there would be a bounded sequence $\{x_n\}$ with
$\|\mathcal{T}x_n\|_2\geq n$, so $\{\mathcal{T}x_n\}$ could have no
convergent subsequence. Compactness is a genuine restriction, however.

\begin{theorem}
\label{thm:identity_not_compact}
The identity operator on an infinite-dimensional normed space is not
compact.
\end{theorem}

\begin{proof}
Consider first a Hilbert space $\mathbb{H}$ with CONS $\{e_j\}$. Since
$\|e_i-e_j\| = \sqrt{2}$ for $i\neq j$, the bounded sequence $\{e_j\}$ has
no Cauchy (hence no convergent) subsequence, so $\{Ie_j\}$ has none either.

For a general infinite-dimensional normed space $\mathbb{X}$, it suffices
to construct a bounded sequence $\{e_j\}$ with
$\inf_{i\neq j}\|e_i-e_j\|>0$; the same argument then applies. Let $e_1$
have unit norm and $\mathbb{Y}_1=\mathrm{span}\{e_1\}$. Pick any
$x\notin\mathbb{Y}_1$ and set $d=\inf\{\|x-y\|:y\in\mathbb{Y}_1\}>0$
(finite since $\mathbb{Y}_1$ is closed, being finite-dimensional). For any
$\alpha\in(1,\infty)$ choose $z\in\mathbb{Y}_1$ with $\|x-z\|<\alpha d$ and
set $e_2 = (x-z)/\|x-z\|$. For any $y\in\mathbb{Y}_1$,
\[
\|e_2-y\| = \frac{\|x-(z+\|x-z\|y)\|}{\|x-z\|} > \frac{d}{\alpha d} =
\alpha^{-1},
\]
since $z+\|x-z\|y\in\mathbb{Y}_1$; in particular $\|e_1-e_2\|>\alpha^{-1}$.
Repeating this construction with $\mathbb{Y}_2=\mathrm{span}\{e_1,e_2\}$,
and inductively thereafter, produces $\{e_j\}$ with $\|e_i-e_j\|>\alpha^{-1}$
for all $i\neq j$.
\end{proof}

\begin{theorem}
\label{thm:compact_props}
Let $\mathcal{T}$ be a compact operator between normed spaces
$\mathbb{X}_1,\mathbb{X}_2$. Then: (1) $\overline{\mathrm{Im}(\mathcal{T})}$
is separable; (2) any finite-rank bounded operator is compact; (3) the
composition of two operators is compact if either factor is compact; (4)
if $\mathbb{X}_2$ is a Banach space, the set of compact operators in
$\mathcal{B}(\mathbb{X}_1,\mathbb{X}_2)$ is closed under the operator norm.
\end{theorem}

\begin{proof}
(1) Let $B(0;r) = \{x\in\mathbb{X}_1:\|x\|_1\leq r\}$ and $\{y_n\}\subset
\mathcal{T}(B(0;r))$. Choose $x_n\in B(0;r)$ with $\|y_n-\mathcal{T}x_n\|_2
<n^{-1}$. Compactness of $\mathcal{T}$ gives a convergent subsequence of
$\{\mathcal{T}x_n\}$, hence of $\{y_n\}$, so $\mathcal{T}(B(0;r))$ is
sequentially compact, hence compact and separable. Since
$\mathrm{Im}(\mathcal{T})\subset\bigcup_{r=1}^\infty \mathcal{T}(B(0;r))$,
$\mathrm{Im}(\mathcal{T})$, and thus its closure, is separable.

(2) If $\mathrm{Im}(\mathcal{T})$ is finite-dimensional and $\{x_n\}$ is
bounded in $\mathbb{X}_1$, the Bolzano--Weierstrass Theorem gives a
convergent subsequence of $\{\mathcal{T}x_n\}$.

(3) Let $\mathcal{T}_1\in\mathcal{B}(\mathbb{X}_1,\mathbb{X}_2)$,
$\mathcal{T}_2\in\mathcal{B}(\mathbb{X}_2,\mathbb{X}_3)$ and $\{x_n\}$
bounded in $\mathbb{X}_1$. If $\mathcal{T}_1$ is compact, $\{\mathcal{T}_1
x_n\}$ has a convergent subsequence, whose image under the continuous map
$\mathcal{T}_2$ also converges. If instead $\mathcal{T}_2$ is compact, then
since $\{\mathcal{T}_1x_n\}$ is bounded, $\{\mathcal{T}_2\mathcal{T}_1x_n\}$
has a convergent subsequence.

(4) Let $\{\mathcal{T}_n\}$ be compact with $\|\mathcal{T}_n-\mathcal{T}\|
\to 0$. Let $\{x_k\}$ be bounded in $\mathbb{X}_1$. By a diagonal argument,
extract nested subsequences $\{x_{nk}\}_k\subset\{x_{(n-1)k}\}_k$ (with
$\{x_{0k}\}:=\{x_k\}$) such that $\mathcal{T}_nx_{nk}\to y_n$ as
$k\to\infty$, for each $n$. Choose $k_n$ increasing in $n$ with
$\|\mathcal{T}_nx_{nk}-y_n\|_2<n^{-1}$ for $k\geq k_n$. For $n'>n$,
\[
y_n-y_{n'} = (y_n-\mathcal{T}_nx_{n'k_{n'}}) +
(\mathcal{T}_{n'}x_{n'k_{n'}}-y_{n'}) +
(\mathcal{T}_n-\mathcal{T})x_{n'k_{n'}} +
(\mathcal{T}-\mathcal{T}_{n'})x_{n'k_{n'}},
\]
so $\|y_n-y_{n'}\|_2 \leq n^{-1}+n'^{-1}+\|\mathcal{T}_n-\mathcal{T}\|+
\|\mathcal{T}_{n'}-\mathcal{T}\|$, which shows $\{y_n\}$ is Cauchy, hence
convergent to some $y$ (as $\mathbb{X}_2$ is complete). Finally
\[
\|\mathcal{T}x_{nk_n}-y\|_2 \leq \|\mathcal{T}-\mathcal{T}_n\|\|x_{nk_n}\|_1
+ \|\mathcal{T}_nx_{nk_n}-y_n\|_2 + \|y_n-y\|_2 \to 0,
\]
so $\{\mathcal{T}x_{nk_n}\}$ converges and $\mathcal{T}$ is compact.
\end{proof}

\begin{theorem}
\label{thm:compact_finite_approx}
Let $\mathbb{H}_1,\mathbb{H}_2$ be Hilbert spaces and
$\mathcal{T}\in\mathcal{B}(\mathbb{H}_1,\mathbb{H}_2)$. Then $\mathcal{T}$
is compact if and only if there exist finite-rank operators $\mathcal{T}_n$
with $\|\mathcal{T}_n-\mathcal{T}\|\to 0$; and $\mathcal{T}$ is compact if
and only if $\mathcal{T}^*$ is compact.
\end{theorem}

\begin{proof}
Sufficiency of the finite-rank approximation follows from parts (2) and (4)
of Theorem~\ref{thm:compact_props}. For necessity, part (1) of Theorem
\ref{thm:compact_props} gives a CONS $\{e_j\}$ for
$\overline{\mathrm{Im}(\mathcal{T})}$; define $\mathcal{T}_nx :=
\sum_{j=1}^n\langle\mathcal{T}x,e_j\rangle_2 e_j = \mathcal{P}_n\mathcal{T}x$
with $\mathcal{P}_n$ the projection onto $\mathrm{span}\{e_1,\dots,e_n\}$.
If $\|\mathcal{T}_n-\mathcal{T}\|\not\to 0$, there is $\epsilon>0$ and
unit-norm $\{x_{n'}\}$ with $\|(\mathcal{T}-\mathcal{T}_{n'})x_{n'}\|_2
>\epsilon$. By compactness $\mathcal{T}x_{n''}\to y$ along a subsequence, and
\[
(\mathcal{T}-\mathcal{T}_{n''})x_{n''} = (I-\mathcal{P}_{n''})y +
(I-\mathcal{P}_{n''})(\mathcal{T}x_{n''}-y) \to 0,
\]
a contradiction; hence $\|\mathcal{T}_n-\mathcal{T}\|\to 0$.

For the adjoint claim: $\mathcal{T}_n^*$ is finite-rank (Theorem
\ref{thm:adjoint}) with $\|\mathcal{T}_n^*-\mathcal{T}^*\| =
\|\mathcal{T}_n-\mathcal{T}\|$, so if $\mathcal{T}$ is compact, so is
$\mathcal{T}^*$ by the first part of the theorem; apply this twice, using
$(\mathcal{T}^*)^*=\mathcal{T}$, for the converse.
\end{proof}

As a consequence of Theorem~\ref{thm:compact_finite_approx} and part (1) of
Theorem~\ref{thm:compact_props}, for a compact
$\mathcal{T}\in\mathcal{B}(\mathbb{H}_1,\mathbb{H}_2)$,
$\mathrm{Im}(\mathcal{T}^*)$ is separable, and since $\mathbb{H}_1 =
\mathrm{Ker}(\mathcal{T})\oplus\mathrm{Im}(\mathcal{T}^*)$ (Theorem
\ref{thm:adjoint_properties} below), there is no loss of generality in
assuming both $\mathbb{H}_1,\mathbb{H}_2$ separable, an assumption we
maintain throughout.

\subsection{Spectral Decomposition of Compact Self-Adjoint Operators}
\label{subsec:spectral_decomp}

\begin{theorem}
\label{thm:adjoint_properties}
Let $\mathcal{T}\in\mathcal{B}(\mathbb{H}_1,\mathbb{H}_2)$. Then: (1)
$\mathrm{Ker}(\mathcal{T}) = (\mathrm{Im}(\mathcal{T}^*))^\perp$; (2)
$\mathrm{Ker}(\mathcal{T}^*\mathcal{T})=\mathrm{Ker}(\mathcal{T})$ and
$\mathrm{Im}(\mathcal{T}^*\mathcal{T})=\mathrm{Im}(\mathcal{T}^*)$; (3)
$\mathbb{H}_1 = \mathrm{Ker}(\mathcal{T})\oplus\mathrm{Im}(\mathcal{T}^*)$.
\end{theorem}

\begin{proof}
(1) If $x_1\in\mathrm{Ker}(\mathcal{T})$ then $\langle x_1,\mathcal{T}^*
x_2\rangle_1=\langle\mathcal{T}x_1,x_2\rangle_2=0$ for all $x_2$, so
$x_1\in(\mathrm{Im}(\mathcal{T}^*))^\perp$. Conversely if
$x_1\in(\mathrm{Im}(\mathcal{T}^*))^\perp$ then $\|\mathcal{T}x_1\|_2^2 =
\langle x_1,\mathcal{T}^*\mathcal{T}x_1\rangle_1=0$ since
$\mathcal{T}^*\mathcal{T}x_1\in\mathrm{Im}(\mathcal{T}^*)$.

(2) If $x_1\in\mathrm{Ker}(\mathcal{T})$ then trivially
$\mathcal{T}^*\mathcal{T}x_1=0$. Conversely if
$\mathcal{T}^*\mathcal{T}x_1=0$ then $0 =
\langle\mathcal{T}^*\mathcal{T}x_1,x_1\rangle_1 = \|\mathcal{T}x_1\|_2^2$,
so $x_1\in\mathrm{Ker}(\mathcal{T})$; this proves the kernel identity, and
the image identity follows from it together with part (1) and Theorem
\ref{thm:projection}, part 3.

(3) By Theorem~\ref{thm:projection}, $\mathbb{H}_1 = \mathrm{Ker}
(\mathcal{T})\oplus(\mathrm{Ker}(\mathcal{T}))^\perp = \mathrm{Ker}
(\mathcal{T})\oplus\mathrm{Im}(\mathcal{T}^*)$, using part (1).
\end{proof}

\begin{definition}
\label{def:eigen}
For $\mathcal{T}\in\mathcal{B}(\mathbb{H})$, a scalar $\lambda\in\mathbb{R}$
and nonzero $e\in\mathbb{H}$ with $\mathcal{T}e=\lambda e$ are, respectively,
an \emph{eigenvalue} and corresponding \emph{eigenvector} (or
\emph{eigenfunction}, when $\mathbb{H}$ is a function space) of
$\mathcal{T}$. $\mathrm{Ker}(\mathcal{T}-\lambda I)$ is the
\emph{eigenspace} of $\lambda$.
\end{definition}

\begin{theorem}
\label{thm:eigen_orthogonal}
Let $\mathcal{T}\in\mathcal{B}(\mathbb{H})$ and $e_j\in\mathrm{Ker}
(\mathcal{T}-\lambda_jI)$, $j=1,2,\dots$, with the $\lambda_j$ distinct and
nonzero. Then: (1) the $e_j$ are linearly independent; (2) if $\mathcal{T}$
is self-adjoint, the $e_j$ are mutually orthogonal.
\end{theorem}

\begin{proof}
(1) By induction on $n$. For $n=2$: if $c_1e_1+c_2e_2=0$ then also
$\lambda_1c_1e_1+\lambda_2c_2e_2=0$, so $\lambda_1(c_1e_1+c_2e_2) -
(\lambda_1c_1e_1+\lambda_2c_2e_2) = (\lambda_1-\lambda_2)c_2e_2=0$, forcing
$c_2=0$ (as $\lambda_1\neq\lambda_2$) and then $c_1=0$. Assume the claim for
$n=k$ and suppose $\sum_{j=1}^{k+1}c_je_j=0$. Then
$\lambda_{k+1}\sum_{j=1}^{k+1}c_je_j - \sum_{j=1}^{k+1}\lambda_jc_je_j =
\sum_{j=1}^k(\lambda_{k+1}-\lambda_j)c_je_j = 0$, so $c_j=0$ for $j\leq k$
by the induction hypothesis, and then $c_{k+1}=0$.

(2) $\langle e_i,e_j\rangle = \langle e_i,\lambda_j^{-1}\mathcal{T}e_j
\rangle = \lambda_j^{-1}\langle\mathcal{T}e_i,e_j\rangle =
\lambda_j^{-1}\lambda_i\langle e_i,e_j\rangle$, which forces
$\langle e_i,e_j\rangle=0$ since $\lambda_i\neq\lambda_j$.
\end{proof}

\begin{theorem}
\label{thm:compact_eigen_finite}
Let $\mathcal{T}\in\mathcal{B}(\mathbb{H})$ be compact. Then: (1)
$\mathrm{Ker}(\mathcal{T}-\lambda I)$ is finite-dimensional for every
$\lambda\neq 0$; (2) for every $\epsilon>0$, only finitely many distinct
eigenvalues of $\mathcal{T}$ satisfy $|\lambda|\geq\epsilon$; (3) the set of
nonzero eigenvalues of $\mathcal{T}$ is countable.
\end{theorem}

\begin{proof}
(1) If $\mathrm{Ker}(\mathcal{T}-\lambda I)$ were infinite-dimensional for
some $\lambda\neq 0$, the construction of Theorem
\ref{thm:identity_not_compact} yields unit-norm $\{e_j\}\subset\mathrm{Ker}
(\mathcal{T}-\lambda I)$ with $\inf_{i\neq j}\|e_i-e_j\|>0$, so
$\inf_{i\neq j}\|\mathcal{T}e_i-\mathcal{T}e_j\| = |\lambda|\inf_{i\neq j}
\|e_i-e_j\|>0$, contradicting compactness.

(2) Let $\lambda_1,\lambda_2,\dots$ be distinct eigenvalues with
$|\lambda_j|\geq\lambda>0$, $e_j\in\mathrm{Ker}(\mathcal{T}-\lambda_jI)$,
$\mathbb{M}_n=\mathrm{span}\{e_1,\dots,e_n\}$, $\mathbb{M}_0:=\{0\}$. By
part (1) of Theorem~\ref{thm:eigen_orthogonal} the $e_j$ are linearly
independent, so $\mathbb{M}_n\cap\mathbb{M}_{n-1}^\perp$ is
one-dimensional; Gram--Schmidt produces an orthonormal sequence $\{x_n\}$
with $x_n\in\mathrm{Ker}(\mathcal{T}-\lambda_nI)\cap\mathbb{M}_{n-1}^\perp$,
and for $n>m$, $\|\mathcal{T}x_n-\mathcal{T}x_m\|^2 = \|\mathcal{T}x_n\|^2+
\|\mathcal{T}x_m\|^2 \geq 2\lambda^2$, again contradicting compactness if
there were infinitely many such $\lambda_j$.

(3) Immediate from (2), taking $\epsilon\downarrow 0$ along a sequence.
\end{proof}

We now state and prove the central result of this section: every compact
self-adjoint operator admits an eigenvalue/eigenvector decomposition that
is the direct infinite-dimensional analog of the spectral decomposition of
a symmetric positive semidefinite matrix.

\begin{theorem}[Spectral Theorem for compact self-adjoint operators]
\label{thm:spectral_theorem}
Let $\mathcal{T}$ be a compact, self-adjoint operator on $\mathbb{H}$. The
nonzero eigenvalues of $\mathcal{T}$ form a finite set or a sequence
tending to $0$, each with finite-dimensional eigenspace, and eigenvectors
for distinct eigenvalues are orthogonal. Order the eigenvalues
$|\lambda_1|\geq|\lambda_2|\geq\cdots$ with corresponding orthonormal
eigenvectors $e_1,e_2,\dots$ (obtained via Gram--Schmidt within repeated
eigenvalues). Then $\{e_j\}$ is a CONS for $\overline{\mathrm{Im}
(\mathcal{T})}$ and
\begin{equation}
\mathcal{T} = \sum_{j\geq 1}\lambda_j\, e_j\otimes e_j, \qquad\text{i.e.}
\qquad \mathcal{T}x = \sum_{j\geq 1}\lambda_j\langle x,e_j\rangle e_j,
\quad x\in\mathbb{H}.
\label{eq:spectral_decomp}
\end{equation}
\end{theorem}

\begin{proof}
Finiteness/decay of the eigenvalue sequence, finite multiplicity, and
orthogonality of eigenvectors are Theorems~\ref{thm:eigen_orthogonal} and
\ref{thm:compact_eigen_finite}. It remains to prove
\eqref{eq:spectral_decomp}. Since $\mathcal{T}$ is self-adjoint, Theorem
\ref{thm:adjoint_properties}, part (3) with $\mathcal{T}^*=\mathcal{T}$
gives $\mathbb{H} = \mathrm{Ker}(\mathcal{T})\oplus\mathrm{Im}(\mathcal{T})$,
and \eqref{eq:spectral_decomp} clearly holds for
$x\in\mathrm{span}\{e_j:j\geq 1\}$ and for $x\in\mathrm{Ker}(\mathcal{T})$.
So it suffices to show
\begin{equation}
\mathrm{Im}(\mathcal{T}) = \overline{\mathrm{span}}\{e_j:j\geq 1\}.
\label{eq:im_span}
\end{equation}
For any finite $n$ and $c_1,\dots,c_n\in\mathbb{R}$,
$\sum_{j=1}^n c_je_j = \mathcal{T}\big(\sum_{j=1}^n\lambda_j^{-1}c_je_j
\big)\in\mathrm{Im}(\mathcal{T})$, so $\overline{\mathrm{span}}\{e_j:j\geq
1\}\subset\overline{\mathrm{Im}(\mathcal{T})}=\mathrm{Im}(\mathcal{T})$
(the closure equality using Theorem~\ref{thm:compact_props}(1) and
completeness of finite-rank projections, as in Theorem
\ref{thm:compact_finite_approx}). Write, by Theorem~\ref{thm:projection},
$\mathrm{Im}(\mathcal{T}) = \overline{\mathrm{span}}\{e_j:j\geq
1\}\oplus\mathbb{N}$ for $\mathbb{N}$ the orthogonal complement of
$\overline{\mathrm{span}}\{e_j\}$ within $\mathrm{Im}(\mathcal{T})$. For
$x\in\mathbb{N}$, $y\in\overline{\mathrm{span}}\{e_j\}$, self-adjointness
gives $\langle\mathcal{T}x,y\rangle=\langle x,\mathcal{T}y\rangle=0$, so
$\mathcal{T}$ maps $\mathbb{N}$ into $\mathbb{N}$. Let $\mathcal{T}_
{\mathbb{N}}$ denote this restriction. Any nonzero eigenvalue of
$\mathcal{T}_\mathbb{N}$ is an eigenvalue of $\mathcal{T}$; but every
nonzero eigenvalue of $\mathcal{T}$ already appears among $\{\lambda_j\}$
with eigenvector in $\mathrm{span}\{e_j\}$, so $\mathcal{T}_\mathbb{N}$ can
have no nonzero eigenvalue. Since $\mathcal{T}_\mathbb{N}$ is a compact
self-adjoint operator on $\mathbb{N}$, if it were not the zero operator
then $\pm\|\mathcal{T}_\mathbb{N}\|$ would be one of its eigenvalues
(Theorem~\ref{thm:eigenvalue_max} below), a contradiction; hence
$\mathcal{T}_\mathbb{N}=0$ and $\mathbb{N}\subset\mathrm{Im}(\mathcal{T})
\cap\mathrm{Ker}(\mathcal{T}) = \{0\}$, proving \eqref{eq:im_span}.
\end{proof}

Every eigenvalue of a nonnegative definite operator is nonnegative, since
$\langle\mathcal{T}e,e\rangle=\lambda\|e\|^2$; conversely, by Theorem
\ref{thm:spectral_theorem}, a compact self-adjoint $\mathcal{T}$ is
nonnegative definite if and only if all its eigenvalues are nonnegative.
This is the property that will identify covariance operators
(Section~\ref{sec:random_elements}) as having entirely nonnegative spectra.

\begin{theorem}
\label{thm:eigenvalue_max}
If $\mathcal{T}$ is compact and nonnegative definite with eigenpairs
$\{(\lambda_j,e_j)\}$ ordered $\lambda_1\geq\lambda_2\geq\cdots\geq 0$, then
for every $k\geq 1$,
\begin{equation}
\lambda_k = \max_{e\in\mathrm{span}\{e_1,\dots,e_{k-1}\}^\perp}
\frac{\langle\mathcal{T}e,e\rangle}{\|e\|^2},
\label{eq:rayleigh_max}
\end{equation}
with $\mathrm{span}\{e_1,\dots,e_{k-1}\}^\perp=\mathbb{H}$ when $k=1$; more
generally, for a compact self-adjoint $\mathcal{T}$ with eigenpairs ordered
$\lambda_1^2\geq\lambda_2^2\geq\cdots$,
\begin{equation}
|\lambda_k| = \max_{e\in\mathrm{span}\{e_1,\dots,e_{k-1}\}^\perp}
\frac{\|\mathcal{T}e\|}{\|e\|},
\qquad\text{so that}\qquad
\|\mathcal{T}\| = \max_j|\lambda_j|.
\label{eq:norm_eigenvalue}
\end{equation}
\end{theorem}

\begin{proof}
For $e=\sum_{j\geq k}a_je_j\in\mathrm{span}\{e_1,\dots,e_{k-1}\}^\perp$,
$\langle\mathcal{T}e,e\rangle/\|e\|^2 = \sum_{j\geq k}\lambda_ja_j^2\big/
\sum_{j\geq k}a_j^2 \leq \lambda_k$ (as $\lambda_j\leq\lambda_k$ for
$j\geq k$), with equality at $e=e_k$; this proves
\eqref{eq:rayleigh_max}. For \eqref{eq:norm_eigenvalue}, apply
\eqref{eq:rayleigh_max} to the nonnegative definite compact operator
$\mathcal{T}^2$ (whose eigenpairs are $(\lambda_j^2,e_j)$, since $\mathcal{T}
^2=\sum_j\lambda_j^2e_j\otimes e_j$ by \eqref{eq:spectral_decomp} applied
twice) to get $\lambda_k^2 = \max_{e\perp\mathrm{span}\{e_1,\dots,e_{k-1}\}}
\|\mathcal{T}e\|^2/\|e\|^2$, and take square roots.
\end{proof}

\begin{theorem}[Courant--Fischer minimax principle]
\label{thm:courant_fischer}
Let $\mathcal{T}$ be a nonnegative, compact operator on $\mathbb{H}$ with
eigenvalues $\lambda_1\geq\lambda_2\geq\cdots\geq 0$ and eigenvectors
$\{e_j\}$. Then, for every $k\geq 1$,
\begin{equation}
\lambda_k = \max_{v_1,\dots,v_k\in\mathbb{H}}\
\min_{v\in\mathrm{span}\{v_1,\dots,v_k\}}\frac{\langle\mathcal{T}v,v
\rangle}{\|v\|^2}
= \min_{v_1,\dots,v_{k-1}\in\mathbb{H}}\
\max_{v\in\mathrm{span}\{v_1,\dots,v_{k-1}\}^\perp}\frac{\langle\mathcal{T}
v,v\rangle}{\|v\|^2},
\label{eq:courant_fischer}
\end{equation}
with both extrema attained at $v_k=e_k$.
\end{theorem}

\begin{proof}
Let $\mathbb{M}_{k-1}=\mathrm{span}\{e_1,\dots,e_{k-1}\}$. For the second
equality in \eqref{eq:courant_fischer}: given any $(k-1)$-dimensional
$\mathbb{V}_{k-1}\subset\mathbb{H}$, the subspace $\mathbb{V}_{k-1}^\perp
\cap\mathbb{M}_k$ (with $\mathbb{M}_k=\mathrm{span}\{e_1,\dots,e_k\}$) is
nontrivial; for $v=\sum_{j=1}^k a_je_j$ in this intersection,
$\langle\mathcal{T}v,v\rangle/\|v\|^2 = \sum_{j=1}^k\lambda_ja_j^2\big/
\sum_{j=1}^k a_j^2\geq\lambda_k$, so
$\max_{v\in\mathbb{V}_{k-1}^\perp}\langle\mathcal{T}v,v\rangle/\|v\|^2
\geq\lambda_k$ for every $\mathbb{V}_{k-1}$, with equality at
$\mathbb{V}_{k-1}=\mathbb{M}_{k-1}$ by Theorem~\ref{thm:eigenvalue_max}.

For the first equality: given any $k$-dimensional $\mathbb{V}_k\subset
\mathbb{H}$, the subspace $\mathbb{V}_k\cap\mathbb{M}_{k-1}^\perp$ is
nontrivial; for $v=\sum_{j\geq k}a_je_j$ in this intersection,
$\langle\mathcal{T}v,v\rangle/\|v\|^2\leq\lambda_k$ (as in the proof of
Theorem~\ref{thm:eigenvalue_max}), so $\min_{v\in\mathbb{V}_k}\langle
\mathcal{T}v,v\rangle/\|v\|^2 \leq \min_{v\in\mathbb{V}_k\cap
\mathbb{M}_{k-1}^\perp}\langle\mathcal{T}v,v\rangle/\|v\|^2\leq\lambda_k$
for every $\mathbb{V}_k$, with equality at $\mathbb{V}_k=\mathbb{M}_k$.
\end{proof}

The Courant--Fischer principle yields an elementary but essential
perturbation bound: the eigenvalues of a compact self-adjoint operator are
a $1$-Lipschitz function of the operator itself, in operator norm. This is
the first (and simplest) instance of the perturbation theory that will be
extended to eigenvectors in Section~\ref{sec:perturbation} and applied to
the estimated covariance operator in Section~\ref{sec:fpca_theory}.

\begin{theorem}
\label{thm:eigenvalue_perturbation}
Let $\mathcal{T},\widetilde{\mathcal{T}}$ be nonnegative definite compact
operators on $\mathbb{H}$ with eigenvalue sequences $\{\lambda_j\}$,
$\{\widetilde\lambda_j\}$. Then
\begin{equation}
\sup_{k\geq 0}|\lambda_{k+1}-\widetilde\lambda_{k+1}| \leq
\|\mathcal{T}-\widetilde{\mathcal{T}}\|.
\label{eq:eigenvalue_perturbation}
\end{equation}
\end{theorem}

\begin{proof}
By Theorem~\ref{thm:courant_fischer}, for any $v_1,\dots,v_k\in\mathbb{H}$,
\[
\max_{v\in\mathrm{span}\{v_1,\dots,v_k\}^\perp}\frac{\langle\mathcal{T}v,v
\rangle}{\|v\|^2} \leq
\max_{v\in\mathrm{span}\{v_1,\dots,v_k\}^\perp}\frac{\langle\widetilde{
\mathcal{T}}v,v\rangle}{\|v\|^2} +
\max_{v\in\mathbb{H}}\frac{\langle(\mathcal{T}-\widetilde{\mathcal{T}})v,v
\rangle}{\|v\|^2},
\]
and the last term is bounded by $\|\mathcal{T}-\widetilde{\mathcal{T}}\|$.
Taking the minimum over $v_1,\dots,v_k$ on both sides gives
$\lambda_{k+1}\leq\widetilde\lambda_{k+1}+\|\mathcal{T}-
\widetilde{\mathcal{T}}\|$; exchanging the roles of $\mathcal{T}$ and
$\widetilde{\mathcal{T}}$ gives the reverse inequality.
\end{proof}

\subsection{The Singular Value Decomposition}
\label{subsec:svd}

For a compact operator $\mathcal{T}\in\mathcal{B}(\mathbb{H}_1,\mathbb{H}_2)$
between two (possibly distinct) Hilbert spaces, self-adjointness is not
available, but a closely related expansion --- the singular value
decomposition (SVD) --- still holds, obtained by applying the Spectral
Theorem to the nonnegative definite compact operators $\mathcal{T}^*
\mathcal{T}$ and $\mathcal{T}\mathcal{T}^*$.

\begin{theorem}[Singular Value Decomposition]
\label{thm:svd}
Let $\mathcal{T}\in\mathcal{B}(\mathbb{H}_1,\mathbb{H}_2)$ be compact. Let
$\{(\lambda_j^2,f_{1j})\}$ be the eigenpairs of the compact, nonnegative
definite operator $\mathcal{T}^*\mathcal{T}$ (Theorem
\ref{thm:spectral_theorem}), with $\lambda_1\geq\lambda_2\geq\cdots\geq 0$,
and set $f_{2j}:=\mathcal{T}f_{1j}/\lambda_j$ for $\lambda_j>0$. Then
$\{(\lambda_j^2,f_{2j})\}$ are eigenpairs of $\mathcal{T}\mathcal{T}^*$,
$\{f_{2j}\}$ is orthonormal, $\mathcal{T}^*f_{2j}=\lambda_jf_{1j}$, and
\begin{equation}
\mathcal{T} = \sum_{j\geq 1}\lambda_j\,(f_{1j}\otimes_1 f_{2j}), \qquad
\text{i.e.}\qquad \mathcal{T}x = \sum_{j\geq 1}\lambda_j\langle
x,f_{1j}\rangle_1 f_{2j}, \quad x\in\mathbb{H}_1.
\label{eq:op_svd}
\end{equation}
The triples $(\lambda_j,f_{1j},f_{2j})$ are the \emph{singular system} of
$\mathcal{T}$; the $\lambda_j$ are its \emph{singular values}.
\end{theorem}

\begin{proof}
For $\lambda_j>0$, $(\mathcal{T}\mathcal{T}^*)(\mathcal{T}f_{1j}) =
\mathcal{T}(\mathcal{T}^*\mathcal{T}f_{1j}) = \lambda_j^2\mathcal{T}f_{1j}$,
so $f_{2j}=\mathcal{T}f_{1j}/\lambda_j$ is a unit-norm eigenvector of
$\mathcal{T}\mathcal{T}^*$ for eigenvalue $\lambda_j^2$; orthonormality of
$\{f_{2j}\}$ follows from Theorem~\ref{thm:eigen_orthogonal}(2) applied to
$\mathcal{T}\mathcal{T}^*$. And $\mathcal{T}^*f_{2j} = \mathcal{T}^*
\mathcal{T}f_{1j}/\lambda_j = \lambda_jf_{1j}$.

For \eqref{eq:op_svd}: by Theorem~\ref{thm:adjoint_properties} and Theorem
\ref{thm:projection}, part 3 of Theorem~\ref{thm:projection}(applied twice)
gives $(\mathrm{Ker}(\mathcal{T}))^\perp = (\mathrm{Ker}(\mathcal{T}^*
\mathcal{T}))^\perp = (\mathrm{Im}(\mathcal{T}^*\mathcal{T}))^{\perp\perp} =
\overline{\mathrm{Im}(\mathcal{T}^*\mathcal{T})}$, which by
\eqref{eq:im_span} (applied to $\mathcal{T}^*\mathcal{T}$) equals
$\overline{\mathrm{span}}\{f_{1j}:j\geq 1\}$. So it suffices to verify
\eqref{eq:op_svd} for $x\in\overline{\mathrm{span}}\{f_{1j}\}$, where it holds
by Theorem~\ref{thm:parseval} and the definition of $f_{2j}$, and for
$x\in\mathrm{Ker}(\mathcal{T})$, where both sides vanish.
\end{proof}

\begin{theorem}
\label{thm:svd_operator_norm}
If $\mathcal{T}\in\mathcal{B}(\mathbb{H}_1,\mathbb{H}_2)$ is compact with
singular values $\lambda_1\geq\lambda_2\geq\cdots\geq 0$, then
$\|\mathcal{T}\|=\lambda_1$, and, more generally, for every $k$,
\begin{equation}
\lambda_k = \max_{v_1,\dots,v_k\in\mathbb{H}_1}\
\min_{f\in\mathrm{span}\{v_1,\dots,v_k\}}\frac{\|\mathcal{T}f\|_2}{\|f\|_1}
=\min_{v_1,\dots,v_{k-1}\in\mathbb{H}_1}\
\max_{f\in\mathrm{span}\{v_1,\dots,v_{k-1}\}^\perp}\frac{\|\mathcal{T}f\|_2
}{\|f\|_1}.
\label{eq:svd_minimax}
\end{equation}
\end{theorem}

\begin{proof}
Immediate from Theorems~\ref{thm:eigenvalue_max} and
\ref{thm:courant_fischer} applied to $\mathcal{T}^*\mathcal{T}$, since
$\|\mathcal{T}f\|_2^2/\|f\|_1^2 = \langle\mathcal{T}^*\mathcal{T}f,f
\rangle_1/\|f\|_1^2$ and $\mathcal{T}^*\mathcal{T}$ has eigenvalues
$\lambda_j^2$.
\end{proof}

\begin{theorem}
\label{thm:compact_iff_svd}
An operator $\mathcal{T}\in\mathcal{B}(\mathbb{H}_1,\mathbb{H}_2)$ is
compact if and only if it admits the expansion \eqref{eq:op_svd}.
\end{theorem}

\begin{proof}
Necessity is Theorem~\ref{thm:svd}. For sufficiency, if \eqref{eq:op_svd}
holds, let $\mathcal{T}_n=\sum_{j=1}^n\lambda_j(f_{1j}\otimes_1f_{2j})$,
which is finite-rank, hence compact. By Theorem
\ref{thm:svd_operator_norm} applied to $\mathcal{T}-\mathcal{T}_n$ (whose
singular values are $\lambda_{n+1},\lambda_{n+2},\dots$),
$\|\mathcal{T}-\mathcal{T}_n\|=\lambda_{n+1}\to 0$, so $\mathcal{T}$ is
compact by Theorem~\ref{thm:compact_props}(4).
\end{proof}

\begin{theorem}
\label{thm:im_not_closed}
If $\mathcal{T}\in\mathcal{B}(\mathbb{H}_1,\mathbb{H}_2)$ is a compact
operator with infinitely many nonzero singular values, then
$\mathrm{Im}(\mathcal{T})$ is not closed.
\end{theorem}

\begin{proof}
Let $\{(\lambda_j,f_{1j},f_{2j})\}$ be the singular system, $\lambda_j
\downarrow 0$. Choose a subsequence $\{j_k\}$ with $\lambda_{j_k}\leq
k^{-2}$ and set $y=\sum_{k\geq 1}\lambda_{j_k}f_{2j_k}\in\mathbb{H}_2$
(convergent since $\sum_k\lambda_{j_k}^2<\infty$). Each partial sum lies in
$\mathrm{Im}(\mathcal{T})$, since $f_{2j_k}=\mathcal{T}(f_{1j_k}/
\lambda_{j_k})$, but $\sum_k\langle y,f_{2j_k}\rangle_2^2/\lambda_{j_k}^2 =
\sum_k 1 = \infty$ shows $y$ fails the Picard condition
$\sum_j\langle y,f_{2j}\rangle_2^2/\lambda_j^2<\infty$ that characterizes
$\mathrm{Im}(\mathcal{T})$ (via $\mathcal{T}^\dagger y :=
\sum_j\lambda_j^{-1}\langle y,f_{2j}\rangle_2 f_{1j}$, the Moore--Penrose
inverse, which converges precisely when this sum is finite). Hence
$y\notin\mathrm{Im}(\mathcal{T})$, even though $y$ is a limit of elements
of $\mathrm{Im}(\mathcal{T})$.
\end{proof}

This last fact --- that a compact operator with infinite rank never has
closed range --- is the root cause of the ill-posedness that pervades
inverse problems in FDA (e.g.\ functional linear regression); it is not
needed for the FPCA development that follows, but is recorded here because
it is an immediate and instructive consequence of the SVD.

\subsection{Hilbert--Schmidt and Trace-Class Operators}
\label{subsec:hs_trace}

Within the class of compact operators, two subclasses defined by summability
of the singular values are of particular importance for FDA: Hilbert--Schmidt
operators, whose squared singular values are summable, and trace-class
operators, whose singular values themselves are summable. Covariance
operators (Section~\ref{sec:random_elements}) will always be trace class,
hence Hilbert--Schmidt, hence compact.

\begin{theorem}
\label{thm:hs_cons_independence}
Let $\mathcal{T}_1,\mathcal{T}_2\in\mathcal{B}(\mathbb{H}_1,\mathbb{H}_2)$
and suppose that for some CONSs $\{e_{1i}\}$, $\{e_{2j}\}$ of
$\mathbb{H}_1,\mathbb{H}_2$, either $\sum_i(\|\mathcal{T}_1e_{1i}\|_2^2+
\|\mathcal{T}_2e_{1i}\|_2^2)<\infty$ or $\sum_j(\|\mathcal{T}_1^*e_{2j}
\|_1^2+\|\mathcal{T}_2^*e_{2j}\|_1^2)<\infty$. Then
\begin{equation}
\sum_i\langle\mathcal{T}_1e_{1i},\mathcal{T}_2e_{1i}\rangle_2 =
\sum_j\langle\mathcal{T}_1^*e_{2j},\mathcal{T}_2^*e_{2j}\rangle_1,
\label{eq:hs_polarization}
\end{equation}
and this common value does not depend on the choice of CONSs.
\end{theorem}

\begin{proof}
First take $\mathcal{T}_1=\mathcal{T}_2=\mathcal{T}$. Expanding
$\mathcal{T}e_{1i}=\sum_j\langle\mathcal{T}e_{1i},e_{2j}\rangle_2e_{2j}$ and
using Parseval's relation (Theorem~\ref{thm:parseval}) twice,
\[
\sum_i\|\mathcal{T}e_{1i}\|_2^2 = \sum_i\sum_j\langle\mathcal{T}e_{1i},
e_{2j}\rangle_2^2 = \sum_i\sum_j\langle e_{1i},\mathcal{T}^*e_{2j}
\rangle_1^2 = \sum_j\|\mathcal{T}^*e_{2j}\|_1^2,
\]
which shows the common double-sum value does not depend on $\{e_{1i}\}$,
and, by symmetry of the middle expression, not on $\{e_{2j}\}$ either. For
general $\mathcal{T}_1,\mathcal{T}_2$, apply this identity to
$\mathcal{T}_1+\mathcal{T}_2$, $\mathcal{T}_1$, and $\mathcal{T}_2$ and
subtract, using either summability hypothesis to justify the term-by-term
subtraction.
\end{proof}

\begin{definition}
\label{def:hs_operator}
$\mathcal{T}\in\mathcal{B}(\mathbb{H}_1,\mathbb{H}_2)$ is
\emph{Hilbert--Schmidt} (HS) if, for some (any) CONS $\{e_i\}$ of
$\mathbb{H}_1$, $\|\mathcal{T}\|_{HS}^2 := \sum_i\|\mathcal{T}e_i\|_2^2 <
\infty$. The HS inner product of $\mathcal{T}_1,\mathcal{T}_2\in
\mathcal{B}_{HS}(\mathbb{H}_1,\mathbb{H}_2)$ is $\langle\mathcal{T}_1,
\mathcal{T}_2\rangle_{HS} := \sum_i\langle\mathcal{T}_1e_i,\mathcal{T}_2e_i
\rangle_2$. By Theorem~\ref{thm:hs_cons_independence}, choosing $\{e_i\}$
to be the singular vectors $\{f_{1j}\}$ of a compact $\mathcal{T}$ gives
$\|\mathcal{T}\|_{HS}^2 = \sum_j\lambda_j^2$.
\end{definition}

\begin{theorem}
\label{thm:hs_compact}
A Hilbert--Schmidt operator is compact.
\end{theorem}

\begin{proof}
Let $\mathcal{T}$ be HS, $\{e_{2i}\}$ a CONS for $\mathbb{H}_2$, and
$\mathcal{T}_nx:=\sum_{i=1}^n\langle\mathcal{T}x,e_{2i}\rangle_2e_{2i}$,
which is finite-rank. For $\|x\|_1\leq 1$, by Cauchy--Schwarz,
\[
\|(\mathcal{T}_n-\mathcal{T})x\|_2^2 = \sum_{i=n+1}^\infty\langle
\mathcal{T}x,e_{2i}\rangle_2^2 = \sum_{i=n+1}^\infty\langle x,\mathcal{T}^*
e_{2i}\rangle_1^2 \leq \sum_{i=n+1}^\infty\|\mathcal{T}^*e_{2i}\|_1^2 \to 0,
\]
where the tail sum $\to 0$ since $\sum_i\|\mathcal{T}^*e_{2i}\|_1^2 =
\|\mathcal{T}\|_{HS}^2<\infty$ (Theorem~\ref{thm:hs_cons_independence}).
So $\|\mathcal{T}_n-\mathcal{T}\|\to 0$ and $\mathcal{T}$ is compact by
Theorem~\ref{thm:compact_props}(4).
\end{proof}

\begin{theorem}
\label{thm:hs_hilbert_space}
$\mathcal{B}_{HS}(\mathbb{H}_1,\mathbb{H}_2)$ is a separable Hilbert space
under $\langle\cdot,\cdot\rangle_{HS}$, with CONS $\{e_{1i}\otimes_1
e_{2j}\}$ for any CONSs $\{e_{1i}\},\{e_{2j}\}$ of $\mathbb{H}_1,
\mathbb{H}_2$.
\end{theorem}

\begin{proof}
Writing $a_{ij}=\langle\mathcal{T}e_{1i},e_{2j}\rangle_2$, Theorem
\ref{thm:hs_cons_independence} gives $\|\mathcal{T}\|_{HS}^2 = \sum_i\sum_j
a_{ij}^2$, identifying $\mathcal{B}_{HS}(\mathbb{H}_1,\mathbb{H}_2)$
isometrically with $\ell^2(\mathbb{N}\times\mathbb{N})$; completeness
follows from that of $\ell^2$ (Theorem~\ref{thm:parseval}/Riesz--Fischer).
Orthonormality of $\{e_{1i}\otimes_1e_{2j}\}$ is immediate. If
$\langle\mathcal{T},e_{1i}\otimes_1e_{2j}\rangle_{HS}=\langle\mathcal{T}
e_{1i},e_{2j}\rangle_2=0$ for all $i,j$, then $\mathcal{T}=0$ by
\eqref{eq:hs_polarization}, so $\{e_{1i}\otimes_1e_{2j}\}$ is a CONS by
Theorem~\ref{thm:parseval}.
\end{proof}

The next result shows that the rank-$k$ truncation of the SVD is the best
possible rank-$k$ approximation to a Hilbert--Schmidt operator, in HS norm
--- the operator-theoretic form of the classical Eckart--Young theorem.
This is the result that will justify truncating an estimated FPCA expansion
to its leading $k$ components as the optimal $k$-term summary of the
estimated covariance structure (cf.\ the variance-explained criterion of
Section~\ref{sec:metrics}).

\begin{theorem}[Schmidt--Mirsky / Eckart--Young]
\label{thm:eckart_young}
Let $\mathcal{T}\in\mathcal{B}_{HS}(\mathbb{H}_1,\mathbb{H}_2)$ have
singular system $\{(\lambda_j,f_{1j},f_{2j})\}$. Then for any finite $k$
and any $x_1,\dots,x_k\in\mathbb{H}_1$, $y_1,\dots,y_k\in\mathbb{H}_2$,
\begin{equation}
\Big\|\mathcal{T}-\sum_{j=1}^kx_j\otimes_1y_j\Big\|_{HS} \geq
\Big\|\mathcal{T}-\sum_{j=1}^k\lambda_jf_{1j}\otimes_1f_{2j}\Big\|_{HS} =
\Big(\sum_{j>k}\lambda_j^2\Big)^{1/2}.
\label{eq:eckart_young}
\end{equation}
\end{theorem}

\begin{proof}
It suffices to show $\|\mathcal{T}-\sum_{j=1}^kx_j\otimes_1y_j\|_{HS}^2
\geq \|\mathcal{T}\|_{HS}^2 - \sum_{j=1}^k\lambda_j^2$. Without loss the
$(y_j,x_j)$ may be taken orthonormal (an orthonormal reparametrization of
$\mathrm{span}\{x_j\}$, $\mathrm{span}\{y_j\}$ can only decrease the norm
of the difference). Extending $\{y_j\}_{j=1}^k$ to a CONS $\{y_j\}$ of
$\mathbb{H}_1$ and writing the HS norm via Definition
\ref{def:hs_operator},
\[
\Big\|\mathcal{T}-\sum_{i=1}^kx_i\otimes_1y_i\Big\|_{HS}^2 =
\sum_j\Big\|\Big(\mathcal{T}-\sum_{i=1}^kx_i\otimes_1y_i\Big)y_j\Big\|_2^2
= \sum_{j=1}^k\|\mathcal{T}y_j-x_j\|_2^2 + \sum_{j>k}\|\mathcal{T}y_j\|_2^2.
\]
Since $\sum_{j>k}\|\mathcal{T}y_j\|_2^2 = \|\mathcal{T}\|_{HS}^2 -
\sum_{j=1}^k\|\mathcal{T}y_j\|_2^2$, it remains to show $\sum_{j=1}^k
\|\mathcal{T}y_j\|_2^2\leq\sum_{j=1}^k\lambda_j^2$. Expanding
$\mathcal{T}y_j=\sum_{l}\lambda_l\langle f_{1l},y_j\rangle_1f_{2l}$ gives
$\|\mathcal{T}y_j\|_2^2 = \sum_l\lambda_l^2\langle f_{1l},y_j\rangle_1^2$,
and, using that $\lambda_l^2$ is nonincreasing in $l$ together with
Parseval's relation applied to the orthonormal $\{y_j\}_{j=1}^k$ (so that
$\sum_{j=1}^k\langle f_{1l},y_j\rangle_1^2\leq 1$ for every $l$, and
$\sum_l\sum_{j=1}^k\langle f_{1l},y_j\rangle_1^2 = k$),
\[
\sum_{j=1}^k\|\mathcal{T}y_j\|_2^2 = \sum_{l}\lambda_l^2\sum_{j=1}^k
\langle f_{1l},y_j\rangle_1^2 \leq \sum_{l=1}^k\lambda_l^2 +
\lambda_k^2\Big(k - \sum_{l=1}^k\sum_{j=1}^k\langle f_{1l},y_j
\rangle_1^2\Big) \leq \sum_{l=1}^k\lambda_l^2,
\]
where the first inequality bounds $\lambda_l^2\leq\lambda_k^2$ for $l>k$ and
adds and subtracts $\lambda_k^2\sum_{l=1}^k(\cdot)$, and the second uses
$\sum_{l=1}^k\sum_{j=1}^k\langle f_{1l},y_j\rangle_1^2\leq k$. Equality
throughout at $x_j=\lambda_jf_{2j}$, $y_j=f_{1j}$ gives
$\|\mathcal{T}-\sum_{j=1}^k\lambda_jf_{1j}\otimes_1f_{2j}\|_{HS}^2 =
\sum_{j>k}\lambda_j^2$.
\end{proof}

\begin{definition}
\label{def:trace_class}
$\mathcal{T}\in\mathcal{B}(\mathbb{H}_1,\mathbb{H}_2)$ is \emph{trace
class} if $\|\mathcal{T}\|_{TR} := \sum_j\langle(\mathcal{T}^*
\mathcal{T})^{1/2}e_j,e_j\rangle_1 <\infty$ for some (any) CONS $\{e_j\}$ of
$\mathbb{H}_1$, where $(\mathcal{T}^*\mathcal{T})^{1/2}$ is the nonnegative
self-adjoint square root furnished by Theorem~\ref{thm:spectral_theorem}
applied with fractional powers of the eigenvalues. Equivalently,
$\|\mathcal{T}\|_{TR}=\sum_j\lambda_j$, the sum of the singular values of
$\mathcal{T}$.
\end{definition}

\begin{theorem}
\label{thm:trace_is_hs}
A trace-class operator is Hilbert--Schmidt (hence compact).
\end{theorem}

\begin{proof}
If $\mathcal{T}$ is trace class then $(\mathcal{T}^*\mathcal{T})^{1/4}$ is
HS, since $\|\mathcal{T}\|_{TR}=\|(\mathcal{T}^*\mathcal{T})^{1/4}\|_{HS}^2$
(both sides equal $\sum_j\lambda_j$, comparing Definitions
\ref{def:hs_operator} and \ref{def:trace_class} via the eigenvalues
$\lambda_j^{1/2}$ of $(\mathcal{T}^*\mathcal{T})^{1/4}$), hence compact
(Theorem~\ref{thm:hs_compact}), so $\mathcal{T}^*\mathcal{T}=
((\mathcal{T}^*\mathcal{T})^{1/4})^4$ is compact (Theorem
\ref{thm:compact_props}(3)) with the same singular values $\lambda_j$ as
$\mathcal{T}$ (Theorem~\ref{thm:svd}). Then $\|\mathcal{T}\|_{HS}^2 =
\sum_j\lambda_j^2 \leq \lambda_1\sum_j\lambda_j = \lambda_1\|\mathcal{T}
\|_{TR}<\infty$.
\end{proof}

For a self-adjoint $\mathcal{T}$ with eigenvalues $\{\lambda_j\}$,
$\|\mathcal{T}\|_{TR}=\sum_j|\lambda_j|$, and if $\mathcal{T}$ is trace
class we call $\mathrm{trace}(\mathcal{T}):=\sum_j\lambda_j=\sum_j\langle
\mathcal{T}e_j,e_j\rangle$ (for any CONS $\{e_j\}$) its \emph{trace}; when
$\mathcal{T}$ is also nonnegative definite, $\mathrm{trace}(\mathcal{T}) =
\|\mathcal{T}\|_{TR}$. This is the quantity that Section~\ref{sec:metrics}
computes, for the empirical covariance operator, as total variance; the
percentage of variance explained by the leading eigenvalues is precisely
$\sum_{j\leq k}\lambda_j/\mathrm{trace}(\mathcal{T})$.

The following bound --- a trace-norm analog of Theorem
\ref{thm:eigenvalue_perturbation} for the whole singular value sequence at
once --- will be the key technical tool for transferring HS-norm
consistency of an estimated covariance operator into $\ell^2$ consistency
of its estimated eigenvalues in Section~\ref{sec:fpca_theory}.

\begin{theorem}
\label{thm:von_neumann}
Let $\mathcal{T},\widetilde{\mathcal{T}}\in\mathcal{B}_{HS}(\mathbb{H}_1,
\mathbb{H}_2)$ with nonascending singular values $\{\lambda_j\},
\{\widetilde\lambda_j\}$. Then
\begin{equation}
\sum_{j\geq 1}(\lambda_j-\widetilde\lambda_j)^2 \leq \mathrm{trace}\big\{
(\mathcal{T}-\widetilde{\mathcal{T}})^*(\mathcal{T}-\widetilde{\mathcal{T}})
\big\} = \|\mathcal{T}-\widetilde{\mathcal{T}}\|_{HS}^2.
\label{eq:von_neumann}
\end{equation}
\end{theorem}

\begin{proof}
Let $\{(\lambda_j,f_{1j},f_{2j})\}$, $\{(\widetilde\lambda_j,
\widetilde{f}_{1j},\widetilde{f}_{2j})\}$ be the singular systems of
$\mathcal{T},\widetilde{\mathcal{T}}$. Expanding $\widetilde{\mathcal{T}}
f_{1j} = \sum_k\widetilde\lambda_k\langle f_{1j},\widetilde{f}_{1k}
\rangle_1\widetilde{f}_{2k}$ and using Definition~\ref{def:hs_operator},
\[
\|\mathcal{T}-\widetilde{\mathcal{T}}\|_{HS}^2 = \sum_j\lambda_j^2 +
\sum_j\widetilde\lambda_j^2 - 2\sum_{j,k}\lambda_j\widetilde\lambda_k
\langle f_{1j},\widetilde{f}_{1k}\rangle_1\langle f_{2j},
\widetilde{f}_{2k}\rangle_2.
\]
By Cauchy--Schwarz and Parseval's relation, for every fixed $j$,
$\big|\sum_k\langle f_{1j},\widetilde{f}_{1k}\rangle_1\langle f_{2j},
\widetilde{f}_{2k}\rangle_2\big|\leq 1$, and likewise for every fixed $k$
summing over $j$; so, writing $\mathcal{A}=\{\langle f_{1j},
\widetilde{f}_{1k}\rangle_1\langle f_{2j},\widetilde{f}_{2k}\rangle_2\}
_{j,k}$, the row and column sums of $\mathcal{A}$ are bounded by $1$ in
absolute value. In the finite-rank case (rank $\leq n$ for both operators;
the general case follows by truncating and taking limits), writing
$\lambda=(\lambda_1,\dots,\lambda_n)^T$,
$\widetilde\lambda=(\widetilde\lambda_1,\dots,\widetilde\lambda_n)^T$,
\eqref{eq:von_neumann} is equivalent to $\lambda^T\mathcal{A}
\widetilde\lambda \leq \lambda^T\widetilde\lambda$. Since $\lambda,
\widetilde\lambda$ have nonincreasing nonnegative entries, each can be
written as a nonnegative combination of the vectors $\{e_i\}_{i=1}^n$
($e_i$ having its first $i$ entries equal to $1$, the rest $0$), so it
suffices to verify $e_k^T\mathcal{A}e_l \leq e_k^Te_l$ for all $k,l$, which
holds because $|e_k^T\mathcal{A}e_l| = \big|\sum_{j\leq k,i\leq l}
\mathcal{A}_{ji}\big|\leq\min(k,l)=e_k^Te_l$, using the row/column sum
bound.
\end{proof}

\subsection{Integral Operators and Mercer's Theorem}
\label{subsec:mercer}

We specialize now to the setting most relevant for FDA: $\mathbb{H} =
\mathbb{L}^2(E,\mathcal{B},\mu)$ for a finite measure $\mu$ on a compact
metric space $E$ (with $\mathcal{B}=\mathcal{B}(E)$ the Borel
$\sigma$-field and, without loss, $\mathrm{supp}(\mu)=E$), and operators
that arise as integral operators against a kernel. This is the setting in
which the covariance operator of a stochastic process will live
(Section~\ref{sec:random_elements}), and Mercer's Theorem is what makes the
covariance kernel and covariance operator two faces of the same object.

\begin{definition}
\label{def:integral_operator}
For a measurable $\kappa$ on $E\times E$ with $\iint_{E\times E}\kappa^2(s,
t)\,d\mu(s)d\mu(t)<\infty$, the associated \emph{integral operator}
$\mathcal{K}\in\mathcal{B}(\mathbb{L}^2(E,\mathcal{B},\mu))$ is
\begin{equation}
(\mathcal{K}f)(\cdot) := \int_E \kappa(s,\cdot)f(s)\,d\mu(s), \qquad
f\in\mathbb{L}^2(E,\mathcal{B},\mu),
\label{eq:integral_operator}
\end{equation}
with $\kappa$ its \emph{kernel}; by Cauchy--Schwarz,
$\|\mathcal{K}\|\leq\big(\iint_{E\times E}\kappa^2\,d\mu\,d\mu\big)^{1/2}$.
We assume throughout that $\kappa$ is continuous on $E\times E$ (hence
uniformly continuous, by compactness of $E$).
\end{definition}

\begin{lemma}
\label{lem:kf_continuous}
For every $f\in\mathbb{L}^2(E,\mathcal{B},\mu)$, $(\mathcal{K}f)(\cdot)$ is
uniformly continuous.
\end{lemma}

\begin{proof}
Given $\epsilon>0$, uniform continuity of $\kappa$ gives $\delta>0$ with
$|\kappa(s,t_2)-\kappa(s,t_1)|<\epsilon$ whenever $|t_2-t_1|<\delta$, for
all $s$. Then $|(\mathcal{K}f)(t_2)-(\mathcal{K}f)(t_1)|\leq\int_E|\kappa(s,
t_2)-\kappa(s,t_1)||f(s)|\,d\mu(s)\leq\epsilon\|f\|\mu(E)^{1/2}$ by
Cauchy--Schwarz.
\end{proof}

\begin{theorem}
\label{thm:integral_operator_compact}
The integral operator $\mathcal{K}$ of Definition
\ref{def:integral_operator} is compact.
\end{theorem}

\begin{proof}
By the Stone--Weierstrass Theorem, for any $\epsilon>0$ there are
continuous $g_1,\dots,g_n,h_1,\dots,h_n$ with $\sup_{s,t}|\kappa(s,t)-
\kappa_n(s,t)|<\epsilon$ for $\kappa_n(s,t)=\sum_{i=1}^ng_i(s)h_i(t)$. The
integral operator $\mathcal{K}_n$ with kernel $\kappa_n$ has range in
$\mathrm{span}\{h_1,\dots,h_n\}$, hence is finite-rank (compact), and by
Cauchy--Schwarz, $\|(\mathcal{K}_n-\mathcal{K})f\|^2 = \int_E\big(\int_E
[\kappa_n(s,t)-\kappa(s,t)]f(s)\,d\mu(s)\big)^2d\mu(t) \leq
\epsilon^2\|f\|^2\mu(E)^2$, so $\|\mathcal{K}_n-\mathcal{K}\|\to 0$ and
$\mathcal{K}$ is compact by Theorem~\ref{thm:compact_props}(4).
\end{proof}

If $\kappa$ is symmetric, $\mathcal{K}$ is self-adjoint (Example
\ref{ex:integral_operator_adjoint} below), and Theorem
\ref{thm:spectral_theorem} gives $\mathcal{K}=\sum_{j\geq 1}\lambda_je_j
\otimes e_j$; by Lemma~\ref{lem:kf_continuous}, the version of $e_j$ given
by $e_j(t)=\lambda_j^{-1}\int_E\kappa(s,t)e_j(s)\,d\mu(s)$ (for
$\lambda_j\neq 0$) is continuous, and we always adopt this continuous
version.

\begin{example}
\label{ex:integral_operator_adjoint}
For $\mathcal{K}$ as in \eqref{eq:integral_operator}, $\langle\mathcal{K}f,
g\rangle = \iint\kappa(t,s)f(s)g(t)\,ds\,dt = \langle f,\int_E\kappa(t,
\cdot)g(t)\,dt\rangle$, so $\mathcal{K}^*$ has kernel $\kappa(t,s)$
(arguments reversed); $\mathcal{K}$ is self-adjoint exactly when $\kappa$
is symmetric.
\end{example}

\begin{theorem}
\label{thm:kernel_nonneg_def}
The integral operator $\mathcal{K}$ is nonnegative definite if and only if
its kernel $\kappa$ is a nonnegative definite function, i.e.\
$\sum_{i,j}a_ia_j\kappa(t_i,t_j)\geq 0$ for all finite $\{a_i\}\subset
\mathbb{R}$, $\{t_i\}\subset E$.
\end{theorem}

\begin{proof}
By compactness of $E$ and uniform continuity of $\kappa$, for any $n$
there is a finite partition $\{E_{ni}\}$ of $E$ of mesh $<\delta_n$ (with
$\delta_n$ chosen so that $\kappa$ varies by less than $n^{-1}$ within a
set of that diameter, via the Heine--Borel Theorem) and points $v_i\in
E_{ni}$ such that the step kernel $\kappa_n(s,t):=\kappa(v_i,v_j)$ for
$(s,t)\in E_{ni}\times E_{nj}$ satisfies $\sup_{s,t}|\kappa(s,t)-
\kappa_n(s,t)|<n^{-1}$. Then $|\langle\mathcal{K}f,f\rangle-\langle
\mathcal{K}_nf,f\rangle|\leq n^{-1}\|f\|^2\to 0$, and $\langle\mathcal{K}_n
f,f\rangle = \sum_{i,j}\kappa(v_i,v_j)\int_{E_{ni}}f\,d\mu\int_{E_{nj}}
f\,d\mu$, which is $\geq 0$ for every $f$ whenever $\kappa$ is nonnegative
definite; letting $n\to\infty$ gives $\langle\mathcal{K}f,f\rangle\geq 0$.

Conversely, suppose $\sum_{i,j=1}^ma_ia_j\kappa(v_i,v_j)<0$ for some
$\{a_i\},\{v_i\}$. By continuity there are disjoint $E_1,\dots,E_m$ with
$\mu(E_i)>0$, $v_i\in E_i$, and $\sum_{i,j}a_ia_j\kappa(u_i,w_j)<0$
uniformly for $u_i,w_i\in E_i$. Averaging over $E_i\times E_j$ (mean value
theorem) gives $\sum_{i,j}a_ia_j\mu(E_i)^{-1}\mu(E_j)^{-1}\int_{E_i}
\int_{E_j}\kappa\,d\mu\,d\mu <0$, which is $\langle\mathcal{K}f,f\rangle$
for $f=\sum_ia_i\mu(E_i)^{-1}I_{E_i}$; so $\mathcal{K}$ is not nonnegative
definite.
\end{proof}

We now come to Mercer's Theorem, which identifies the eigen-decomposition
of a symmetric, nonnegative-definite integral operator with a convergent
series expansion of its kernel --- exactly the representation that, in
Section~\ref{sec:random_elements}, connects a covariance kernel
$\gamma(s,t)$ to the eigenfunctions that FPCA estimates.

\begin{lemma}
\label{lem:mercer_lemma}
Let $\kappa$ be continuous, symmetric, and nonnegative definite, with
$\mathcal{K}$ its integral operator and eigenpairs $\{(\lambda_j,e_j)\}$.
Then: (1) $\sum_j\lambda_je_j^2(t)\leq\kappa(t,t)$ for all $t$; (2)
$\sum_j|\lambda_je_j(s)e_j(t)|\leq\{\kappa(t,t)\kappa(s,s)\}^{1/2}$ for all
$s,t$; (3) $\lim_{n\to\infty}\sup_{s,t}\sum_{j>n}|\lambda_je_j(s)e_j(t)|=0$;
(4) $\sum_j\lambda_je_j(s)e_j(t)$ is well defined, and converges absolutely
and uniformly to a (uniformly) continuous function of $(s,t)$.
\end{lemma}

\begin{proof}
(1) Let $\kappa_n(s,t)=\kappa(s,t)-\sum_{j=1}^n\lambda_je_j(s)e_j(t)$, with
integral operator $\mathcal{K}_n$. For any $f$, $\langle\mathcal{K}_nf,f
\rangle=\langle\mathcal{K}f,f\rangle-\sum_{j=1}^n\lambda_j\langle f,e_j
\rangle^2=\sum_{j>n}\lambda_j\langle f,e_j\rangle^2\geq 0$, so
$\mathcal{K}_n$ is nonnegative definite, hence so is $\kappa_n$ by Theorem
\ref{thm:kernel_nonneg_def}; in particular $\kappa_n(t,t)\geq 0$, which is
the claim.

(2) By Cauchy--Schwarz applied to part (1), for any finite index set $J$,
$\sum_{j\in J}|\lambda_je_j(s)e_j(t)|\leq\big(\sum_{j\in J}\lambda_j
e_j^2(s)\big)^{1/2}\big(\sum_{j\in J}\lambda_je_j^2(t)\big)^{1/2}\leq
\{\kappa(s,s)\kappa(t,t)\}^{1/2}$.

(3) By the same Cauchy--Schwarz bound with $J=\{n+1,n+2,\dots\}$,
$\sum_{j>n}|\lambda_je_j(s)e_j(t)| \leq \big(\sum_{j>n}\lambda_je_j^2(s)
\big)^{1/2}\big(\sum_{j>n}\lambda_je_j^2(t)\big)^{1/2}$; by part (1) and
dominated convergence, $\sum_{j>n}\lambda_je_j^2(t)\downarrow 0$ for every
$t$ as $n\to\infty$, and by Dini's Theorem (monotone convergence of
continuous functions to a continuous limit, on a compact space) this
convergence is uniform in $t$; the claim follows.

(4) Given $\epsilon>0$, part (3) gives $n_\epsilon$ with $\sup_{s,t}
\sum_{j>n_\epsilon}|\lambda_je_j(s)e_j(t)|<\epsilon$, and uniform
continuity of $e_1,\dots,e_{n_\epsilon}$ gives $\delta>0$ such that
$\big|\sum_{j\leq n_\epsilon}\lambda_je_j(s)e_j(t)-\sum_{j\leq n_\epsilon}
\lambda_je_j(s')e_j(t')\big|<\epsilon$ whenever $d((s,t),(s',t'))<\delta$;
combining these two facts gives uniform continuity of the (absolutely,
uniformly convergent, by (3)) limit.
\end{proof}

\begin{theorem}[Mercer's Theorem]
\label{thm:mercer}
Let $\kappa$ be continuous, symmetric, and nonnegative definite on
$E\times E$, with integral operator $\mathcal{K}$ and eigenpairs
$\{(\lambda_j,e_j)\}$. Then
\begin{equation}
\kappa(s,t) = \sum_{j\geq 1}\lambda_j e_j(s)e_j(t), \qquad s,t\in E,
\label{eq:mercer}
\end{equation}
with the series converging absolutely and uniformly.
\end{theorem}

\begin{proof}
By Lemma~\ref{lem:mercer_lemma}(4), $\kappa'(s,t):=\sum_j\lambda_je_j(s)
e_j(t)$ is a well-defined continuous kernel, and the integral operator
$\mathcal{K}'$ with kernel $\kappa'$ satisfies $\mathcal{K}'e_j = \lambda_j
e_j$ for every $j$ and $\mathcal{K}'e = 0$ for $e\perp\overline{
\mathrm{span}}\{e_j\}$ (termwise integration, justified by uniform
convergence), i.e.\ $\mathcal{K}'$ has the same eigen-decomposition as
$\mathcal{K}$, so $\mathcal{K}'=\mathcal{K}$ by
\eqref{eq:spectral_decomp}. Since distinct continuous kernels define
distinct integral operators (if $\kappa_1\neq\kappa_2$ on a set of
positive $\mu\times\mu$-measure, some $f$ separates $\int\kappa_1(s,\cdot)
f(s)d\mu(s)$ from $\int\kappa_2(s,\cdot)f(s)d\mu(s)$), $\kappa=\kappa'$.
\end{proof}

\begin{theorem}
\label{thm:mercer_trace_hs}
Under the conditions of Theorem~\ref{thm:mercer}, $\mathcal{K}$ is trace
class with
\begin{equation}
\mathrm{trace}(\mathcal{K}) = \int_E\kappa(s,s)\,d\mu(s), \qquad
\|\mathcal{K}\|_{HS}^2 = \iint_{E\times E}\kappa^2(s,t)\,d\mu(s)d\mu(t).
\label{eq:mercer_trace_hs}
\end{equation}
\end{theorem}

\begin{proof}
By \eqref{eq:mercer} and Fubini's Theorem (justified by uniform
convergence), $\mathrm{trace}(\mathcal{K})=\sum_j\lambda_j = \sum_j
\lambda_j\int_Ee_j^2(t)\,d\mu(t) = \int_E\sum_j\lambda_je_j^2(t)\,d\mu(t) =
\int_E\kappa(t,t)\,d\mu(t)$, finite by Lemma~\ref{lem:mercer_lemma}(1); so
$\mathcal{K}$ is trace class. For the HS norm, let $\kappa_n(s,t)=
\sum_{j=1}^n\lambda_je_j(s)e_j(t)$ with integral operator $\mathcal{K}_n$.
Then $\mathrm{trace}(\mathcal{K}_n^2)=\sum_{j=1}^n\lambda_j^2\uparrow
\sum_j\lambda_j^2=\|\mathcal{K}\|_{HS}^2$ (by \eqref{eq:mercer_trace_hs}'s
first identity applied to $\mathcal{K}_n^2$, which has kernel
$\int_E\kappa_n(s,u)\kappa_n(u,t)\,d\mu(u)$, and equals $\iint\kappa_n^2\,
d\mu\,d\mu$ by orthonormality of $\{e_j\}$), and $\iint\kappa_n^2\,d\mu\,
d\mu\to\iint\kappa^2\,d\mu\,d\mu$ by Theorem~\ref{thm:mercer} and dominated
convergence.
\end{proof}

Combining Mercer's Theorem with the Eckart--Young Theorem
(\ref{thm:eckart_young}) yields the fundamental optimality property of the
truncated eigen-expansion of a covariance kernel that underlies FPCA
dimension reduction: no rank-$r$ kernel approximates $\kappa$ better, in
mean-square, than the truncated Mercer expansion.

\begin{theorem}
\label{thm:mercer_optimal_rank}
Under the conditions of Theorem~\ref{thm:mercer}, for any $r$ with
$\lambda_r>0$,
\begin{equation}
\min_{\mathrm{rank}(W)=r}\iint_{E\times E}\{\kappa(s,t)-W(s,t)\}^2\,d\mu(s)
d\mu(t) = \sum_{j>r}\lambda_j^2,
\label{eq:mercer_optimal_rank}
\end{equation}
where the minimum, over kernels $W$ of rank $r$ (i.e.\ whose integral
operator has rank $r$), is attained at $W(s,t)=\sum_{j=1}^r\lambda_je_j(s)
e_j(t)$.
\end{theorem}

\begin{proof}
By Theorem~\ref{thm:eckart_young}, the operator $\mathcal{W}=\sum_{j=1}^r
\lambda_je_j\otimes e_j$ minimizes $\|\mathcal{K}-\mathcal{W}\|_{HS}^2$
over rank-$r$ operators $\mathcal{W}$, with minimal value $\sum_{j>r}
\lambda_j^2$. Since $\mathcal{W}$ is an integral operator with (rank-$r$)
kernel $W(s,t)=\sum_{j=1}^r\lambda_je_j(s)e_j(t)$, and by
\eqref{eq:mercer_trace_hs} applied to $\mathcal{K}-\mathcal{W}$ (whose
kernel is $\kappa-W$), $\|\mathcal{K}-\mathcal{W}\|_{HS}^2 = \iint\{\kappa-
W\}^2\,d\mu\,d\mu$, the result follows.
\end{proof}

In Section~\ref{sec:metrics}, this theorem is the exact theoretical
counterpart of the empirical observation that FPCA truncated to $n_{80}$
components is the best rank-$n_{80}$ mean-square summary of the estimated
cross-country covariance structure available from any basis; it is what
makes ``variance explained'' a meaningful optimality criterion rather than
a merely descriptive one.

\section{Perturbation Theory}
\label{sec:perturbation}

Theorem~\ref{thm:eigenvalue_perturbation} already showed that eigenvalues
of a compact self-adjoint operator vary in a $1$-Lipschitz way under
perturbation of the operator, in operator norm. This section develops the
finer perturbation theory needed for the FPCA asymptotics of
Section~\ref{sec:fpca_theory}: expansions, to first order, of the
eigenprojections, eigenvalues, and eigenvectors of a compact self-adjoint
operator $\mathcal{T}$ under a perturbation $\Delta = \widetilde{\mathcal{T}}
-\mathcal{T}$, together with explicit remainder bounds. Throughout,
$\mathcal{T}$ plays the role of a population operator (e.g.\ a true
covariance operator) and $\widetilde{\mathcal{T}}$ the role of an
estimator; the results below are what let us convert operator-norm
consistency of $\widetilde{\mathcal{T}}$ for $\mathcal{T}$
(Section~\ref{sec:mean_cov_estimation}) into consistency, and asymptotic
normality, of the corresponding estimated eigenvalues and eigenfunctions.
We follow \cite{hsingeubank2015}, Section~5.1, restricting throughout to
the self-adjoint case (the non-self-adjoint singular-value perturbation
theory of that source's Section~5.2 is developed there for canonical
correlation analysis and is not needed for FPCA).

\subsection{The Resolvent and Contour Integration}
\label{subsec:resolvent}

The technical device underlying this section is contour integration of an
operator-valued function, which requires the scalar field to be extended
from $\mathbb{R}$ to $\mathbb{C}$; this complexification is a purely
technical device (Definition~\ref{def:hilbert_complex} below is the only
place it is used), and all operators, eigenvalues, and eigenvectors of
interest remain real.

\begin{definition}
\label{def:hilbert_complex}
A \emph{complex} Hilbert space is defined as in Definition~\ref{def:Lp}
(Hilbert space), with the scalar field $\mathbb{R}$ replaced by
$\mathbb{C}$ and condition (iv) of the inner product (symmetry) replaced
by conjugate symmetry, $\langle x,y\rangle=\overline{\langle y,x\rangle}$.
Every real Hilbert space $\mathbb{H}$ embeds in its complexification in the
natural way, and a (real) self-adjoint $\mathcal{T}\in\mathcal{B}
(\mathbb{H})$ extends to a self-adjoint operator on the complexification.
\end{definition}

\begin{definition}
\label{def:spectrum_resolvent}
For $\mathcal{T}\in\mathcal{B}(\mathbb{H})$ (on the complexification), the
\emph{spectrum} $\sigma(\mathcal{T})$ is the set of $z\in\mathbb{C}$ for
which $\mathcal{T}-zI$ is not invertible in $\mathcal{B}(\mathbb{H})$; its
complement $\rho(\mathcal{T})$ is the \emph{resolvent set}. For a compact
operator on an infinite-dimensional $\mathbb{H}$, $\sigma(\mathcal{T})$
consists exactly of the eigenvalues of $\mathcal{T}$ together with $0$
(whether or not $0$ is itself an eigenvalue). For $z\in\rho(\mathcal{T})$,
using the spectral decomposition $\mathcal{T}=\sum_j\lambda_j\mathcal{P}_j$
(with $\mathcal{P}_j$ the projection onto the eigenspace of the
nonrepeated eigenvalue $\lambda_j$, so $\mathcal{P}_j = \sum_{k=1}^{n_j}
e_{jk}\otimes e_{jk}$ for an orthonormal basis $e_{j1},\dots,e_{jn_j}$ of
that eigenspace of dimension $n_j$), the \emph{resolvent} of $\mathcal{T}$
is
\begin{equation}
\mathcal{R}(z) := (\mathcal{T}-zI)^{-1} = \sum_{j\geq 1}\frac{1}{\lambda_j
-z}\mathcal{P}_j,
\label{eq:resolvent}
\end{equation}
a bounded operator satisfying $\|\mathcal{R}(z)\| = 1/\min_j|z-\lambda_j|$.
\end{definition}

The bound on $\|\mathcal{R}(z)\|$ follows since, for
$f=\sum_j\mathcal{P}_jf$, $\|\mathcal{R}(z)f\|^2 = \sum_j|z-\lambda_j|^{-2}
\|\mathcal{P}_jf\|^2 \leq \max_j|z-\lambda_j|^{-2}\|f\|^2$, with equality
approached by concentrating $f$ near the eigenspace closest to $z$.

\begin{definition}
\label{def:contour_integral}
Let $\Omega$ be a Jordan curve (simple closed curve) in $\mathbb{C}$ of
length $l_\Omega$, traversed counterclockwise, with $d_\Omega(z,z')$ the
arc length from $z$ to $z'$ along $\Omega$. For an indexed family
$\{\mathcal{T}(z):z\in\Omega\}\subset\mathcal{B}(\mathbb{H})$, and ordered
points $z_0,\dots,z_n=z_0$ on $\Omega$ with $\xi_j$ on the arc from $z_j$
to $z_{j+1}$, define the Riemann-type sum $\mathcal{I}(\{z_j,\xi_j\}) :=
\sum_{j=0}^{n-1}\mathcal{T}(\xi_j)(z_{j+1}-z_j)$. If $\{z_{jn},\xi_{jn}\}$
is a sequence of such partitions with mesh $\to 0$ and
$\sup_{d_\Omega(z,z')\leq\delta}\|\mathcal{T}(z)-\mathcal{T}(z')\|\to 0$ as
$\delta\downarrow 0$, the \emph{contour integral} is
\begin{equation}
\oint_\Omega\mathcal{T}(z)\,dz := \lim_{n\to\infty}\mathcal{I}(\{z_{jn},
\xi_{jn}\}).
\label{eq:contour_integral}
\end{equation}
\end{definition}

\begin{lemma}
\label{lem:contour_cauchy}
Under the modulus-of-continuity condition of Definition
\ref{def:contour_integral}, the limit \eqref{eq:contour_integral} exists
and does not depend on the approximating sequence of partitions.
\end{lemma}

\begin{proof}
For two partitions $\{z_j,\xi_j\}$, $\{\widetilde{z}_j,\widetilde\xi_j\}$
with mesh $\leq\delta$, merge them into a single ordered partition
$\{\check{z}_j,\check\xi_j\}$ refining both; writing the two Riemann sums
as sums over the refined partition and subtracting term by term (each
original arc is split into subarcs on which the corresponding $\xi_j$ or
$\widetilde\xi_j$ still lies), $\|\mathcal{I}(\{z_j,\xi_j\})-\mathcal{I}
(\{\widetilde{z}_j,\widetilde\xi_j\})\|\leq\sum_j\|\mathcal{T}(\check\xi_j)
-\mathcal{T}(\check\xi_j')\|\,d_\Omega(\check{z}_{j-1},\check{z}_j)\leq
\sup_{d_\Omega(z,z')\leq\delta}\|\mathcal{T}(z)-\mathcal{T}(z')\|\,l_\Omega$
by the triangle inequality. Applying this to a sequence of partitions with
mesh $\to 0$ shows $\{\mathcal{I}(\{z_{jn},\xi_{jn}\})\}$ is Cauchy in the
Banach space $\mathcal{B}(\mathbb{H})$, hence convergent; applying it to
two such sequences merged together shows the limit is independent of the
sequence.
\end{proof}

Contour integration of an operator-valued function inherits linearity from
\eqref{eq:contour_integral} directly. Applying it to the resolvent yields a
contour-integral representation of an eigenprojection, via the residue
calculus fact $\frac{1}{2\pi i}\oint_\Omega(\lambda_k-z)^{-1}dz =
-I(\lambda_k\text{ inside }\Omega)$ (Cauchy's integral formula, taking the
contour to enclose no other singularity).

\begin{theorem}
\label{thm:projection_contour}
Let $\mathcal{T}$ be compact self-adjoint and let $\Omega$ be the boundary
of a disk containing exactly the eigenvalues $\{\lambda_k:p\leq k\leq q\}$
(repeated eigenvalues, i.e.\ counted with multiplicity, all strictly
interior to $\Omega$) and no others. Then the projection $\mathcal{P}$ onto
the union of the corresponding eigenspaces is
\begin{equation}
\mathcal{P} = -\frac{1}{2\pi i}\oint_\Omega \mathcal{R}(z)\,dz.
\label{eq:projection_contour}
\end{equation}
\end{theorem}

\begin{proof}
Since $\|\mathcal{R}(z)-\mathcal{R}(z')\| = \max_j\big|\frac{1}{z-\lambda_j}
-\frac{1}{z'-\lambda_j}\big| = \max_j\frac{|z-z'|}{|z-\lambda_j||z'-
\lambda_j|} \leq |z-z'|/\epsilon^2$ for $\epsilon:=\inf_{z\in\Omega,j}
|z-\lambda_j|>0$, the modulus-of-continuity condition of Definition
\ref{def:contour_integral} holds and \eqref{eq:projection_contour} is
well defined. By \eqref{eq:resolvent} and linearity of the contour
integral, $-\frac{1}{2\pi i}\oint_\Omega\mathcal{R}(z)\,dz = -\frac{1}
{2\pi i}\sum_k\mathcal{P}_k\oint_\Omega(\lambda_k-z)^{-1}\,dz =
\sum_{p\leq k\leq q}\mathcal{P}_k = \mathcal{P}$, applying Cauchy's
integral formula termwise (justified by the uniform bound $\|\mathcal{R}
(z)\|\leq 1/\epsilon$ on $\Omega$, which makes the interchange of the
(uniformly convergent) sum and the contour integral valid).
\end{proof}

\subsection{Perturbation of Eigenprojections}
\label{subsec:eigenprojection_perturbation}

We now compare the eigenprojections of a compact self-adjoint
$\mathcal{T}$, with nonrepeated eigenvalues $\lambda_1>\lambda_2>\cdots$
and multiplicities $n_j$, to those of a perturbation
$\widetilde{\mathcal{T}}=\mathcal{T}+\Delta$. When $\mathcal{T}$ has
repeated eigenvalues, the eigenspaces of $\widetilde{\mathcal{T}}$ must be
matched to those of $\mathcal{T}$ by dimension: writing $\lambda_1\geq
\lambda_2\geq\cdots$ for the \emph{repeated} eigenvalues of
$\widetilde{\mathcal{T}}$ (i.e.\ listed with multiplicity) with
eigenvectors $\{\widetilde{e}_k\}$ and $N_j:=\sum_{k=1}^jn_k$, define
\begin{equation}
\widetilde{\mathcal{P}}_j := \sum_{k=N_{j-1}+1}^{N_j}\widetilde{e}_k
\otimes\widetilde{e}_k,
\label{eq:matched_projection}
\end{equation}
i.e.\ the projection onto the span of exactly $n_j$ eigenvectors of
$\widetilde{\mathcal{T}}$, matching the dimension of $\mathcal{P}_j$; this
is the pairing used throughout.

\begin{theorem}
\label{thm:projection_perturbation}
Let $\mathcal{T}$ be compact self-adjoint with nonrepeated eigenvalue
$\lambda_j$, eigenprojection $\mathcal{P}_j$, and $\eta_j := \frac{1}{2}
\min_{k\neq j}|\lambda_k-\lambda_j|$. Let $\widetilde{\mathcal{T}}=
\mathcal{T}+\Delta$ be compact self-adjoint with $\|\Delta\|<\eta_j$, and
let $\widetilde{\mathcal{P}}_j$ be as in \eqref{eq:matched_projection}.
Then
\begin{equation}
\widetilde{\mathcal{P}}_j - \mathcal{P}_j = \mathcal{S}_j\Delta\mathcal{P}_j
+ \mathcal{P}_j\Delta\mathcal{S}_j + \frac{1}{2\pi i}\oint_{\Omega_j}
\mathcal{M}(z)\,dz,
\label{eq:projection_perturbation}
\end{equation}
where $\Omega_j$ is the circle centered at $\lambda_j$ of radius $\eta_j$,
$\mathcal{S}_j := \sum_{k\neq j}(\lambda_k-\lambda_j)^{-1}\mathcal{P}_k$,
and $\mathcal{M}(z) := \mathcal{R}(z)\sum_{k\geq 2}\{-\Delta\mathcal{R}
(z)\}^k$. Consequently,
\begin{equation}
\Big\|\widetilde{\mathcal{P}}_j-\mathcal{P}_j - (\mathcal{S}_j\Delta
\mathcal{P}_j+\mathcal{P}_j\Delta\mathcal{S}_j)\Big\| \leq
\frac{\delta_j^2}{1-\delta_j}, \qquad \delta_j := \|\Delta\|/\eta_j,
\label{eq:projection_perturbation_bound}
\end{equation}
and, since $\|\mathcal{S}_j\|\leq(2\eta_j)^{-1}$,
\begin{equation}
\|\mathcal{S}_j\Delta\mathcal{P}_j+\mathcal{P}_j\Delta\mathcal{S}_j\|\leq
\delta_j/\sqrt2, \qquad\text{so that}\qquad
\|\widetilde{\mathcal{P}}_j-\mathcal{P}_j\| \leq \delta_j/\sqrt2 +
\delta_j^2/(1-\delta_j) \leq \delta_j/(1-\delta_j).
\label{eq:projection_norm_bound}
\end{equation}
\end{theorem}

\begin{proof}
Let $\widetilde{\mathcal{R}}$ denote the resolvent of
$\widetilde{\mathcal{T}}$. By Theorem~\ref{thm:eigenvalue_perturbation},
$|\widetilde\lambda-\lambda_j|\leq\|\Delta\|<\eta_j$ for every eigenvalue
$\widetilde\lambda$ of $\widetilde{\mathcal{T}}$ matched to $\lambda_j$, so
$\Omega_j$ encircles exactly the eigenvalues of $\widetilde{\mathcal{T}}$
corresponding to $\lambda_j$ and no others; Theorem
\ref{thm:projection_contour} applies to both $\mathcal{T}$ and
$\widetilde{\mathcal{T}}$ with this contour, giving
\begin{equation}
\widetilde{\mathcal{P}}_j-\mathcal{P}_j = -\frac{1}{2\pi i}\oint_{\Omega_j}
\big(\widetilde{\mathcal{R}}(z)-\mathcal{R}(z)\big)\,dz.
\label{eq:proj_diff_resolvents}
\end{equation}
Since $\widetilde{\mathcal{R}}(z) = (\Delta+\mathcal{T}-zI)^{-1} =
\mathcal{R}(z)(\Delta\mathcal{R}(z)+I)^{-1}$ and, on $\Omega_j$,
$\|\Delta\mathcal{R}(z)\|\leq\|\Delta\|/\eta_j=\delta_j<1$ (using
$\|\mathcal{R}(z)\|=1/\eta_j$ on $\Omega_j$ by Definition
\ref{def:spectrum_resolvent}), the Neumann series gives $\widetilde{
\mathcal{R}}(z)-\mathcal{R}(z) = \mathcal{R}(z)\sum_{k\geq 1}\{-\Delta
\mathcal{R}(z)\}^k = -\mathcal{R}(z)\Delta\mathcal{R}(z)+\mathcal{M}(z)$.
Substituting into \eqref{eq:proj_diff_resolvents},
\[
\widetilde{\mathcal{P}}_j-\mathcal{P}_j = \frac{1}{2\pi i}\oint_{\Omega_j}
\mathcal{R}(z)\Delta\mathcal{R}(z)\,dz + \frac{1}{2\pi i}\oint_{\Omega_j}
\mathcal{M}(z)\,dz,
\]
and, using $\mathcal{R}(z)\Delta\mathcal{R}(z) = \sum_{k,l}(\lambda_k-z)^{-1}
(\lambda_l-z)^{-1}\mathcal{P}_k\Delta\mathcal{P}_l$ and the elementary
contour integral $\frac{1}{2\pi i}\oint_{\Omega_j}(\lambda_k-z)^{-1}
(\lambda_l-z)^{-1}dz$, which equals $-(\lambda_k-\lambda_j)^{-1}$ if
$k\neq j,l=j$; $(\lambda_l-\lambda_j)^{-1}$ if $k=j,l\neq j$; and $0$
otherwise (by Cauchy's integral formula, since exactly one of the two
simple poles $\lambda_k,\lambda_l$ lies inside $\Omega_j$ in these cases,
and none or both do in the remaining cases, the latter contributing $0$ by
residues at infinity / direct partial-fraction cancellation), the first
term equals $\sum_{k\neq j}(\lambda_k-\lambda_j)^{-1}(\mathcal{P}_k\Delta
\mathcal{P}_j+\mathcal{P}_j\Delta\mathcal{P}_k) = \mathcal{S}_j\Delta
\mathcal{P}_j+\mathcal{P}_j\Delta\mathcal{S}_j$, proving
\eqref{eq:projection_perturbation}.

For \eqref{eq:projection_perturbation_bound}: $\|\mathcal{M}(z)\| \leq
\|\mathcal{R}(z)\|\sum_{k\geq 2}\|\Delta\mathcal{R}(z)\|^k \leq
\eta_j^{-1}\delta_j^2/(1-\delta_j)$ on $\Omega_j$, so $\big\|\frac{1}{2\pi
i}\oint_{\Omega_j}\mathcal{M}(z)\,dz\big\| \leq \frac{l_{\Omega_j}}{2\pi}
\eta_j^{-1}\delta_j^2/(1-\delta_j) = \delta_j^2/(1-\delta_j)$, since
$l_{\Omega_j}=2\pi\eta_j$.

For \eqref{eq:projection_norm_bound}: $\|\mathcal{S}_j\Delta\mathcal{P}_j+
\mathcal{P}_j\Delta\mathcal{S}_j\|^2 = \|\mathcal{S}_j\Delta\mathcal{P}_j
\|^2+\|\mathcal{P}_j\Delta\mathcal{S}_j\|^2 \leq 2\|\mathcal{S}_j\|^2
\|\Delta\|^2$ (orthogonality of the two terms' ranges/domains), and
$\|\mathcal{S}_j\|=\max_{k\neq j}|\lambda_k-\lambda_j|^{-1}\leq(2\eta_j)^{-1}$
gives the stated bound; combining with
\eqref{eq:projection_perturbation_bound} via the triangle inequality gives
the norm bound on $\widetilde{\mathcal{P}}_j-\mathcal{P}_j$.
\end{proof}

Theorem~\ref{thm:projection_perturbation} shows that
$\mathcal{S}_j\Delta\mathcal{P}_j+\mathcal{P}_j\Delta\mathcal{S}_j$ is a
first-order (linear in $\Delta$) approximation to
$\widetilde{\mathcal{P}}_j-\mathcal{P}_j$, with remainder $O(\delta_j^2)$.
This linearization is the technical heart of the asymptotic normality
results in Section~\ref{sec:fpca_theory}: it expresses the perturbation of
an eigenprojection as an explicit \emph{linear} functional of $\Delta$
(hence, when $\Delta$ is itself asymptotically Gaussian, so is the
leading-order term of $\widetilde{\mathcal{P}}_j-\mathcal{P}_j$).

\subsection{Perturbation of Eigenvalues}
\label{subsec:eigenvalue_perturbation}

\begin{theorem}
\label{thm:eigenvalue_perturbation_exact}
Under the conditions of Theorem~\ref{thm:projection_perturbation}, with
$\widetilde\lambda_{N_{j-1}+1},\dots,\widetilde\lambda_{N_j}$ the (repeated)
eigenvalues of $\widetilde{\mathcal{T}}$ matched to $\lambda_j$ and
$\widetilde{e}_{N_{j-1}+1},\dots,\widetilde{e}_{N_j}$ the corresponding
orthonormal eigenvectors used in \eqref{eq:matched_projection},
\begin{equation}
\widetilde{\mathcal{P}}_j\widetilde{\mathcal{T}}\widetilde{\mathcal{P}}_j -
\lambda_j\widetilde{\mathcal{P}}_j = \widetilde{\mathcal{P}}_j\Delta
\widetilde{\mathcal{P}}_j + (\widetilde{\mathcal{P}}_j-\mathcal{P}_j)
(\mathcal{T}-\lambda_jI)(\widetilde{\mathcal{P}}_j-\mathcal{P}_j),
\label{eq:eigenvalue_perturbation_operator}
\end{equation}
and, for each $k=N_{j-1}+1,\dots,N_j$,
\begin{equation}
\widetilde\lambda_k - \lambda_j = \langle\Delta\widetilde{e}_k,
\widetilde{e}_k\rangle + \big\langle(\widetilde{\mathcal{P}}_j-\mathcal{P}_j)
(\mathcal{T}-\lambda_jI)(\widetilde{\mathcal{P}}_j-\mathcal{P}_j)
\widetilde{e}_k,\widetilde{e}_k\big\rangle,
\label{eq:eigenvalue_perturbation_exact}
\end{equation}
where the second term is $O(\|\widetilde{\mathcal{P}}_j-\mathcal{P}_j\|^2)
= O(\delta_j^2)$ by \eqref{eq:projection_norm_bound}, since it is bounded
in norm by $\|\mathcal{T}\|\,\|\widetilde{\mathcal{P}}_j-\mathcal{P}_j\|^2$.
\end{theorem}

\begin{proof}
Expanding $\widetilde{\mathcal{P}}_j\widetilde{\mathcal{T}}\widetilde{
\mathcal{P}}_j - \lambda_j\widetilde{\mathcal{P}}_j = \widetilde{
\mathcal{P}}_j\Delta\widetilde{\mathcal{P}}_j + \widetilde{\mathcal{P}}_j
(\mathcal{T}-\lambda_jI)\widetilde{\mathcal{P}}_j$, and using
$\mathcal{P}_j(\mathcal{T}-\lambda_jI)=0$ (since $\mathcal{P}_j$ projects
onto the $\lambda_j$-eigenspace) to write $\widetilde{\mathcal{P}}_j
(\mathcal{T}-\lambda_jI)\widetilde{\mathcal{P}}_j = (\widetilde{
\mathcal{P}}_j-\mathcal{P}_j)(\mathcal{T}-\lambda_jI)(\widetilde{
\mathcal{P}}_j-\mathcal{P}_j)$ (the cross terms vanish since
$(\mathcal{T}-\lambda_jI)\mathcal{P}_j=0=\mathcal{P}_j(\mathcal{T}-
\lambda_jI)$), proves \eqref{eq:eigenvalue_perturbation_operator}. Taking
the inner product of both sides with $\widetilde{e}_k$ and using
$\widetilde{\mathcal{P}}_j\widetilde{\mathcal{T}}\widetilde{\mathcal{P}}_j
= \sum_{k=N_{j-1}+1}^{N_j}\widetilde\lambda_k\widetilde{e}_k\otimes
\widetilde{e}_k$ and $\widetilde{\mathcal{P}}_j\widetilde{e}_k=
\widetilde{e}_k$ gives \eqref{eq:eigenvalue_perturbation_exact}.
\end{proof}

Identity \eqref{eq:eigenvalue_perturbation_exact} refines the crude bound
of Theorem~\ref{thm:eigenvalue_perturbation}: to first order in $\Delta$,
$\widetilde\lambda_k-\lambda_j\approx\langle\Delta\widetilde{e}_k,
\widetilde{e}_k\rangle$, with an explicit, quadratic-in-$\Delta$ remainder.
In the leading case of interest for FPCA --- $\lambda_j$ of multiplicity
one for every $j$, so $\mathcal{P}_j=e_j\otimes e_j$ and
$\widetilde{\mathcal{P}}_j=\widetilde{e}_j\otimes\widetilde{e}_j$ for a
single matched eigenvalue $\widetilde\lambda_j$ --- this reduces to
\begin{equation}
\widetilde\lambda_j - \lambda_j = \langle\Delta e_j,e_j\rangle +
\big\langle(\widetilde{\mathcal{P}}_j-\mathcal{P}_j)(\mathcal{T}-
\lambda_jI)(\widetilde{\mathcal{P}}_j-\mathcal{P}_j)e_j,e_j\big\rangle,
\label{eq:eigenvalue_perturbation_simple}
\end{equation}
obtained from \eqref{eq:eigenvalue_perturbation_exact} by symmetry of the
argument in $e_j,\widetilde{e}_j$ (an identical derivation with the roles
of $\mathcal{T},\widetilde{\mathcal{T}}$ exchanged in the projection
identity gives the version stated with $e_j$ in place of $\widetilde{e}_j$
in the leading term, up to the same $O(\delta_j^2)$ remainder). This
first-order expansion, $\widetilde\lambda_j-\lambda_j\approx\langle\Delta
e_j,e_j\rangle$, is precisely what will drive the asymptotic normality of
the estimated FPCA eigenvalues in Section~\ref{sec:fpca_theory}: since
$\Delta = \widetilde{\mathcal{T}}-\mathcal{T}$ is a linear functional of
the estimated covariance operator, a central limit theorem for
$\widetilde{\mathcal{T}}$ transfers, via this linearization, into one for
$\widetilde\lambda_j$.

\subsection{Perturbation of Eigenvectors}
\label{subsec:eigenvector_perturbation}

We first relate the norm of the difference between two unit eigenvectors
to the norm of the difference between their (rank-one) projections ---
which is the quantity directly controlled by Theorem
\ref{thm:projection_perturbation}.

\begin{lemma}
\label{lem:eigenvector_projection_norm}
Let $\mathcal{P}=e\otimes e$, $\widetilde{\mathcal{P}}=\widetilde{e}\otimes
\widetilde{e}$ for unit vectors $e,\widetilde{e}$. Then: (1) the
eigenvalues of $\mathcal{P}-\widetilde{\mathcal{P}}$ are $\pm(1-\langle
e,\widetilde{e}\rangle^2)^{1/2}$; (2) if $\langle e,\widetilde{e}\rangle
\geq 0$, $\|e-\widetilde{e}\|^2 = 2\big[1-(1-\|\mathcal{P}-
\widetilde{\mathcal{P}}\|^2)^{1/2}\big]$; (3) $\|\mathcal{P}-
\widetilde{\mathcal{P}}\|_{HS}^2 = 2(1-\langle e,\widetilde{e}\rangle^2)
= 2\|\mathcal{P}-\widetilde{\mathcal{P}}\|^2$.
\end{lemma}

\begin{proof}
(1) The range of $\mathcal{P}-\widetilde{\mathcal{P}}$ is contained in
$\mathrm{span}\{e,\widetilde{e}\}$, on which, in the basis $(e,
\widetilde{e})$, $\mathcal{P}-\widetilde{\mathcal{P}}$ acts as the matrix
$\left(\begin{smallmatrix}1&-\langle e,\widetilde{e}\rangle\\ \langle
e,\widetilde{e}\rangle&-1\end{smallmatrix}\right)$ (since $\mathcal{P}e=e$,
$\mathcal{P}\widetilde{e}=\langle\widetilde{e},e\rangle e$, and
similarly for $\widetilde{\mathcal{P}}$), whose characteristic polynomial
is $\mu^2-(1-\langle e,\widetilde{e}\rangle^2)=0$, giving the stated
eigenvalues.

(2) By part (1) and Theorem~\ref{thm:svd_operator_norm} (operator norm
equals largest singular value, and $\mathcal{P}-\widetilde{\mathcal{P}}$
is self-adjoint so its singular values are $|\text{eigenvalues}|$),
$\|\mathcal{P}-\widetilde{\mathcal{P}}\|^2 = 1-\langle e,\widetilde{e}
\rangle^2$. Combined with $\|e-\widetilde{e}\|^2 = 2(1-\langle e,
\widetilde{e}\rangle)$ and, for $\langle e,\widetilde{e}\rangle\geq 0$,
$\langle e,\widetilde{e}\rangle = (1-\|\mathcal{P}-\widetilde{\mathcal{P}}
\|^2)^{1/2}$, the claim follows.

(3) Immediate from part (1) (both nonzero eigenvalues of $\mathcal{P}-
\widetilde{\mathcal{P}}$ contribute equally to the HS norm) and the
operator-norm identity from part (2)'s proof.
\end{proof}

For $x\in[0,1]$, $1-(1-x)^{1/2}\leq x/2+x^2/(4(1-x))$ (Taylor expansion
with remainder of $(1-x)^{1/2}$), so part (2) of Lemma
\ref{lem:eigenvector_projection_norm} gives, whenever $\langle e,
\widetilde{e}\rangle\geq 0$,
\begin{equation}
\|e-\widetilde{e}\|^2 \leq \|\mathcal{P}-\widetilde{\mathcal{P}}\|^2 +
\frac{1}{2}\frac{\|\mathcal{P}-\widetilde{\mathcal{P}}\|^4}{1-\|\mathcal{P}
-\widetilde{\mathcal{P}}\|^2},
\label{eq:eigenvector_norm_sq_bound}
\end{equation}
and, since $(1+x)^{1/2}\leq 1+x/2$ for $x\in[0,1]$, taking square roots and
simplifying,
\begin{equation}
\|e-\widetilde{e}\| \leq \|\mathcal{P}-\widetilde{\mathcal{P}}\| +
\frac{1}{4}\frac{\|\mathcal{P}-\widetilde{\mathcal{P}}\|^3}{1-\|\mathcal{P}
-\widetilde{\mathcal{P}}\|^2}.
\label{eq:eigenvector_norm_bound}
\end{equation}

We can now state the eigenvector perturbation expansion that, together
with \eqref{eq:eigenvalue_perturbation_simple}, will be the direct input to
the FPCA asymptotic normality results of Section~\ref{sec:fpca_theory}.

\begin{theorem}
\label{thm:eigenvector_perturbation}
Let $\lambda_j$ be a simple (multiplicity-one) eigenvalue of $\mathcal{T}$
with unit eigenvector $e_j$; the other eigenvalues of $\mathcal{T}$ need
not be simple. Let $\widetilde{\mathcal{T}}=\mathcal{T}+\Delta$ be compact
self-adjoint with matched eigenpair $(\widetilde\lambda_j,\widetilde{e}_j)$
chosen with $\langle e_j,\widetilde{e}_j\rangle\geq 0$, and suppose
$\inf_{k\neq j}|\widetilde\lambda_j-\lambda_k|>0$. Then
\begin{equation}
\widetilde{e}_j - e_j = \sum_{k\neq j}(\widetilde\lambda_j-\lambda_k)^{-1}
\mathcal{P}_k\Delta\widetilde{e}_j + \mathcal{P}_j(\widetilde{e}_j-e_j).
\label{eq:eigenvector_expansion}
\end{equation}
Moreover, with $\eta_j=\frac{1}{2}\inf_{k\neq j}|\lambda_j-\lambda_k|$ and
$\delta_j=\|\Delta\|/\eta_j$, if $\delta_j<1$ there is a finite function
$\psi$ with $\psi(\delta)\downarrow 1$ as $\delta\downarrow 0$ such that
\begin{equation}
\|e_j - \widetilde{e}_j - \mathcal{S}_j\Delta e_j\| \leq \psi(\delta_j)
\delta_j^2, \qquad \mathcal{S}_j := \sum_{k\neq j}(\lambda_k-\lambda_j)^{-1}
\mathcal{P}_k.
\label{eq:eigenvector_perturbation_bound}
\end{equation}
\end{theorem}

\begin{proof}
For $k\neq j$, $(\widetilde\lambda_j-\lambda_k)\mathcal{P}_k(\widetilde{e}_j
-e_j) = (\widetilde\lambda_j-\lambda_k)\mathcal{P}_k\widetilde{e}_j =
\mathcal{P}_k\big(\widetilde{\mathcal{T}}\widetilde{e}_j - \mathcal{T}
\mathcal{P}_k\widetilde{e}_j\big) = \mathcal{P}_k\Delta\widetilde{e}_j +
\mathcal{P}_k(\mathcal{T}\widetilde{e}_j-\mathcal{T}\mathcal{P}_k
\widetilde{e}_j)$; since $\mathcal{P}_k\mathcal{T}=\lambda_k\mathcal{P}_k$
this simplifies (writing $\mathcal{T}\widetilde e_j=\sum_l\lambda_l
\mathcal{P}_l\widetilde e_j$) to $\mathcal{P}_k\Delta\widetilde{e}_j$, using
$\widetilde{\mathcal{T}}\widetilde{e}_j=\widetilde\lambda_j\widetilde{e}_j$.
Summing over $k\neq j$ and adding $\mathcal{P}_j(\widetilde{e}_j-e_j)$
recovers $\widetilde{e}_j-e_j=\sum_k\mathcal{P}_k(\widetilde{e}_j-e_j)$ by
Theorem~\ref{thm:spectral_theorem} (completeness of $\{\mathcal{P}_k\}$),
proving \eqref{eq:eigenvector_expansion}.

For \eqref{eq:eigenvector_perturbation_bound}: Theorem
\ref{thm:eigenvalue_perturbation} gives $|\widetilde\lambda_j-\lambda_j|
\leq\|\Delta\|<\eta_j\leq\inf_{k\neq j}|\lambda_j-\lambda_k|$, so
$\inf_{k\neq j}|\widetilde\lambda_j-\lambda_k| \geq \inf_{k\neq j}
|\lambda_j-\lambda_k|-\|\Delta\|>0$ and \eqref{eq:eigenvector_expansion}
applies. Expanding $(\widetilde\lambda_j-\lambda_k)^{-1}$ as a geometric
series in $(\lambda_j-\widetilde\lambda_j)/(\lambda_j-\lambda_k)$
(convergent since $|\lambda_j-\widetilde\lambda_j|<|\lambda_j-\lambda_k|$)
and substituting into \eqref{eq:eigenvector_expansion} decomposes
$\widetilde{e}_j-e_j$ into four terms: (i) $\sum_{k\neq j}(\lambda_j-
\lambda_k)^{-1}\mathcal{P}_k\Delta e_j = \mathcal{S}_j\Delta e_j$; (ii)
$\sum_{k\neq j}(\lambda_j-\lambda_k)^{-1}\mathcal{P}_k\Delta(\widetilde{e}_j
-e_j)$; (iii) the $s\geq 1$ terms of the geometric expansion, $\sum_{k\neq
j}\sum_{s\geq 1}(\lambda_j-\widetilde\lambda_j)^s(\lambda_j-\lambda_k)^{
-s-1}\mathcal{P}_k\Delta\widetilde{e}_j$; (iv) $\mathcal{P}_j(\widetilde{e}
_j-e_j)$.

Term (ii) is bounded, by Bessel's inequality, by $\|\Delta\|\|
\widetilde{e}_j-e_j\|/(2\eta_j) = (\delta_j/2)\|\widetilde{e}_j-e_j\|$.
Term (iii): by Parseval's relation (orthogonality of the $\mathcal{P}_k$
ranges), its squared norm is $\sum_{k\neq j}\big(\sum_{s\geq
1}(\lambda_j-\widetilde\lambda_j)^s(\lambda_j-\lambda_k)^{-s-1}\big)^2\|
\mathcal{P}_k\Delta\widetilde{e}_j\|^2 \leq \sum_{k\neq j}\frac{(\lambda_j-
\widetilde\lambda_j)^2}{(\lambda_j-\lambda_k)^2(|\lambda_j-\lambda_k|-
|\lambda_j-\widetilde\lambda_j|)^2}\|\mathcal{P}_k\Delta\widetilde{e}_j
\|^2 \leq \frac{\|\Delta\|^4}{4\eta_j^2(2\eta_j-\|\Delta\|)^2} \leq
\delta_j^4/4$ (bounding each denominator factor by $\eta_j\cdot\eta_j$,
using $\|\Delta\|<\eta_j$, and $\sum_{k\neq j}\|\mathcal{P}_k\Delta
\widetilde{e}_j\|^2\leq\|\Delta\|^2$ by Bessel). Term (iv): $\mathcal{P}_j
(\widetilde{e}_j-e_j) = \langle\widetilde{e}_j-e_j,e_j\rangle e_j =
(\langle\widetilde{e}_j,e_j\rangle-1)e_j$, so $\|\mathcal{P}_j(\widetilde{
e}_j-e_j)\| = 1-\langle\widetilde{e}_j,e_j\rangle = \frac{1}{2}\|
\widetilde{e}_j-e_j\|^2$.

Combining these four bounds with the triangle inequality, and using
\eqref{eq:eigenvector_norm_bound} together with
\eqref{eq:projection_norm_bound} to control $\|\widetilde{e}_j-e_j\|$ (of
order $\delta_j$) in the terms (ii) and (iv) that are quadratic or higher
in $\|\widetilde{e}_j-e_j\|$, yields $\|e_j-\widetilde{e}_j-\mathcal{S}_j
\Delta e_j\| \leq \psi(\delta_j)\delta_j^2$ for an explicit finite $\psi$
with $\psi(\delta)\downarrow 1$ as $\delta\downarrow 0$ (each of the four
terms other than $\mathcal{S}_j\Delta e_j$ contributes $O(\delta_j^2)$,
with constants that remain bounded as $\delta_j\downarrow 0$).
\end{proof}

As with the eigenvalue expansion \eqref{eq:eigenvalue_perturbation_simple},
Theorem~\ref{thm:eigenvector_perturbation} identifies $\mathcal{S}_j\Delta
e_j$ as the first-order (linear-in-$\Delta$) term of the eigenvector
perturbation $\widetilde{e}_j-e_j$, with an explicit $O(\delta_j^2)$
remainder. Together, \eqref{eq:eigenvalue_perturbation_simple} and
\eqref{eq:eigenvector_perturbation_bound} reduce the asymptotic behavior of
the estimated FPCA eigenvalue/eigenfunction pairs to the (much simpler)
asymptotic behavior of the linear functionals $\langle\Delta e_j,e_j
\rangle$ and $\mathcal{S}_j\Delta e_j$ of the covariance operator
perturbation $\Delta=\widetilde{\Gamma}-\Gamma$ --- which is exactly what
the large-sample theory of Sections~\ref{sec:random_elements} and
\ref{sec:mean_cov_estimation} establishes.

\section{Smoothing and Regularization}
\label{sec:smoothing}

Discretely sampled, possibly noisy functional data must be converted into
genuine elements of a function space before the Hilbert-space machinery of
the preceding sections applies. This section develops the general theory of
penalized least-squares (regularization) estimators that accomplishes this,
and specializes it to smoothing splines, which are the theoretical
counterpart of the adaptive B-spline smoothing used in
Section~\ref{subsec:bspline}. We follow \cite{hsingeubank2015}, Chapter~6.

\subsection{The Functional Linear Model}
\label{subsec:functional_linear_model}

\begin{definition}
\label{def:functional_linear_model}
Let $\mathbb{Y},\mathbb{H}$ be Hilbert spaces and $\mathcal{T}\in
\mathcal{B}(\mathbb{H},\mathbb{Y})$. The \emph{functional linear model} is
\begin{equation}
Y = \mathcal{T}m + \varepsilon, \qquad Y,\varepsilon\in\mathbb{Y},\ \ m\in
\mathbb{H},
\label{eq:functional_linear_model}
\end{equation}
where $\varepsilon$ is a mean-zero error element and $\mathcal{T}$ is
treated as known; the goal is to estimate $m$ from an observation of $Y$.
\end{definition}

The case directly relevant to Chapter~\ref{chap:methods} is the discretely
sampled setting: $\mathbb{Y}=\mathbb{R}^n$ with $\|Y\|_{\mathbb{Y}}^2 =
n^{-1}\sum_{i=1}^nY_i^2$, $\mathbb{H}$ an RKHS of functions on $[0,1]$,
and $\mathcal{T}g = (g(t_1),\dots,g(t_n))^T$ for observation
points $t_1,\dots,t_n$; boundedness of $\mathcal{T}$ is precisely
boundedness of the evaluation functionals, which characterizes an RKHS
(Definition~\ref{def:rkhs}, Section~\ref{subsec:rkhs_prelim}). Model
\eqref{eq:functional_linear_model} also subsumes the case where the sample
paths $x_1,\dots,x_n$ of a mean-$m$ process are themselves the object being
smoothed: writing $Y_{ij}=x_i(t_{ij})+\widetilde\varepsilon_{ij} =
m(t_{ij}) + \{x_i(t_{ij})-m(t_{ij})\} + \widetilde\varepsilon_{ij}$ and
stacking, $\varepsilon_{ij}:=x_i(t_{ij})-m(t_{ij})+\widetilde\varepsilon_{ij}$
recovers \eqref{eq:functional_linear_model} with a (possibly correlated)
error term.

\subsection{Penalized Least-Squares Estimation}
\label{subsec:penalized_ls}

Ordinary least squares, $\widehat{m}=\mathrm{argmin}_{g\in\mathbb{H}}\|Y-
\mathcal{T}g\|_{\mathbb{Y}}^2$, is ill-posed when $\mathbb{H}$ is
infinite-dimensional: any $g$ interpolating the data minimizes the
criterion, so the estimator is undetermined. The remedy is to penalize
roughness of $g$, measured by a nonnegative operator $\mathcal{W}\in
\mathcal{B}(\mathbb{H})$ (e.g.\ $\langle g,\mathcal{W}g\rangle_{\mathbb{H}}
=\int_0^1|g^{(m)}(t)|^2\,dt$ for $\mathcal{W}$ the projection onto
$\mathbb{H}_1$ of the Sobolev decomposition of Definition
\ref{def:sobolev}).

\begin{theorem}[Representer theorem]
\label{thm:representer}
Assume $\mathcal{T}^*\mathcal{T}+\eta\mathcal{W}$ is invertible. For
$\eta\in(0,\infty)$, the unique minimizer over $g\in\mathbb{H}$ of
\begin{equation}
f(g;\mathcal{T},\mathcal{W},Y,\eta) := \|Y-\mathcal{T}g\|_{\mathbb{Y}}^2 +
\eta\langle g,\mathcal{W}g\rangle_{\mathbb{H}}
\label{eq:penalized_criterion}
\end{equation}
is
\begin{equation}
m_\eta := (\mathcal{T}^*\mathcal{T}+\eta\mathcal{W})^{-1}\mathcal{T}^*Y.
\label{eq:regularization_estimator}
\end{equation}
\end{theorem}

\begin{proof}
Write a candidate solution as $\widetilde{g}=m_\eta+g$. Expanding,
$f(\widetilde{g};\mathcal{T},\mathcal{W},Y,\eta) = \|Y\|_{\mathbb{Y}}^2 +
\langle\mathcal{T}^*\mathcal{T}\widetilde{g},\widetilde{g}\rangle_{
\mathbb{H}} - 2\langle Y,\mathcal{T}\widetilde{g}\rangle_{\mathbb{Y}} +
\eta\langle\mathcal{W}\widetilde{g},\widetilde{g}\rangle_{\mathbb{H}}$.
Since $\mathcal{T}^*Y=(\mathcal{T}^*\mathcal{T}+\eta\mathcal{W})m_\eta$ by
\eqref{eq:regularization_estimator}, $\langle Y,\mathcal{T}\widetilde{g}
\rangle_{\mathbb{Y}} = \langle(\mathcal{T}^*\mathcal{T}+\eta\mathcal{W})
m_\eta,\widetilde{g}\rangle_{\mathbb{H}}$, so
\[
f(\widetilde{g};\mathcal{T},\mathcal{W},Y,\eta) = \|Y\|_{\mathbb{Y}}^2 +
\big\langle(\mathcal{T}^*\mathcal{T}+\eta\mathcal{W})(\widetilde{g}-2m_\eta),
\widetilde{g}\big\rangle_{\mathbb{H}}.
\]
Writing $\widetilde{g}-2m_\eta = (\widetilde{g}-m_\eta)-m_\eta$ and
expanding the inner product,
\[
\big\langle(\mathcal{T}^*\mathcal{T}+\eta\mathcal{W})(\widetilde{g}-2m_\eta),
\widetilde{g}\big\rangle_{\mathbb{H}} = \big\langle(\mathcal{T}^*\mathcal{T}
+\eta\mathcal{W})g,g\big\rangle_{\mathbb{H}} - \big\langle(\mathcal{T}^*
\mathcal{T}+\eta\mathcal{W})m_\eta,m_\eta\big\rangle_{\mathbb{H}}
\]
(using symmetry of $\mathcal{T}^*\mathcal{T}+\eta\mathcal{W}$ to combine
cross terms). Since $\mathcal{T}^*\mathcal{T}+\eta\mathcal{W}$ is positive
definite (it is invertible and nonnegative, being a sum of nonnegative
operators), the first term, and hence
$f(\widetilde{g};\mathcal{T},\mathcal{W},Y,\eta)$, is uniquely minimized at
$g=0$, i.e.\ $\widetilde{g}=m_\eta$.
\end{proof}

Equation~\eqref{eq:regularization_estimator} is precisely the abstract form
of the B-spline smoothing-spline criterion
\eqref{eq:smoothing_spline}--\eqref{eq:sigma_est} used in
Section~\ref{subsec:bspline}: there, $\mathbb{Y}=\mathbb{R}^n$,
$\mathbb{H}=\mathbb{W}^m(\mathcal{T})$, $\mathcal{W}$ penalizes
$\int_{\mathcal{T}}[X^{(m)}(t)]^2\,dt$, and $\eta=\lambda$ is chosen
data-adaptively per country via $\widehat\sigma$ rather than by the
cross-validation methods of Section~\ref{subsec:reg_param_selection}
below. The behavior of $m_\eta$ at the extremes $\eta\to 0$ (which recovers
the minimum-norm least-squares solution of $\mathcal{T}g=Y$, via the
Moore--Penrose pseudoinverse $\mathcal{T}^\dagger$) and $\eta\to\infty$
(which forces $m_\eta$ into $\mathrm{Ker}(\mathcal{W})$, the ``ideal'',
unpenalized model class) is established in \cite{hsingeubank2015}, Theorems
6.2.2--6.2.3; we omit the proofs here, as they rely on the theory of
generalized (Moore--Penrose) operator inverses developed in that source's
Section~3.5, which falls outside the scope of the preliminaries of
Section~\ref{sec:preliminaries} and is not needed for the fixed-$\eta$,
data-adaptively-chosen smoothing spline actually used in this thesis.

\subsection{Bias and Variance}
\label{subsec:bias_variance}

\begin{theorem}
\label{thm:bias_variance}
Let $m_\eta$ be as in \eqref{eq:regularization_estimator} for the model
\eqref{eq:functional_linear_model}, and suppose a future, independent
observation $Y_{\mathrm{new}}=\mathcal{T}m+\varepsilon_{\mathrm{new}}$ is to
be predicted by $\mathcal{T}m_\eta$, with prediction risk $\mathrm{Risk}
(\eta) = \mathrm{Var}(\eta)+\mathrm{Bias}^2(\eta)$ for
\begin{equation}
\mathrm{Bias}^2(\eta) := \|\mathcal{T}(\mathbb{E}m_\eta-m)\|_{\mathbb{Y}}^2,
\qquad
\mathrm{Var}(\eta) := \mathbb{E}\|\mathcal{T}(m_\eta-\mathbb{E}m_\eta)
\|_{\mathbb{Y}}^2.
\label{eq:bias_variance_def}
\end{equation}
Then $\mathrm{Bias}^2(\eta)\leq\eta\langle m,\mathcal{W}m\rangle_{\mathbb{H}}$
and $\mathrm{Var}(\eta) = \mathbb{E}\|\mathcal{T}(\mathcal{T}^*\mathcal{T}+
\eta\mathcal{W})^{-1}\mathcal{T}^*\varepsilon\|_{\mathbb{Y}}^2$.
\end{theorem}

\begin{proof}
$\mathbb{E}m_\eta = (\mathcal{T}^*\mathcal{T}+\eta\mathcal{W})^{-1}
\mathcal{T}^*\mathcal{T}m$ is, by Theorem~\ref{thm:representer} applied
with $Y$ replaced by $\mathcal{T}m$, the minimizer over $g\in\mathbb{H}$ of
$\|\mathcal{T}m-\mathcal{T}g\|_{\mathbb{Y}}^2+\eta\langle g,\mathcal{W}g
\rangle_{\mathbb{H}}$. Hence $\|\mathcal{T}(\mathbb{E}m_\eta-m)\|_{
\mathbb{Y}}^2 \leq \|\mathcal{T}(\mathbb{E}m_\eta-m)\|_{\mathbb{Y}}^2+
\eta\langle\mathbb{E}m_\eta,\mathcal{W}\mathbb{E}m_\eta\rangle_{\mathbb{H}}
\leq \|\mathcal{T}(g-m)\|_{\mathbb{Y}}^2+\eta\langle g,\mathcal{W}g
\rangle_{\mathbb{H}}$ for every $g$; taking $g=m$ gives the bias bound. The
variance formula follows since $m_\eta-\mathbb{E}m_\eta = (\mathcal{T}^*
\mathcal{T}+\eta\mathcal{W})^{-1}\mathcal{T}^*\varepsilon$ directly from
\eqref{eq:regularization_estimator} and linearity of expectation.
\end{proof}

\begin{corollary}
\label{cor:variance_eigenvalues}
Let $\mathbb{Y}=\mathbb{R}^n$ with $\|Y\|_{\mathbb{Y}}^2=n^{-1}\sum_jY_j^2$,
and suppose $\varepsilon$ has uncorrelated entries with mean $0$ and
variance $\sigma^2$, and $\mathcal{W}$ is invertible. Then
\begin{equation}
\mathrm{Var}(\eta) = \frac{\sigma^2}{n}\sum_{j=1}^n\Big(\frac{\gamma_j}{
\gamma_j+\eta}\Big)^2,
\label{eq:variance_eigenvalues}
\end{equation}
where $\gamma_1\geq\cdots\geq\gamma_n\geq 0$ are the eigenvalues of the
compact operator $\mathcal{W}^{-1/2}\mathcal{T}^*\mathcal{T}\mathcal{W}^{
-1/2}$.
\end{corollary}

\begin{proof}
Since $\mathrm{Im}(\mathcal{T})$ is finite-dimensional, $\mathcal{T}$, and
hence $\widetilde{\mathcal{T}}:=\mathcal{T}\mathcal{W}^{-1/2}$, is compact
(Theorem~\ref{thm:compact_props}(2)), with at most $n$ singular values
$\gamma_1\geq\cdots\geq\gamma_n\geq 0$ (which are the eigenvalues of
$\widetilde{\mathcal{T}}^*\widetilde{\mathcal{T}}=\mathcal{W}^{-1/2}
\mathcal{T}^*\mathcal{T}\mathcal{W}^{-1/2}$, by Theorem~\ref{thm:svd}).
Writing $\mathcal{T}(\mathcal{T}^*\mathcal{T}+\eta\mathcal{W})^{-1}
\mathcal{T}^* = \widetilde{\mathcal{T}}(\widetilde{\mathcal{T}}^*
\widetilde{\mathcal{T}}+\eta I)^{-1}\widetilde{\mathcal{T}}^*$ (substituting
$\mathcal{T}=\widetilde{\mathcal{T}}\mathcal{W}^{1/2}$) and using
$\mathbb{E}\langle\varepsilon,e\rangle_{\mathbb{Y}}^2=\sigma^2e^Te/n^2$ for
any $e\in\mathbb{R}^n$ together with the SVD of $\widetilde{\mathcal{T}}$
in Theorem~\ref{thm:bias_variance}'s variance formula gives
\eqref{eq:variance_eigenvalues}.
\end{proof}

\subsection{A Computational Formula}
\label{subsec:computational_formula}

\begin{theorem}
\label{thm:computational_formula}
Let $\mathbb{Y}=\mathbb{R}^n$, $\mathbb{H}=\mathbb{H}_0\oplus\mathbb{H}_1$
with $\dim(\mathbb{H}_0)=q<\infty$, basis $\{\phi_1,\dots,\phi_q\}$ of
$\mathbb{H}_0$, and $\mathcal{W}=\mathcal{P}_1$ the projection onto
$\mathbb{H}_1$. Write $\mathcal{T}g=(\mathcal{T}_1g,\dots,\mathcal{T}_ng)^T$
for bounded linear functionals $\mathcal{T}_i$ with representers $\tau_i$
(i.e.\ $\mathcal{T}_ig=\langle\tau_i,g\rangle_{\mathbb{H}}$), and let
$\xi_i:=\mathcal{P}_1\tau_i$. Define $\Phi:=\{\mathcal{T}_i\phi_j\}_{i,j}$,
$\Psi:=\{\mathcal{T}_i\xi_j\}_{i,j}=\{\langle\xi_i,\xi_j\rangle_{\mathbb{H}}
\}_{i,j}$, $\mathcal{U}:=(\Phi,\Psi)$, and $\mathcal{V}:=
\left(\begin{smallmatrix}0_{q\times q}&0_{q\times n}\\0_{n\times q}&
\Psi_{n\times n}\end{smallmatrix}\right)$. If $\mathcal{U}^T\mathcal{U}+
\eta\mathcal{V}$ is invertible, the minimizer of
\eqref{eq:penalized_criterion} is $m_\eta = \sum_{j=1}^q\widehat{c}_j
\phi_j+\sum_{i=1}^n\widehat{b}_i\xi_i$, where
\begin{equation}
\binom{\widehat{c}}{\widehat{b}} = (\mathcal{U}^T\mathcal{U}+\eta
\mathcal{V})^{-1}\mathcal{U}^TY.
\label{eq:computational_formula}
\end{equation}
\end{theorem}

\begin{proof}
Write $g=\sum_jc_j\phi_j+\sum_ib_i\xi_i+h$ for $h\perp\mathrm{span}\{
\phi_1,\dots,\phi_q,\xi_1,\dots,\xi_n\}$. Since $\tau_i-\xi_i\in
\mathbb{H}_0=\mathrm{span}\{\phi_j\}$, $\mathcal{T}_ih=\langle\tau_i,h
\rangle_{\mathbb{H}} = \langle\tau_i-\xi_i,h\rangle_{\mathbb{H}}+\langle
\xi_i,h\rangle_{\mathbb{H}}=0$, so $\mathcal{T}_ig$ does not depend on $h$,
while $\langle g,\mathcal{P}_1g\rangle_{\mathbb{H}} = \|\mathcal{P}_1g\|^2$
strictly increases with $\|h\|$ whenever $h\neq 0$ has a nonzero
$\mathbb{H}_1$-component (since $h\perp\xi_i$ does not preclude $h\in
\mathbb{H}_1$ in general, but any such component only adds to
$\langle g,\mathcal{P}_1g\rangle_{\mathbb{H}}$ without changing $\mathcal{T}
g$); hence the minimizer has $h=0$, i.e.\ $m_\eta\in\mathrm{span}\{\phi_j,
\xi_i\}$. For such $g$, $\mathcal{T}g=\Phi c+\Psi b$ and $\langle g,
\mathcal{P}_1g\rangle_{\mathbb{H}}=b^T\Psi b$ (since $\mathcal{P}_1\phi_j=0$
and $\mathcal{P}_1\xi_i=\xi_i$), so
\eqref{eq:penalized_criterion} becomes $\|Y-\mathcal{U}\binom{c}{b}\|_{
\mathbb{Y}}^2+\eta\binom{c}{b}^T\mathcal{V}\binom{c}{b}$, and Theorem
\ref{thm:representer} (applied in $\mathbb{R}^{q+n}$ with the identity
map replaced by $\mathcal{U}$) gives \eqref{eq:computational_formula}.
\end{proof}

\begin{theorem}
\label{thm:rkhs_computational}
If $\mathbb{H}_0=\mathbb{H}(K_0)$, $\mathbb{H}_1=\mathbb{H}(K_1)$,
$\mathbb{H}=\mathbb{H}(K)$ for reproducing kernels with $K=K_0+K_1$ and
$\mathbb{H}_0\cap\mathbb{H}_1=\{0\}$, then $\mathbb{H}=\mathbb{H}_0\oplus
\mathbb{H}_1$ and, in the notation of Theorem
\ref{thm:computational_formula}, $\tau_i(t)=\mathcal{T}_i^{(\cdot)}K(t,
\cdot)$ and $\xi_i(t)=\mathcal{T}_i^{(\cdot)}K_1(t,\cdot)$ (applying
$\mathcal{T}_i$ to the indicated argument).
\end{theorem}

\begin{proof}
$\mathbb{H}=\mathbb{H}_0\oplus\mathbb{H}_1$ is the kernel-sum decomposition
of Definition~\ref{def:rkhs}. By the reproducing
property, $\tau_i(t)=\langle\tau_i,K(t,\cdot)\rangle_{\mathbb{H}} =
\mathcal{T}_i^{(\cdot)}K(t,\cdot)$. Similarly, $\xi_i(t) = \langle\xi_i,
K(t,\cdot)\rangle_{\mathbb{H}} = \langle\mathcal{P}_1\tau_i,K(t,\cdot)
\rangle_{\mathbb{H}} = \langle\tau_i,\mathcal{P}_1K(t,\cdot)\rangle_{
\mathbb{H}} = \langle\tau_i,K_1(t,\cdot)\rangle_{\mathbb{H}} =
\mathcal{T}_i^{(\cdot)}K_1(t,\cdot)$, using that $\mathcal{P}_1$ projects
$K(t,\cdot)=K_0(t,\cdot)+K_1(t,\cdot)$ onto its $\mathbb{H}_1$-component.
\end{proof}

Theorem~\ref{thm:rkhs_computational} applies directly to $\mathbb{H}=
\mathbb{W}^q[0,1]$ with $K_0,K_1$ as in Definition~\ref{def:sobolev} (i.e.\
Theorem 2.8.1 of \cite{hsingeubank2015}).

\subsection{Regularization Parameter Selection}
\label{subsec:reg_param_selection}

We now specialize to $\mathbb{Y}=\mathbb{R}^n$, $\mathbb{H}=\mathbb{H}_0
\oplus\mathbb{H}_1$ (finite-dimensional $\mathbb{H}_0$) of Theorem
\ref{thm:computational_formula}, and describe three data-adaptive methods
for choosing $\eta$ that do not require the closed-form, noise-level-based
rule used in Section~\ref{subsec:bspline}.

Let $m_\eta^{[k]}$ minimize $n^{-1}\sum_{i\neq k}(Y_i-\mathcal{T}_ig)^2+
\eta\langle g,\mathcal{P}_1g\rangle_{\mathbb{H}}$ (leaving out
observation $k$), and define ordinary cross-validation as $\mathrm{CV}
(\eta) := n^{-1}\sum_{k=1}^n(Y_k-\mathcal{T}_km_\eta^{[k]})^2$, chosen to
minimize $\mathrm{CV}(\eta)$. Direct evaluation requires refitting $n$
times per candidate $\eta$; the following device avoids this.

\begin{lemma}[Wahba's leave-one-out lemma]
\label{lem:leave_one_out}
For $y\in\mathbb{R}$, let $m_\eta[k,y]$ minimize $n^{-1}(y-\mathcal{T}_kg)^2
+n^{-1}\sum_{i\neq k}(Y_i-\mathcal{T}_ig)^2+\eta\langle g,\mathcal{P}_1g
\rangle_{\mathbb{H}}$ (so $m_\eta=m_\eta[k,Y_k]$). Then
$m_\eta[k,\mathcal{T}_km_\eta^{[k]}] = m_\eta^{[k]}$.
\end{lemma}

\begin{proof}
For any $g\in\mathbb{H}$, using that $m_\eta^{[k]}$ minimizes the
leave-one-out criterion, and abbreviating $z:=\mathcal{T}_km_\eta^{[k]}$,
\begin{align*}
n^{-1}(\mathcal{T}_km_\eta^{[k]}-z)^2 &+ n^{-1}\sum_{i\neq k}(Y_i-
\mathcal{T}_im_\eta^{[k]})^2+\eta\langle m_\eta^{[k]},\mathcal{P}_1
m_\eta^{[k]}\rangle_{\mathbb{H}} \\
&= n^{-1}\sum_{i\neq k}(Y_i-\mathcal{T}_im_\eta^{[k]})^2+\eta\langle
m_\eta^{[k]},\mathcal{P}_1m_\eta^{[k]}\rangle_{\mathbb{H}} \\
&\leq n^{-1}\sum_{i\neq k}(Y_i-\mathcal{T}_ig)^2+\eta\langle g,\mathcal{P}_1
g\rangle_{\mathbb{H}} \\
&\leq n^{-1}(z-\mathcal{T}_kg)^2+n^{-1}\sum_{i\neq k}(Y_i-\mathcal{T}_ig)^2
+\eta\langle g,\mathcal{P}_1g\rangle_{\mathbb{H}},
\end{align*}
i.e.\ $m_\eta^{[k]}$ minimizes the criterion defining $m_\eta[k,z]$, so
$m_\eta[k,z]=m_\eta^{[k]}$.
\end{proof}

Define the ``leverage'' $a_{kk}(\eta) := (\mathcal{T}_km_\eta-\mathcal{T}_k
m_\eta^{[k]})/(Y_k-\mathcal{T}_km_\eta^{[k]})$, so that $Y_k-\mathcal{T}_k
m_\eta^{[k]} = (Y_k-\mathcal{T}_km_\eta)/(1-a_{kk}(\eta))$ and
\begin{equation}
\mathrm{CV}(\eta) = n^{-1}\sum_{k=1}^n\frac{(Y_k-\mathcal{T}_km_\eta)^2}
{[1-a_{kk}(\eta)]^2}.
\label{eq:cv_formula}
\end{equation}
By Lemma~\ref{lem:leave_one_out}, writing $Y_k'=\mathcal{T}_km_\eta^{[k]}$,
$a_{kk}(\eta) = \{\mathcal{T}_km_\eta[k,Y_k]-\mathcal{T}_km_\eta[k,Y_k']\}/
(Y_k-Y_k')$; since $m_\eta[k,y]$ is, by
Theorem~\ref{thm:computational_formula} applied with the $k$th response
replaced by $y$, linear in $y$, this ratio equals the derivative
$\partial\mathcal{T}_km_\eta[k,y]/\partial y|_{y=Y_k}$, i.e.\ the $k$th
diagonal entry of the ``hat matrix'' $\mathcal{A}(\eta):=\mathcal{U}
(\mathcal{U}^T\mathcal{U}+\eta\mathcal{V})^{-1}\mathcal{U}^T$ satisfying
$(\mathcal{T}_1m_\eta,\dots,\mathcal{T}_nm_\eta)^T=\mathcal{A}(\eta)Y$. So
$\mathrm{CV}(\eta)$ is computable directly from the residuals $Y-
\mathcal{T}m_\eta$ and $\mathrm{diag}(\mathcal{A}(\eta))$, without
refitting.

Generalized cross-validation replaces the individual leverages $a_{kk}
(\eta)$ in \eqref{eq:cv_formula} by their average, minimizing
\begin{equation}
\mathrm{GCV}(\eta) := \frac{\|Y-\mathcal{T}m_\eta\|_{\mathbb{Y}}^2}{
\big(1-n^{-1}\mathrm{trace}\,\mathcal{A}(\eta)\big)^2}.
\label{eq:gcv_formula}
\end{equation}
Suppose $\varepsilon$ has uncorrelated entries, mean $0$, variance
$\sigma^2$, and let $\mu:=\mathbb{E}Y=(\mathcal{T}_1m,\dots,\mathcal{T}_nm)
^T$. The prediction risk of $\mathcal{A}(\eta)Y$ for $\mu$ is $\mathrm{Risk}
(\eta):=n^{-1}\mathbb{E}\|\mathcal{A}(\eta)Y-\mu\|_2^2 = n^{-1}\mu^T
(I-\mathcal{A}(\eta))^T(I-\mathcal{A}(\eta))\mu + \sigma^2n^{-1}
\mathrm{trace}\,\mathcal{A}(\eta)\mathcal{A}(\eta)^T$ (bias-variance
decomposition, as in Theorem~\ref{thm:bias_variance} specialized to
$\mathbb{R}^n$). Writing $\mathrm{MSE}(\eta):=\|(I-\mathcal{A}(\eta))Y\|_{
\mathbb{Y}}^2$, so that $\mathrm{GCV}(\eta)=\mathrm{MSE}(\eta)/(1-n^{-1}
\mathrm{trace}\,\mathcal{A}(\eta))^2$, an identical bias-variance expansion
of $\mathbb{E}[\mathrm{MSE}(\eta)]$ gives
\begin{equation}
\mathbb{E}[\mathrm{MSE}(\eta)] = \mathrm{Risk}(\eta) + \sigma^2 -
\frac{2\sigma^2}{n}\mathrm{trace}\,\mathcal{A}(\eta).
\label{eq:mse_risk_identity}
\end{equation}

\begin{theorem}[GCV Theorem]
\label{thm:gcv}
Let $\tau_j(\eta):=\mathrm{trace}\,\mathcal{A}^j(\eta)/n$, $j=1,2$, and
$g(\eta):=\{2\tau_1(\eta)+\tau_1(\eta)^2/\tau_2(\eta)\}/(1-\tau_1(\eta))^2$.
Then
\begin{equation}
\frac{|\mathbb{E}\,\mathrm{GCV}(\eta)-\sigma^2-\mathrm{Risk}(\eta)|}{
\mathrm{Risk}(\eta)} \leq g(\eta).
\label{eq:gcv_theorem}
\end{equation}
\end{theorem}

\begin{proof}
From \eqref{eq:mse_risk_identity} and $\mathrm{GCV}(\eta)=\mathbb{E}[
\mathrm{MSE}(\eta)]/(1-\tau_1(\eta))^2$ (in expectation), algebra gives
\[
\mathbb{E}\,\mathrm{GCV}(\eta)-\sigma^2-\mathrm{Risk}(\eta) =
\frac{\tau_1(\eta)(2-\tau_1(\eta))}{(1-\tau_1(\eta))^2}\mathrm{Risk}(\eta)
- \frac{\sigma^2\tau_1(\eta)^2}{(1-\tau_1(\eta))^2}.
\]
Dividing by $\mathrm{Risk}(\eta)$ and applying the triangle inequality,
\[
\frac{|\mathbb{E}\,\mathrm{GCV}(\eta)-\sigma^2-\mathrm{Risk}(\eta)|}{
\mathrm{Risk}(\eta)} \leq \frac{\tau_1(\eta)(2-\tau_1(\eta))}{(1-\tau_1
(\eta))^2} + \frac{\tau_1(\eta)^2/\tau_2(\eta)}{(1-\tau_1(\eta))^2} =
g(\eta),
\]
using $\sigma^2/\mathrm{Risk}(\eta)\leq 1/\tau_2(\eta)$ (from
\eqref{eq:mse_risk_identity} and nonnegativity of the bias term, since
$\mathrm{trace}\,\mathcal{A}^2(\eta) \le n\sigma^2/\mathrm{Risk}(\eta)\cdot$
after rearrangement -- more directly, $\sigma^2\mathrm{trace}\,
\mathcal{A}^2(\eta)/n \leq \mathrm{Risk}(\eta)$ is the variance component of
$\mathrm{Risk}(\eta)$ in $\mathbb{R}^n$).
\end{proof}

When $g(\eta)$ is small, $\mathrm{GCV}(\eta)$ is thus a near-unbiased
estimator of $\sigma^2+\mathrm{Risk}(\eta)$. A related, exactly unbiased
estimator of $\mathrm{Risk}(\eta)$, from \eqref{eq:mse_risk_identity}, is
Mallows' criterion $C_L(\eta):=\mathrm{MSE}(\eta)+2\sigma^2n^{-1}
\mathrm{trace}\,\mathcal{A}(\eta)$, minimized in place of
$\mathrm{GCV}(\eta)$ when an estimate of $\sigma^2$ is available; GCV
avoids this requirement at the cost of a small bias governed by $g(\eta)$.

\subsection{Splines}
\label{subsec:splines_theory}

We close this section with the theory of polynomial splines, which
resolves the penalized least-squares problem
\eqref{eq:penalized_criterion} explicitly for the roughness penalty
$\langle g,\mathcal{W}g\rangle_{\mathbb{H}}=\int_0^1|g^{(q)}(t)|^2\,dt$ on
$\mathbb{H}=\mathbb{W}^q[0,1]$, and underlies the B-spline basis of
Section~\ref{subsec:bspline}.

\begin{definition}
\label{def:spline}
A \emph{spline of order $q\geq 1$} with knots $0<t_1<\cdots<t_J<1$ is
\begin{equation}
g(t) = \sum_{i=0}^{q-1}\theta_it^i + \sum_{j=1}^J\delta_j(t-t_j)_+^{q-1},
\qquad \theta_i,\delta_j\in\mathbb{R},
\label{eq:spline_def}
\end{equation}
where $x_+:=\max(x,0)$. Equivalently, $g$ is a piecewise polynomial of
order $q$ on each $[t_i,t_{i+1})$, with $q-2$ continuous derivatives (for
$q\geq 2$) and $(q-1)$-st derivative a step function jumping at
$t_1,\dots,t_J$. The functions $1,t,\dots,t^{q-1},(t-t_1)_+^{q-1},\dots,
(t-t_J)_+^{q-1}$ are linearly independent, so the space
$\mathcal{S}_q(t_1,\dots,t_J)$ of such splines has dimension $q+J$.
\end{definition}

For computation, an equivalent, nearly-orthogonal basis is preferable to
the truncated-power basis of \eqref{eq:spline_def}: the B-spline basis,
constructed via divided differences.

\begin{definition}
\label{def:divided_difference}
The $q$th-order \emph{divided difference} of $g$ at $t_i,\dots,t_{i+q}$ is
defined recursively by $[t_i]g:=g(t_i)$ and
\begin{equation}
[t_i,\dots,t_{i+q}]g := \frac{[t_{i+1},\dots,t_{i+q}]g -
[t_i,\dots,t_{i+q-1}]g}{t_{i+q}-t_i}.
\label{eq:divided_difference}
\end{equation}
\end{definition}

\begin{theorem}
\label{thm:divided_difference_interp}
Let $p_q(x):=\sum_{i=1}^qg(t_i)\ell_i(x)$ for $\ell_i(x):=\prod_{j\neq i}
(x-t_j)/(t_i-t_j)$. Then: (1) $p_q$ is the unique polynomial of order $q$
agreeing with $g$ at $t_1,\dots,t_q$; (2) for $q=1,2,\dots$, the coefficient
of $x^q$ in $p_{q+1}$ (the interpolant of $g$ at $t_i,\dots,t_{i+q}$) is
$[t_i,\dots,t_{i+q}]g$.
\end{theorem}

\begin{proof}
(1) $\ell_i$ has order $q$ and vanishes at every $t_j$, $j\neq i$, except
$\ell_i(t_i)=1$; uniqueness of interpolating polynomials of order $q$
through $q$ points is standard.

(2) By induction on $q$; trivial for $q=1$. Let $p_q,\widetilde{p}_q$
interpolate $g$ at $t_i,\dots,t_{i+q-1}$ and $t_{i+1},\dots,t_{i+q}$
respectively. Then $p(x):=\frac{x-t_i}{t_{i+q}-t_i}\widetilde{p}_q(x)+
\frac{t_{i+q}-x}{t_{i+q}-t_i}p_q(x)$ has order $q+1$ and agrees with $g$ at
$t_i,\dots,t_{i+q}$ (direct check at each node), so $p_{q+1}=p$ by
uniqueness, and the coefficient of $x^q$ in $p_{q+1}$ is $\{[t_{i+1},
\dots,t_{i+q}]g-[t_i,\dots,t_{i+q-1}]g\}/(t_{i+q}-t_i) = [t_i,\dots,
t_{i+q}]g$ by the induction hypothesis and \eqref{eq:divided_difference}.
\end{proof}

\begin{corollary}
\label{cor:divided_difference_poly}
If $g$ is a polynomial of order $q$ on $[t_i,t_{i+q}]$, then
$[t_i,\dots,t_{i+q}]g=0$.
\end{corollary}

\begin{proof}
With $p_l$ as in Theorem~\ref{thm:divided_difference_interp}, $p_l=g$ for
$l\geq q$ (both interpolate $g$, which already has order $\leq q\leq l$),
so $[t_i,\dots,t_{i+q}]g$, the coefficient of $x^q$ in $p_{q+1}=g$, is $0$
since $g$ has degree $<q$.
\end{proof}

Introduce phantom knots $t_{-(q-1)}<\cdots<t_{-1}\leq 0=t_0$ and
$1=t_{J+1}\leq t_{J+2}<\cdots<t_{J+q}$ (arbitrary, e.g.\ all equal to $0$
or $1$ respectively), and define, for $i=-(q-1),\dots,J$, the \emph{B-spline}
\begin{equation}
N_{iq}(t) := (t_{i+q}-t_i)[t_i,\dots,t_{i+q}](\cdot-t)_+^{q-1}.
\label{eq:bspline_def}
\end{equation}

\begin{theorem}
\label{thm:bspline_properties}
The B-splines $\{N_{iq}\}_{i=-(q-1)}^J$ satisfy: (1) $N_{iq}$ is a
polynomial of order $q$ on each $(t_i,t_{i+1})$; (2) $N_{iq}(t)=0$ for
$t\notin[t_i,t_{i+q}]$ (local support); (3) with $N_{i1}:=I_{[t_i,t_{i+1})}
$, $N_{iq}(t) = \frac{t-t_i}{t_{i+q-1}-t_i}N_{i(q-1)}(t) +
\frac{t_{i+q}-t}{t_{i+q}-t_{i+1}}N_{(i+1)(q-1)}(t)$.
\end{theorem}

\begin{proof}
(1) is immediate from \eqref{eq:bspline_def}. For (2): $N_{iq}(t)=0$ for
$t>t_{i+q}$ since every term $(t_{i+r}-t)_+^{q-1}$, $r=0,\dots,q$, then
vanishes. For $t<t_i$, every $(t_{i+r}-t)_+^{q-1}=(t_{i+r}-t)^{q-1}$, so
$N_{iq}(t) = (t_{i+q}-t_i)[t_i,\dots,t_{i+q}](\cdot-t)^{q-1}$, and since
$(\cdot-t)^{q-1}$ is a polynomial of order $q$ on $[t_i,t_{i+q}]$, this
vanishes by Corollary~\ref{cor:divided_difference_poly}. (Part (3), the
recursion, is a direct algebraic identity relating divided differences at
nested knot sets; we omit the routine verification, cf.\
\cite{hsingeubank2015}, Theorem 6.6.5.)
\end{proof}

By the local support property, $N_{iq}$ and $N_{jq}$ have disjoint supports
whenever $|i-j|>q-1$, so $\langle N_{iq},N_{jq}\rangle_2=0$ in that case
(``near orthogonality''); this, together with the fact that the $N_{iq}$
therefore have full rank on $\mathrm{span}\{N_{iq}\}$, shows they form a
basis for $\mathcal{S}_q(t_1,\dots,t_J)$, and this is the basis
\texttt{scipy.interpolate.UnivariateSpline} constructs internally for the
B-spline fits of Section~\ref{subsec:bspline}.

\begin{definition}
\label{def:natural_spline}
A spline $g$ of order $2q$ with knots $t_1<\cdots<t_J$ is a \emph{natural
spline} (notation $\mathcal{N}^{2q}(t_1,\dots,t_J)$) if $g$ reduces to a
polynomial of order $q$ outside $[t_1,t_J]$; equivalently, $g^{(j)}(0)=
g^{(j)}(1)=0$ for $j=q,\dots,2q-1$ (the \emph{natural boundary
conditions}), which impose $2q$ linear constraints, giving
$\dim\mathcal{N}^{2q}(t_1,\dots,t_J) = J$.
\end{definition}

\begin{theorem}
\label{thm:natural_spline_interp}
For any $0\leq t_1<\cdots<t_J\leq 1$ and $z_1,\dots,z_J\in\mathbb{R}$,
there is a unique natural spline $g\in\mathcal{N}^{2q}(t_1,\dots,t_J)$ with
$g(t_j)=z_j$ for all $j$.
\end{theorem}

\begin{proof}
This is a linear system: the $J$-dimensional space $\mathcal{N}^{2q}
(t_1,\dots,t_J)$ (Definition~\ref{def:natural_spline}) is mapped to
$\mathbb{R}^J$ by evaluation at $t_1,\dots,t_J$; existence and uniqueness
of the interpolant is equivalent to this linear map being a bijection,
which holds since a natural spline vanishing at all $t_j$ must (being a
polynomial of order $q\leq J$ on each subinterval, pieced together with
$q-2$ continuous derivatives, and vanishing together with its top $q$
derivatives at the boundary) vanish identically -- the standard, purely
linear-algebraic verification for natural spline interpolation
(\cite{hsingeubank2015}, Theorem 6.6.7; de Boor, 1978).
\end{proof}

The following optimality property of the natural interpolating spline is
the key structural fact that reduces the (infinite-dimensional) smoothing
problem in $\mathbb{W}^q[0,1]$ to a finite-dimensional one.

\begin{theorem}
\label{thm:natural_spline_optimal}
Let $\widetilde{g}\in\mathcal{N}^{2q}(t_1,\dots,t_J)$ interpolate
$z_1,\dots,z_J$ (with $0<t_1<\cdots<t_J<1$). Then for every $g\in
\mathbb{W}^q[0,1]$ that also interpolates the $z_j$,
\begin{equation}
\int_0^1|\widetilde{g}^{(q)}(t)|^2\,dt \leq \int_0^1|g^{(q)}(t)|^2\,dt,
\label{eq:natural_spline_optimal}
\end{equation}
with equality iff $g=\widetilde{g}+h$ for $h$ a $q$th-order polynomial
vanishing at $t_1,\dots,t_J$ (so $h=0$ identically if $J\geq q$).
\end{theorem}

\begin{proof}
Let $h:=g-\widetilde{g}$, so $h(t_j)=0$ for all $j$. Repeated integration
by parts, using the natural boundary conditions on $\widetilde{g}$ (which
kill the boundary terms at $0,1$) and that $\widetilde{g}^{(2q-1)}$ is
piecewise constant (being the $(q-1)$-st derivative of the order-$q$
piecewise polynomial $\widetilde{g}^{(q)}$... more precisely, constant on
each $(t_j,t_{j+1})$ since $\widetilde{g}$ has order $2q$), gives
\begin{align}
\int_0^1\widetilde{g}^{(q)}h^{(q)}\,dt &= -\int_0^1\widetilde{g}^{(q+1)}
h^{(q-1)}\,dt = \cdots = (-1)^{q-1}\int_0^1\widetilde{g}^{(2q-1)}h^{(1)}\,dt
\notag \\
&= (-1)^{q-1}\sum_{j=1}^{J-1}\widetilde{g}^{(2q-1)}(t_j^+)\,[h(t_{j+1})-
h(t_j)] = 0,
\label{eq:integration_by_parts}
\end{align}
using $h(t_j)=0$ for every $j$ in the last step and that $\widetilde{g}$
is a polynomial of order $q$ (so $\widetilde{g}^{(2q-1)}=0$) outside
$[t_1,t_J]$. Hence
\[
\int_0^1|g^{(q)}|^2\,dt = \int_0^1|\widetilde{g}^{(q)}|^2\,dt + 2\int_0^1
\widetilde{g}^{(q)}h^{(q)}\,dt + \int_0^1|h^{(q)}|^2\,dt =
\int_0^1|\widetilde{g}^{(q)}|^2\,dt + \int_0^1|h^{(q)}|^2\,dt \geq
\int_0^1|\widetilde{g}^{(q)}|^2\,dt,
\]
proving \eqref{eq:natural_spline_optimal}. Equality holds iff
$h^{(q)}=0$ a.e., i.e.\ (since $h\in\mathbb{W}^q[0,1]$, via the Taylor
representation of Section~\ref{sec:preliminaries}) iff $h$ is a
polynomial of order $q$; as $h(t_j)=0$ for $j=1,\dots,J$ and a nonzero
polynomial of order $q$ has at most $q-1$ roots, $h=0$ once $J\geq q$.
\end{proof}

\begin{theorem}
\label{thm:natural_spline_finite_dim}
Let $n>q$, $\eta\in(0,\infty)$, and $L:\mathbb{R}^n\to[0,\infty)$. If
\begin{equation}
L(g(t_1),\dots,g(t_n)) + \eta\int_0^1|g^{(q)}(t)|^2\,dt
\label{eq:general_penalized_criterion}
\end{equation}
has a minimizer over $g\in\mathbb{W}^q[0,1]$, that minimizer lies in
$\mathcal{N}^{2q}(t_1,\dots,t_n)$.
\end{theorem}

\begin{proof}
Let $\widehat{g}$ minimize \eqref{eq:general_penalized_criterion} and let
$\widetilde{g}\in\mathcal{N}^{2q}(t_1,\dots,t_n)$ interpolate $\widehat{g}$
at $t_1,\dots,t_n$ (Theorem~\ref{thm:natural_spline_interp}). Then
$L(\widehat{g}(t_1),\dots)=L(\widetilde{g}(t_1),\dots)$ by construction,
while Theorem~\ref{thm:natural_spline_optimal} gives $\int|\widetilde{g}
^{(q)}|^2\leq\int|\widehat{g}^{(q)}|^2$, so $\widetilde{g}$ attains a value
of \eqref{eq:general_penalized_criterion} no larger than $\widehat{g}$'s;
by minimality of $\widehat{g}$, equality holds and $\widetilde{g}$ is also
a minimizer, necessarily equal to $\widehat{g}$ once $n\geq q$ by the
equality case of Theorem~\ref{thm:natural_spline_optimal}.
\end{proof}

Theorem~\ref{thm:natural_spline_finite_dim} converts the smoothing spline
problem \eqref{eq:penalized_criterion} (with $\mathbb{H}=\mathbb{W}^q
[0,1]$, $\mathcal{T}g=(g(t_1),\dots,g(t_n))^T$, and $L$ the residual sum
of squares) from an infinite-dimensional to an $n$-dimensional problem:
writing $g_j\in\mathcal{N}^{2q}(t_1,\dots,t_n)$ for the natural spline
interpolating the $j$th unit vector at $t_1,\dots,t_n$, every $g\in
\mathcal{N}^{2q}(t_1,\dots,t_n)$ is $g=\sum_jg(t_j)g_j$, and minimizing
\eqref{eq:penalized_criterion} reduces to minimizing $\|Y-g\|_{\mathbb{Y}}
^2+\eta g^T\mathcal{W}g$ over $g=(g(t_1),\dots,g(t_n))^T\in\mathbb{R}^n$,
for $\mathcal{W}=\{\int_0^1g_i^{(q)}(t)g_j^{(q)}(t)\,dt\}_{i,j}$, with
explicit solution
\begin{equation}
m_\eta = (m_\eta(t_1),\dots,m_\eta(t_n))^T = (I+\eta\mathcal{W})^{-1}Y,
\qquad m_\eta(\cdot) = \sum_{j=1}^nm_\eta(t_j)g_j(\cdot).
\label{eq:smoothing_spline_solution}
\end{equation}
This natural spline of order $2q$ is the \emph{smoothing spline}; for
$q=2$ it is the cubic smoothing spline of
\eqref{eq:smoothing_spline}, which is exactly the B-spline smoothing
estimator of Section~\ref{subsec:bspline} (with $\eta=\lambda$).

We close with explicit rates for the bias and variance of the smoothing
spline, under equally spaced design points; these are the theoretical
counterpart of the RMSE--sparsity tradeoff explored empirically in
Section~\ref{subsec:sparsity_results}.

\begin{theorem}
\label{thm:smoothing_spline_rates}
Assume $\mathbb{E}[\varepsilon\varepsilon^T]=\sigma^2I$, $t_j=(2j-1)/2n$,
and let $m=(m(t_1),\dots,m(t_n))^T$. For every $n$ and $\eta>0$,
\begin{equation}
\|\mathbb{E}m_\eta-m\|_{\mathbb{Y}}^2 \leq \frac{\eta}{4}\big(m^Tm+m^T
\mathcal{W}m\big), \qquad
\mathbb{E}\|m_\eta-\mathbb{E}m_\eta\|_{\mathbb{Y}}^2 \leq \frac{q\sigma^2}
{n} + \frac{C\sigma^2}{n}\eta^{-1/(2q)},
\label{eq:smoothing_spline_rates}
\end{equation}
for a constant $C$ depending only on $q$.
\end{theorem}

\begin{proof}
Let $0=\gamma_1=\cdots=\gamma_q<\gamma_{q+1}\leq\cdots\leq\gamma_n$ be the
eigenvalues of $\mathcal{W}$ (known to satisfy $C_1(j-q)^{2q}\leq\gamma_j
\leq C_2(j-q)^{2q}$ for $j>q$ and constants $C_1,C_2$ depending only on
$q$, from the asymptotic theory of natural-spline Green's functions; cf.\
\cite{hsingeubank2015}, citing Utreras, 1983) with eigenvectors
$e_1,\dots,e_n$. From \eqref{eq:smoothing_spline_solution}, $\mathbb{E}
(m_\eta-m) = (I+\eta\mathcal{W})^{-1}m-m$, so
\[
\|\mathbb{E}(m_\eta-m)\|_{\mathbb{Y}}^2 = n^{-1}\sum_j\Big(\frac{\eta
\gamma_j}{1+\eta\gamma_j}\Big)^2(m^Te_j)^2 = n^{-1}\sum_j(1+\gamma_j)(m^T
e_j)^2\cdot(1+\gamma_j)^{-1}\Big(\frac{\eta\gamma_j}{1+\eta\gamma_j}\Big)^2.
\]
Since $\sum_j(1+\gamma_j)(m^Te_j)^2=m^Tm+m^T\mathcal{W}m$ and, substituting
$u=\gamma_j$, $\sup_{u\geq 0}(1+u)^{-1}(\eta u/(1+\eta u))^2 = \eta\sup_
{u\geq 0}\frac{u}{1+u}\cdot\frac{u}{(1+u)(\eta+u)} \le \eta\sup_{u\geq 0}
\frac{u}{(1+u)^2}=\eta/4$ (bounding the second factor by its value at the
maximizer of $u/(1+u)^2$), the bias bound follows.

For the variance, $\mathbb{E}\|m_\eta-\mathbb{E}m_\eta\|_{\mathbb{Y}}^2 =
n^{-1}\mathbb{E}[\varepsilon^T(I+\eta\mathcal{W})^{-2}\varepsilon] =
\sigma^2n^{-1}\sum_j(1+\eta\gamma_j)^{-2}$ (using $\mathbb{E}[\varepsilon
\varepsilon^T]=\sigma^2I$ and the eigendecomposition of $\mathcal{W}$).
The first $q$ terms (with $\gamma_j=0$) contribute $q\sigma^2/n$; the
remaining terms are bounded, using $\gamma_j\geq C_1(j-q)^{2q}$ and
approximating the sum by the integral $\int_0^\infty(1+C_1\eta x^{2q})^{
-2}\,dx = C\eta^{-1/(2q)}$ (substituting $u=C_1\eta x^{2q}$) for a constant
$C=C(q)$, giving the stated rate.
\end{proof}

Theorem~\ref{thm:smoothing_spline_rates} makes explicit the bias-variance
tradeoff invoked qualitatively in Section~\ref{subsec:bspline}: the bias
term grows linearly in $\eta$ while the variance term decays like
$\eta^{-1/(2q)}$, so that the mean-squared error $\mathrm{Bias}^2(\eta)+
\mathrm{Var}(\eta)$ is minimized at an intermediate $\eta=O(n^{2q/(2q+1)})$
-- the same $n^{-2q/(2q+1)}$ rate that appears in the classical
nonparametric regression literature, and that motivates choosing $\lambda$
proportional to the estimated noise level $\widehat\sigma^2$ (rather than
by a fixed rate in $n$) in the country-adaptive scheme of
Equation~\eqref{eq:sigma_est}.

\section{Random Elements in a Hilbert Space}
\label{sec:random_elements}

This section lays the probabilistic foundations of FDA: what it means for
a Hilbert-space-valued object to be a random element, its mean and
covariance operator, the Karhunen--Lo\`eve Theorem (the population-level
statement of FPCA), and the large-sample theory (strong law of large
numbers and central limit theorem) that Section~\ref{sec:mean_cov_estimation}
will use to study estimators of these quantities. We follow
\cite{hsingeubank2015}, Chapter~7, restricting to the material needed for
the FPCA pipeline (the random-element and mean-square-continuous-process
perspectives, and their large-sample theory); the material on
reproducing-kernel-valued processes and canonical-correlation-oriented
congruence relations in that source's Sections~7.5--7.6 is not needed here
and is omitted.

\subsection{Probability Measures on a Hilbert Space}
\label{subsec:prob_measures}

Throughout, $\mathbb{H}$ is a separable Hilbert space with its norm-induced
metric $d(f,g)=\|f-g\|$, and $\mathcal{B}(\mathbb{H})$ denotes its Borel
$\sigma$-field (the smallest $\sigma$-field containing all norm-open
subsets of $\mathbb{H}$) -- an unfortunate notational clash with the space
of bounded operators, resolved by context throughout. Let $\mathcal{M}$
denote the class of sets $\{x\in\mathbb{H}:\langle x,f\rangle\in B\}$ for
$f\in\mathbb{H}$ and $B\subset\mathbb{R}$ open.

\begin{theorem}
\label{thm:borel_sigma_field}
$\sigma(\mathcal{M})=\mathcal{B}(\mathbb{H})$.
\end{theorem}

\begin{proof}
Each set in $\mathcal{M}$ is open or closed (continuity of $x\mapsto
\langle x,f\rangle$), so $\mathcal{M}\subset\mathcal{B}(\mathbb{H})$ and
$\sigma(\mathcal{M})\subset\mathcal{B}(\mathbb{H})$. For the reverse
inclusion, it suffices to show every set $\{x:\|x-f\|>r\}$, $f\in
\mathbb{H}$, $r>0$, is in $\sigma(\mathcal{M})$. Fix a CONS $\{e_j\}$
(Theorem~\ref{thm:parseval}). Since $\|x\|^2=\sum_k\langle x,e_k\rangle^2$,
\begin{equation}
\{x:\|x\|>r\} = \bigcup_{j=1}^\infty\Big\{x:\sum_{k=1}^j\langle x,e_k
\rangle^2>r^2\Big\}.
\label{eq:borel_union}
\end{equation}
For $j=2$: $\{x:\langle x,e_1\rangle^2+\langle x,e_2\rangle^2>r^2\} =
\bigcup_{q\in\mathbb{Q}}\big(\{x:\langle x,e_1\rangle^2>q\}\cap\{x:
\langle x,e_2\rangle^2>r^2-q\}\big) \in\sigma(\mathcal{M})$, since both
sets on the right are in $\mathcal{M}$; induction using the same device
(splitting $\sum_{k\leq j}\langle x,e_k\rangle^2>r^2$ into the first
coordinate versus the rest, over rational thresholds) extends this to
every $j$, so \eqref{eq:borel_union} gives $\{x:\|x\|>r\}\in\sigma
(\mathcal{M})$. For general $f$, $\{x:\|x-f\|>r\} = \{x:\|x\|^2>r^2+
2\langle x,f\rangle-\|f\|^2\} = \bigcup_{q\in\mathbb{Q}}\big(\{x:\|x\|^2>q
\}\cap\{x:2\langle x,f\rangle<q-r^2+\|f\|^2\}\big)$, a countable union of
sets already shown (or immediately seen, for the second) to lie in
$\sigma(\mathcal{M})$, so $\{x:\|x-f\|>r\}\in\sigma(\mathcal{M})$ and
$\mathcal{B}(\mathbb{H})\subset\sigma(\mathcal{M})$.
\end{proof}

\begin{theorem}
\label{thm:random_element_char}
Let $\chi:(\Omega,\mathcal{F},\mathbb{P})\to(\mathbb{H},\mathcal{B}
(\mathbb{H}))$. Then: (1) $\chi$ is measurable if $\langle\chi,f\rangle$
is measurable for every $f\in\mathbb{H}$; (2) if $\chi$ is measurable, its
distribution is uniquely determined by the marginal distributions of
$\langle\chi,f\rangle$, $f\in\mathbb{H}$.
\end{theorem}

\begin{proof}
(1) If $\langle\chi,f\rangle$ is measurable for all $f$, then for any
$\{x:\langle x,f\rangle\in B\}\in\mathcal{M}$ ($B$ open), $\chi^{-1}
(\{x:\langle x,f\rangle\in B\}) = \{\omega:\langle\chi(\omega),f\rangle\in
B\}\in\mathcal{F}$; Theorem~\ref{thm:borel_sigma_field} then gives
measurability of $\chi$ (the converse direction is immediate from
continuity of $x\mapsto\langle x,f\rangle$).

(2) Let $\check{\mathcal{M}}$ be the class of finite intersections of sets
in $\mathcal{M}$; $\sigma(\check{\mathcal{M}}) = \sigma(\mathcal{M}) =
\mathcal{B}(\mathbb{H})$. If probability measures $\mu_1,\mu_2$ agree on
$\check{\mathcal{M}}$, then $\{A\in\mathcal{B}(\mathbb{H}):\mu_1(A)=
\mu_2(A)\}$ is a $\lambda$-system containing the $\pi$-system
$\check{\mathcal{M}}$, so by the $\pi$-$\lambda$ theorem it contains
$\sigma(\check{\mathcal{M}})=\mathcal{B}(\mathbb{H})$, i.e.\
$\check{\mathcal{M}}$ is a determining class. Since $\mathbb{P}\circ
\chi^{-1}(\cap_{i=1}^n\{x:\langle x,f_i\rangle\in B_i\}) = \mathbb{P}
(\langle\chi,f_i\rangle\in B_i,\,1\leq i\leq n)$, and the Cram\'er--Wold
device determines a joint distribution in $\mathbb{R}^n$ from its
one-dimensional projections, the distribution of $\chi$ is determined by
the marginal distributions of $\langle\chi,f\rangle$, $f\in\mathbb{H}$.
\end{proof}

A measurable $\chi:\Omega\to\mathbb{H}$ is called a \emph{random element}
of $\mathbb{H}$; this is the natural infinite-dimensional generalization of
a random vector, and is the object FPCA operates on -- in
Section~\ref{chap:methods}'s application, $\chi$ represents a country's
(centered) GDP growth trajectory viewed as an element of $\mathbb{L}^2
(\mathcal{T})$.

\subsection{Mean and Covariance Operator of a Random Element}
\label{subsec:mean_cov_operator}

\begin{definition}
\label{def:mean_element}
If $\mathbb{E}\|\chi\|<\infty$, the \emph{mean element} of a random
element $\chi$ of $\mathbb{H}$ is the Bochner integral $m=\mathbb{E}(\chi)
:=\int_\Omega\chi\,d\mathbb{P}$ (well defined by
Theorem~\ref{thm:bochner_exists} below); equivalently, since
$f\mapsto\mathbb{E}[\langle\chi,f\rangle]$ is then a bounded linear
functional, $m$ is its representer (Theorem~\ref{thm:adjoint}'s
Riesv-type argument, i.e.\ Theorem~3.2.1 of \cite{hsingeubank2015}), so
that in either case
\begin{equation}
\langle m,f\rangle = \mathbb{E}[\langle\chi,f\rangle], \qquad f\in
\mathbb{H}.
\label{eq:mean_element_char}
\end{equation}
\end{definition}

\begin{theorem}
\label{thm:bochner_exists}
If $\mathbb{E}\|\chi\|<\infty$ for a random element $\chi$ of the
separable Hilbert space $\mathbb{H}$, $\chi$ is Bochner integrable.
\end{theorem}

This is a direct consequence of $\mathbb{H}$ being separable: taking
$\mathbb{H}_n=\mathrm{span}\{e_1,\dots,e_n\}$ for a CONS $\{e_j\}$ and
$g_n$ the projection of $\chi$ onto $\mathbb{H}_n$ satisfies the
approximation condition for Bochner integrability by dominated
convergence; we omit the routine verification (\cite{hsingeubank2015},
Theorem 2.6.5).

\begin{theorem}
\label{thm:mean_variance_identity}
If $\mathbb{E}\|\chi\|^2<\infty$, $\mathbb{E}\|\chi-m\|^2 = \mathbb{E}
\|\chi\|^2-\|m\|^2$.
\end{theorem}

\begin{proof}
Expand $\mathbb{E}\|\chi-m\|^2 = \mathbb{E}\|\chi\|^2 - 2\mathbb{E}\langle
\chi,m\rangle+\|m\|^2$; by \eqref{eq:mean_element_char} with $f=m$,
$\mathbb{E}\langle\chi,m\rangle=\langle m,m\rangle=\|m\|^2$.
\end{proof}

\begin{definition}
\label{def:covariance_operator}
If $\mathbb{E}\|\chi\|^2<\infty$, the \emph{covariance operator} of $\chi$
is the Bochner integral (in $\mathcal{B}_{HS}(\mathbb{H})$, well defined
by Theorem~\ref{thm:bochner_exists} applied to $\mathcal{B}_{HS}
(\mathbb{H})$ since $\|(\chi-m)\otimes(\chi-m)\|_{HS}=\|\chi-m\|^2$ has
finite expectation by Theorem~\ref{thm:mean_variance_identity})
\begin{equation}
\Gamma := \mathbb{E}[(\chi-m)\otimes(\chi-m)] := \int_\Omega(\chi-m)
\otimes(\chi-m)\,d\mathbb{P}.
\label{eq:covariance_operator_def}
\end{equation}
\end{definition}

\begin{theorem}
\label{thm:covariance_identity}
$\mathbb{E}[(\chi-m)\otimes(\chi-m)] = \mathbb{E}(\chi\otimes\chi)-m\otimes
m$.
\end{theorem}

\begin{proof}
$\chi\otimes m$, $m\otimes\chi$, $m\otimes m$ are all HS with $m\otimes m$
constant and $\mathbb{E}\|\chi\otimes m\|_{HS}=\mathbb{E}\|m\otimes\chi
\|_{HS}=\|m\|\mathbb{E}\|\chi\|<\infty$. Since $(\mathbb{E}(m\otimes\chi))f
= m\langle\mathbb{E}\chi,f\rangle$ (Theorem~\ref{thm:bochner_exists}'s
linearity property, applied to the HS-valued Bochner integral) equals
$\langle m,f\rangle m = (m\otimes m)f$ by \eqref{eq:mean_element_char},
and likewise for $\mathbb{E}(\chi\otimes m)$, expanding $(\chi-m)\otimes
(\chi-m) = \chi\otimes\chi-\chi\otimes m-m\otimes\chi+m\otimes m$ and
taking expectations termwise gives the identity.
\end{proof}

We henceforth take $m=0$ without loss of generality (replacing $\chi$ by
$\chi-m$), so $\Gamma=\mathbb{E}(\chi\otimes\chi)$.

\begin{theorem}
\label{thm:covariance_properties}
Suppose $m=0$ and $\mathbb{E}\|\chi\|^2<\infty$. For $f,g\in\mathbb{H}$:
(1) $\langle\Gamma f,g\rangle = \mathbb{E}[\langle\chi,f\rangle\langle
\chi,g\rangle]$; (2) $\Gamma$ is nonnegative definite and trace class with
$\|\Gamma\|_{TR}=\mathbb{E}\|\chi\|^2$; (3) $\mathbb{P}(\chi\in
\mathrm{Im}(\Gamma))=1$.
\end{theorem}

\begin{proof}
(1) By \eqref{eq:covariance_operator_def} and linearity of the Bochner
integral (Theorem~\ref{thm:bochner_exists}'s linearity property applied to
the bounded linear functional $\mathcal{T}f:=\langle f,g\rangle$ on
$\mathcal{B}_{HS}(\mathbb{H})$), $\Gamma f = \int_\Omega\chi\langle\chi,f
\rangle\,d\mathbb{P}$, so $\langle\Gamma f,g\rangle = \mathbb{E}[\langle
\chi,f\rangle\langle\chi,g\rangle]$.

(2) Nonnegative definiteness is immediate from (1) with $f=g$. For the
trace: for any CONS $\{e_j\}$, $\|\Gamma\|_{TR} = \sum_j\langle\Gamma e_j,
e_j\rangle = \sum_j\mathbb{E}\langle\chi,e_j\rangle^2 = \mathbb{E}\|\chi
\|^2<\infty$ (Parseval, Theorem~\ref{thm:parseval}), so $\Gamma$ is trace
class.

(3) By Theorem~\ref{thm:adjoint_properties}(1) with self-adjointness of
$\Gamma$, $(\mathrm{Im}(\Gamma))^\perp = \mathrm{Ker}(\Gamma)$. For $f\in
(\mathrm{Im}(\Gamma))^\perp$, $\mathbb{E}[\langle\chi,f\rangle^2] =
\langle\Gamma f,f\rangle=0$ by part (1), so $\langle\chi,f\rangle=0$
almost surely. Thus, with probability one, $\chi\perp(\mathrm{Im}
(\Gamma))^\perp$, i.e.\ $\chi\in((\mathrm{Im}(\Gamma))^\perp)^\perp =
\mathrm{Im}(\Gamma)$ (closure, by Theorem~\ref{thm:projection}).
\end{proof}

By Theorem~\ref{thm:spectral_theorem} (applicable since, by Theorem
\ref{thm:covariance_properties}(2), $\Gamma$ is trace class, hence
Hilbert--Schmidt, hence compact, and self-adjoint), $\Gamma$ admits the
eigendecomposition
\begin{equation}
\Gamma = \sum_{j\geq 1}\lambda_je_j\otimes e_j,
\label{eq:covariance_eigendecomp}
\end{equation}
with $\{e_j\}$ a CONS for $\mathrm{Im}(\Gamma)$, $\lambda_1\geq\lambda_2
\geq\cdots\geq 0$ either finite in number or tending to $0$, and each
nonzero $\lambda_j$ of finite multiplicity -- extending the spectral
decomposition of a covariance matrix to the functional setting.

\begin{corollary}[Population Karhunen--Lo\`eve expansion]
\label{cor:population_kl}
With probability one, $\chi = \sum_{j\geq 1}\langle\chi,e_j\rangle e_j$,
where $\{\langle\chi,e_j\rangle\}_{j\geq 1}$ are uncorrelated, mean-zero
random variables with $\mathrm{Var}(\langle\chi,e_j\rangle)=\lambda_j$.
\end{corollary}

\begin{proof}
By Theorem~\ref{thm:covariance_properties}(3), $\chi\in\mathrm{Im}
(\Gamma)\subset\overline{\mathrm{span}}\{e_j\}$ almost surely; the
pathwise Fourier expansion of Theorem~\ref{thm:parseval} then gives $\chi=
\sum_j\langle\chi,e_j\rangle e_j$ a.s. Uncorrelatedness and the stated
variances follow from Theorem~\ref{thm:covariance_properties}(1) applied
to $f=e_i,g=e_j$: $\mathbb{E}[\langle\chi,e_i\rangle\langle\chi,e_j
\rangle]=\langle\Gamma e_i,e_j\rangle=\lambda_j\delta_{ij}$.
\end{proof}

Corollary~\ref{cor:population_kl} \emph{is} FPCA at the population level:
the $\langle\chi,e_j\rangle$ are the population FPC scores, and the
$e_j$ the population FPCs. The following result gives the precise sense in
which truncating this expansion at the leading $n$ eigenfunctions is
optimal -- the population counterpart of the variance-explained criterion
of Section~\ref{sec:metrics}.

\begin{theorem}
\label{thm:optimal_truncation}
For any CONS $\{f_j\}$ of $\mathbb{H}$,
\begin{equation}
\mathbb{E}\Big\|\chi-\sum_{j=1}^n\langle\chi,f_j\rangle f_j\Big\|^2 =
\mathbb{E}\|\chi\|^2 - \sum_{j=1}^n\langle\Gamma f_j,f_j\rangle,
\label{eq:optimal_truncation}
\end{equation}
which is minimized, over choices of orthonormal $f_1,\dots,f_n$, by
$f_j=e_j$, $1\leq j\leq n$.
\end{theorem}

\begin{proof}
Expanding the square and using $\mathbb{E}\|\sum_{j\leq n}\langle\chi,f_j
\rangle f_j\|^2 = \sum_{j\leq n}\mathbb{E}\langle\chi,f_j\rangle^2 =
\sum_{j\leq n}\langle\Gamma f_j,f_j\rangle$ (Theorem
\ref{thm:covariance_properties}(1)) and $\mathbb{E}\langle\chi,\sum_{j\leq
n}\langle\chi,f_j\rangle f_j\rangle = \sum_{j\leq n}\mathbb{E}\langle\chi,
f_j\rangle^2$ gives \eqref{eq:optimal_truncation} directly. Since
$\sum_{j=1}^n\langle\Gamma f_j,f_j\rangle\leq\sum_{j=1}^n\lambda_j$ for any
orthonormal $f_1,\dots,f_n$, with equality at $f_j=e_j$ (Theorem
\ref{thm:eigenvalue_max}/Theorem~\ref{thm:courant_fischer}, applied
successively as in the proof of Theorem~\ref{thm:eckart_young}), the
right side of \eqref{eq:optimal_truncation} is minimized at $f_j=e_j$.
\end{proof}

\subsection{Mean-Square Continuous Processes and the Karhunen--Lo\`eve
Theorem}
\label{subsec:kl_theorem}

We now turn to the complementary, stochastic-process perspective: $X=
\{X(t):t\in E\}$ on $(\Omega,\mathcal{F},\mathbb{P})$, for a compact
metric space $E$ (e.g.\ $E=\mathcal{T}$, the rescaled observation window),
with mean function $m(t):=\mathbb{E}[X(t)]$ and covariance function
$\gamma(s,t):=\mathrm{Cov}(X(s),X(t))$, provided these exist (a
\emph{second-order process}).

\begin{theorem}
\label{thm:covariance_kernel_nonneg}
The covariance function $\gamma$ of a second-order process is nonnegative
definite.
\end{theorem}

\begin{proof}
$\sum_{i,j}a_ia_j\gamma(t_i,t_j) = \mathrm{Var}\big(\sum_ja_jX(t_j)\big)
\geq 0$.
\end{proof}

By Theorem~\ref{thm:covariance_kernel_nonneg} and the Moore--Aronszajn
Theorem (Definition~\ref{def:rkhs}), $\gamma$ generates an RKHS
$\mathbb{H}(\gamma)$ for which it is the reproducing kernel.

\begin{definition}
\label{def:ms_continuous}
A second-order process $X$ is \emph{mean-square continuous} if
$\lim_{n\to\infty}\mathbb{E}[X(t_n)-X(t)]^2=0$ whenever $t_n\to t$ in $E$.
\end{definition}

\begin{theorem}
\label{thm:ms_continuity_char}
A second-order process $X$ is mean-square continuous if and only if $m$
and $\gamma$ are both continuous.
\end{theorem}

\begin{proof}
Since $\mathbb{E}[X(s)-X(t)]^2 = \gamma(s,s)+\gamma(t,t)-2\gamma(s,t)+
(m(s)-m(t))^2$, continuity of $m,\gamma$ gives mean-square continuity.
Conversely, mean-square continuity gives $|m(s)-m(t)|\leq\mathbb{E}[X(s)-
X(t)]^2{}^{1/2}\to 0$ (continuity of $m$); assuming WLOG $m\equiv 0$,
write $\gamma(s,t)-\gamma(s',t') = (\gamma(s,t)-\gamma(s',t)) +
(\gamma(s',t)-\gamma(s',t'))$, and bound each difference by
Cauchy--Schwarz: $|\gamma(s,t)-\gamma(s',t)|\leq\gamma(t,t)^{1/2}
\mathbb{E}[X(s)-X(s')]^2{}^{1/2}$ and similarly for the second term, both
of which $\to 0$ as $(s,t)\to(s',t')$ by mean-square continuity.
\end{proof}

For a mean-square continuous, mean-zero process, fix a finite measure
$\mu$ on $(E,\mathcal{B}(E))$; the associated \emph{covariance operator}
$\mathcal{K}$ is the integral operator of Definition
\ref{def:integral_operator} with kernel $\gamma$, and by Mercer's Theorem
(Theorem~\ref{thm:mercer}),
\begin{equation}
\gamma(s,t) = \sum_{j\geq 1}\lambda_je_j(s)e_j(t), \qquad s,t\in E,
\label{eq:covariance_mercer}
\end{equation}
absolutely and uniformly convergent, with $(\lambda_j,e_j)$ the
eigenpairs of $\mathcal{K}$.

We next construct the $\mathbb{L}^2$ stochastic integral of $X$ against
$f\in\mathbb{L}^2(E,\mathcal{B}(E),\mu)$, which is what makes precise the
otherwise formal expression $\int_EX(t)f(t)\,d\mu(t)$. For a partition
$\{E_i\}_{i=1}^{m(n)}$ of $E$ (with each $E_i\in\mathcal{B}(E)$ of diameter
$<1/n$, which exists by total boundedness of $E$) and arbitrary $t_i\in
E_i$, let $I_X(f;\{E_i,t_i\}) := \sum_iX(t_i)\int_{E_i}f\,d\mu$. For two
such partitions with parameters $n,n'$, direct computation of the squared
$\mathbb{L}^2(\Omega,\mathcal{F},\mathbb{P})$ difference between the two
Riemann-type sums, expanded via $\mathbb{E}[X(t_i)X(t_j')]=\gamma(t_i,
t_j')$, expresses it as a difference of double sums that (by uniform
continuity of $\gamma$, on the compact $E\times E$) converges, as $n,n'
\to\infty$, to $0$ once each double sum is recognized as a Riemann-type
approximation to $\iint_{E\times E}\gamma\,f\,f\,d\mu\,d\mu$ minus twice a
cross term converging to the same double integral; this shows
$\{I_X(f;\{E_i,t_i\})\}$ is Cauchy in $\mathbb{L}^2(\Omega,\mathcal{F},
\mathbb{P})$ and hence, by completeness (the Riesz--Fischer Theorem,
Definition~\ref{def:Lp}), converges to a limit $I_X(f)\in\mathbb{L}^2
(\Omega,\mathcal{F},\mathbb{P})$, independent of the approximating
sequence of partitions.

\begin{theorem}
\label{thm:stochastic_integral}
Let $X$ be mean-square continuous with mean zero. For $f,g\in\mathbb{L}^2
(E,\mathcal{B}(E),\mu)$: (1) $\mathbb{E}[I_X(f)]=0$; (2) $\mathbb{E}[I_X
(f)X(t)] = \int_E\gamma(u,t)f(u)\,d\mu(u)$; (3) $\mathbb{E}[I_X(f)I_X(g)]
= \iint_{E\times E}\gamma(u,v)f(u)g(v)\,d\mu(u)d\mu(v)$.
\end{theorem}

\begin{proof}
Each identity holds for the approximating Riemann-type sums by direct
computation (finite linearity and $\mathbb{E}[X(t_i)X(t_j)]=\gamma(t_i,
t_j)$), and passes to the limit $I_X(f)$ since $|\mathbb{E}[I_X(f)]-
\mathbb{E}[I_X(f;\{E_i,t_i\})]| \leq \mathbb{E}[(I_X(f)-I_X(f;\{E_i,t_i
\}))^2]^{1/2}\to 0$ by construction of $I_X(f)$ as the $\mathbb{L}^2$
limit, and similarly (via Cauchy--Schwarz) for the other two identities.
\end{proof}

\begin{corollary}
\label{cor:score_orthogonality}
$\mathbb{E}[I_X(e_i)I_X(e_j)]=\lambda_i\delta_{ij}$.
\end{corollary}

\begin{proof}
By Theorem~\ref{thm:stochastic_integral}(3) and \eqref{eq:covariance_mercer}
(Mercer expansion of $\gamma$), $\mathbb{E}[I_X(e_i)I_X(e_j)] = \iint
\gamma(u,v)e_i(u)e_j(v)\,d\mu(u)d\mu(v) = \langle\mathcal{K}e_j,e_i\rangle
= \lambda_j\delta_{ij}$.
\end{proof}

\begin{theorem}[Karhunen--Lo\`eve Theorem]
\label{thm:karhunen_loeve}
Let $X$ be mean-square continuous with mean zero. Then, with $X_n(t):=
\sum_{j=1}^nI_X(e_j)e_j(t)$,
\begin{equation}
\lim_{n\to\infty}\sup_{t\in E}\mathbb{E}[X(t)-X_n(t)]^2 = 0.
\label{eq:kl_convergence}
\end{equation}
\end{theorem}

\begin{proof}
By Theorem~\ref{thm:stochastic_integral}(2) (with $f=e_j$) and Corollary
\ref{cor:score_orthogonality},
\[
\mathbb{E}[X_n(t)-X(t)]^2 = \mathbb{E}[X_n(t)]^2-2\mathbb{E}[X_n(t)X(t)]+
\mathbb{E}[X(t)]^2 = \sum_{j=1}^n\lambda_je_j^2(t) - 2\sum_{j=1}^n
\lambda_je_j^2(t)+\gamma(t,t) = \gamma(t,t)-\sum_{j=1}^n\lambda_je_j^2(t),
\]
which tends to $0$ uniformly in $t$ by the uniform convergence in Mercer's
Theorem (\eqref{eq:covariance_mercer}, Theorem~\ref{thm:mercer}).
\end{proof}

The $I_X(e_j)$ are the \emph{scores} of $X$, $X_n$ the $n$-term
Karhunen--Lo\`eve expansion, and Theorem~\ref{thm:karhunen_loeve} states
that this expansion converges to $X$ uniformly in mean-square error --
the exact process-perspective analog of Corollary
\ref{cor:population_kl}, and the theoretical statement that the empirical
FPCA of Chapter~\ref{chap:results} approximates.

\begin{theorem}
\label{thm:process_as_element}
If $X$ is jointly measurable (as a function of $(t,\omega)$, w.r.t.\
$\mathcal{B}(E)\times\mathcal{F}$) and $X(\cdot,\omega)\in\mathbb{H}:=
\mathbb{L}^2(E,\mathcal{B}(E),\mu)$ for every $\omega$, then
$\omega\mapsto X(\cdot,\omega)$ is a random element of $\mathbb{H}$.
\end{theorem}

\begin{proof}
Joint measurability implies, by the standard proof of Fubini's Theorem,
that $\langle X(\cdot,\omega),f\rangle=\int_EX(t,\omega)f(t)\,d\mu(t)$ is
$\mathcal{F}$-measurable in $\omega$ for every $f\in\mathbb{H}$; apply
Theorem~\ref{thm:random_element_char}(1).
\end{proof}

Theorem~\ref{thm:process_as_element} is what licenses treating a GDP
growth trajectory -- observed, and mean-square continuous, as a
stochastic process $X(t)$ over the discretely sampled time domain -- as a
single random element $\chi=X(\cdot)$ of $\mathbb{L}^2(\mathcal{T})$ for
the purposes of Sections~\ref{subsec:mean_cov_operator} and
\ref{sec:fpca_theory}: the two perspectives of this section coincide, and
the covariance operator \eqref{eq:covariance_operator_def} of $\chi$
agrees with the integral operator \eqref{eq:integral_operator} built from
$\gamma$.

\subsection{Large Sample Theory}
\label{subsec:large_sample}

We conclude with the strong law of large numbers and central limit
theorem for i.i.d.\ random elements of a Hilbert space, which
Section~\ref{sec:mean_cov_estimation} applies to the sample mean and
sample covariance operator.

\begin{theorem}
\label{thm:pairwise_variance}
If $\chi_1,\dots,\chi_n$ are pairwise independent, mean-zero random
elements of $\mathbb{H}$ and $S_n:=\sum_{i=1}^n\chi_i$, then $\mathbb{E}
\|S_n\|^2 = \sum_{i=1}^n\mathbb{E}\|\chi_i\|^2$.
\end{theorem}

\begin{proof}
Fix a CONS $\{e_k\}$. Pairwise independence gives $\mathbb{E}[\langle
\chi_i,e_k\rangle\langle\chi_j,e_k\rangle] = \langle\mathbb{E}\chi_i,e_k
\rangle\langle\mathbb{E}\chi_j,e_k\rangle=0$ for $i\neq j$. By Parseval
(Theorem~\ref{thm:parseval}) and Tonelli's theorem (nonnegative summands),
$\mathbb{E}\|S_n\|^2 = \sum_k\mathbb{E}\langle S_n,e_k\rangle^2 = \sum_k
\sum_{i=1}^n\mathbb{E}\langle\chi_i,e_k\rangle^2 = \sum_{i=1}^n\sum_k
\mathbb{E}\langle\chi_i,e_k\rangle^2 = \sum_{i=1}^n\mathbb{E}\|\chi_i\|^2$.
\end{proof}

\begin{theorem}[Strong Law of Large Numbers]
\label{thm:slln}
If $\chi_1,\chi_2,\dots$ are pairwise independent and identically
distributed random elements of $\mathbb{H}$ with $\mathbb{E}\|\chi_1\|<
\infty$, then $n^{-1}S_n\to\mathbb{E}(\chi_1)$ almost surely.
\end{theorem}

\begin{proof}[Proof sketch]
Following the truncation argument of Etemadi (1983), set $\chi_i' :=
\chi_iI_{\{\|\chi_i\|\leq i\}}$, $S_n' := \sum_{i=1}^n\chi_i'$, and
$k_n:=\lfloor\alpha^n\rfloor$ for fixed $\alpha>1$. Applying Theorem
\ref{thm:mean_variance_identity}, Theorem~\ref{thm:pairwise_variance}, and
Markov's inequality along the geometric subsequence $\{k_n\}$, together
with the elementary bound $\sum_n k_n^{-2}I_{\{k_n\geq i\}} \leq 4(1-
\alpha^{-2})^{-1}i^{-2}$, gives $\sum_n\mathbb{P}(\|S_{k_n}'-\mathbb{E}
S_{k_n}'\|/k_n>\epsilon) \leq 4(1-\alpha^{-2})^{-1}\epsilon^{-2}\sum_i
\mathbb{E}\|\chi_i'\|^2i^{-2}$, and this last sum is finite (bounded by a
constant multiple of $\mathbb{E}\|\chi_1\|$, via a layer-cake / dyadic
decomposition of $\mathbb{E}\|\chi_1\|^2I_{\{j\leq\|\chi_1\|<j+1\}}$
summed against $\sum_{i>j}i^{-2}\asymp j^{-1}$); the Borel--Cantelli
Lemma then gives $(S_{k_n}'-\mathbb{E}S_{k_n}')/k_n\to 0$ a.s. Dominated
convergence gives $\mathbb{E}\chi_n'\to\mathbb{E}\chi_1$, hence (Cesàro)
$\mathbb{E}S_{k_n}'/k_n\to\mathbb{E}\chi_1$, so $S_{k_n}'/k_n\to
\mathbb{E}\chi_1$ a.s.; a further Borel--Cantelli argument shows
$\chi_i'=\chi_i$ eventually a.s., upgrading this to $S_{k_n}/k_n\to
\mathbb{E}\chi_1$ a.s. Finally, sandwiching a general $n$ between
consecutive $k_{m(n)-1}<n\leq k_{m(n)}$ and bounding $\|S_n/n-
S_{k_{m(n)}}/k_{m(n)}\|$ using the scalar strong law of large numbers
applied to $\|\chi_i\|$ (real-valued) shows $\limsup_n\|S_n/n-
S_{k_{m(n)}}/k_{m(n)}\|\leq 2(\alpha-1)\mathbb{E}\|\chi_1\|$ a.s.; letting
$\alpha\downarrow 1$ completes the proof. (Full details:
\cite{hsingeubank2015}, Theorem 7.7.2, adapting Etemadi, 1983.)
\end{proof}

For the central limit theorem, we need a suitable notion of weak
convergence on $\mathbb{H}$: $\chi_n\xrightarrow{d}\chi$ if $\mathbb{E}
[f(\chi_n)]\to\mathbb{E}[f(\chi)]$ for every bounded continuous
$f:\mathbb{H}\to\mathbb{R}$. The key technical tool is a characterization
of tightness adapted to Hilbert space.

\begin{definition}
\label{def:tightness}
A family $\{\mu_\alpha\}_{\alpha\in I}$ of probability measures on
$(\mathbb{H},\mathcal{B}(\mathbb{H}))$ is \emph{tight} if for every
$\epsilon>0$ there is a compact $W\subset\mathbb{H}$ with $\inf_\alpha
\mu_\alpha(W)\geq 1-\epsilon$. For $S\subset\mathbb{H}$, $\epsilon>0$, let
$S^\epsilon:=\{x\in\mathbb{H}:\inf_{z\in S}\|x-z\|\leq\epsilon\}$.
\end{definition}

\begin{theorem}
\label{thm:tightness_characterization}
Let $\{\mu_\alpha\}_{\alpha\in I}$ be probability measures on
$(\mathbb{H},\mathcal{B}(\mathbb{H}))$. Suppose that for every
$\epsilon,\delta>0$ there is a finite $\{y_1,\dots,y_k\}\subset\mathbb{H}$
such that, with $S:=\mathrm{span}\{y_1,\dots,y_k\}$: (1)
$\inf_\alpha\mu_\alpha(S^\epsilon)\geq 1-\delta$; (2) $\inf_\alpha
\mu_\alpha(\{x:|\langle x,y_j\rangle|\leq r,\,j=1,\dots,k\})\geq 1-\delta$
for some $r>0$. Then $\{\mu_\alpha\}$ is tight.
\end{theorem}

\begin{proof}
Fix $\epsilon>0$ and, applying the hypothesis with this $\epsilon=\delta$,
obtain orthonormal (WLOG) $y_1,\dots,y_k$, $S=\mathrm{span}\{y_j\}$, and
$r$ satisfying (1)--(2); set $A_\epsilon := S^\epsilon\cap\{x:|\langle x,
y_j\rangle|\leq r,\,j\leq k\}$, so $\inf_\alpha\mu_\alpha(A_\epsilon)\geq
1-2\epsilon$.

We show $A_\epsilon\subset W_\epsilon^\epsilon$ for a compact $W_\epsilon$.
For $x\in\mathbb{H}$ write $x=x_1+x_2$ with $x_1$ the projection onto $S$.
For $x\in A_\epsilon$: $\|x_2\|\leq\epsilon$ (definition of $S^\epsilon$),
so for any $z=z_1+z_2\in\mathbb{H}$, $|\langle x_2,z_2\rangle|\leq
\epsilon\|z\|$; and $|\langle x_1,z_1\rangle| = |\sum_{j\leq k}\langle
x,y_j\rangle\langle z,y_j\rangle|\leq kr\|z\|$ (Cauchy--Schwarz on the
$k$-vector of coordinates, using $|\langle x,y_j\rangle|\leq r$). Hence
$\|x\| = \sup_{\|z\|\leq1}|\langle x,z\rangle|\leq\epsilon+kr=:\nu$, so
$A_\epsilon\subset B_\nu\cap S^\epsilon$ for $B_\nu:=\{x:\|x\|\leq\nu\}$.

If $x\in B_\nu\cap S^\epsilon$, there is $y_j\in S$ (finite-dimensional,
so realizing the infimum, or within $\epsilon$ of it) with $\|x-y_j\|\leq
\epsilon$, so $\|y_j\|\leq\nu+\epsilon$, giving $\inf_{z\in S\cap
B_{\nu+\epsilon}}\|x-z\|\leq\epsilon$; thus $B_\nu\cap S^\epsilon\subset
(S\cap B_{\nu+\epsilon})^\epsilon =: W_\epsilon^\epsilon$ for the compact
(finite-dimensional, closed, bounded) set $W_\epsilon:=S\cap B_{\nu+
\epsilon}$. So $A_\epsilon\subset W_\epsilon^\epsilon$ and $\inf_\alpha
\mu_\alpha(W_\epsilon^\epsilon)\geq 1-2\epsilon$.

Repeating this construction at levels $\epsilon/2^j$, $j\geq 2$, and
setting $W:=\cap_{j\geq2}W_{\epsilon/2^j}^{\epsilon/2^j}$ gives $\inf_
\alpha\mu_\alpha(W)\geq 1-\epsilon$ (union bound over the tail
probabilities $2\cdot\epsilon/2^j$). It remains to show $W$ is compact:
$W$ is closed, and totally bounded because, writing $W=\cap_j(W_{\epsilon/
2^j}\cap W)\cup(\text{negligible remainder})$, each $W_{\epsilon/2^j}\cap
W$ is compact (closed subset of the compact $W_{\epsilon/2^j}$), so for
any $\nu'>0$, choosing $j$ with $\epsilon/2^j<\nu'/2$ and a finite
$\nu'/2$-net $F$ of $W_{\epsilon/2^j}\cap W$ gives $W\subset F^{\nu'}$ (any
$x\in W$ is within $\epsilon/2^j<\nu'/2$ of $W_{\epsilon/2^j}\cap W$,
which is within $\nu'/2$ of $F$). By the Heine--Borel Theorem (a complete,
totally bounded metric space is compact; cf.\ Theorem 2.1.17 of
\cite{hsingeubank2015}), $W$ is compact.
\end{proof}

\begin{theorem}
\label{thm:weak_convergence_criterion}
Let $\chi,\chi_n$ be random elements of $\mathbb{H}$ with $\langle\chi_n,f
\rangle\xrightarrow{d}\langle\chi,f\rangle$ in $\mathbb{R}$ for every
$f\in\mathbb{H}$, and suppose that for every $\epsilon,\delta>0$ there is
a finite-dimensional $S$ with $\inf_{n\geq1}\mathbb{P}(\chi_n\in
S^\epsilon)\geq1-\delta$. Then $\chi_n\xrightarrow{d}\chi$.
\end{theorem}

\begin{proof}
The hypothesis verifies Theorem~\ref{thm:tightness_characterization}'s
conditions for $\{\mathbb{P}\circ\chi_n^{-1}\}$ (condition (2) follows
from convergence in distribution of each $\langle\chi_n,y_j\rangle$,
which is in particular bounded in probability), so this family is tight,
hence (Prohorov's Theorem) relatively compact in the weak topology. Let
$\mathbb{P}\circ\chi_{n'}^{-1}\xrightarrow{w}\nu_1$ and $\mathbb{P}\circ
\chi_{n''}^{-1}\xrightarrow{w}\nu_2$ be two weakly convergent
subsequences, with limiting random elements $\widetilde\chi,\check{\chi}$
(i.e.\ $\chi_{n'}\xrightarrow{d}\widetilde\chi$, $\chi_{n''}
\xrightarrow{d}\check{\chi}$). By the continuous mapping theorem
applied to $x\mapsto\langle x,f\rangle$, $\langle\chi_{n'},f\rangle
\xrightarrow{d}\langle\widetilde\chi,f\rangle$ and $\langle\chi_{n''},f
\rangle\xrightarrow{d}\langle\check{\chi},f\rangle$; but both
subsequences of $\langle\chi_n,f\rangle$ converge in distribution to
$\langle\chi,f\rangle$ by hypothesis, so $\langle\widetilde\chi,f\rangle
\stackrel{d}{=}\langle\check{\chi},f\rangle\stackrel{d}{=}\langle\chi,f
\rangle$ for every $f$; by Theorem~\ref{thm:random_element_char}(2),
$\widetilde\chi\stackrel{d}{=}\check{\chi}\stackrel{d}{=}\chi$. As every
weakly convergent subsequence of $\{\mathbb{P}\circ\chi_n^{-1}\}$ has the
same limit $\mathbb{P}\circ\chi^{-1}$, the whole sequence converges
weakly to it.
\end{proof}

\begin{theorem}[Central Limit Theorem]
\label{thm:hilbert_clt}
Let $\chi_1,\chi_2,\dots$ be i.i.d.\ random elements of $\mathbb{H}$ with
mean $0$ and $\mathbb{E}\|\chi_1\|^2<\infty$. Then
\begin{equation}
\xi_n := n^{-1/2}S_n \xrightarrow{d} \xi,
\label{eq:hilbert_clt}
\end{equation}
where $\xi$ is a Gaussian random element of $\mathbb{H}$ with covariance
operator $\mathbb{E}(\chi_1\otimes\chi_1)$.
\end{theorem}

\begin{proof}
We verify the hypotheses of Theorem~\ref{thm:weak_convergence_criterion}.
For any $f\in\mathbb{H}$, $\langle\xi_n,f\rangle = n^{-1/2}\sum_i\langle
\chi_i,f\rangle$ is a normalized sum of i.i.d.\ real random variables with
mean $0$ and variance $\langle\mathbb{E}(\chi_1\otimes\chi_1)f,f\rangle$,
so by the classical (scalar) CLT, $\langle\xi_n,f\rangle\xrightarrow{d}
N(0,\langle\mathbb{E}(\chi_1\otimes\chi_1)f,f\rangle)$, which is the
distribution of $\langle\xi,f\rangle$.

For the tightness condition, fix a CONS $\{e_j\}$, let $S_J:=
\mathrm{span}\{e_1,\dots,e_J\}$, and let $\xi_{nJ}'$ be the projection of
$\xi_n$ onto $S_J^\perp$ (so $\mathbb{P}(\|\xi_{nJ}'\|\leq\epsilon) =
\mathbb{P}(\xi_n\in S_J^\epsilon)$). Writing $\chi_{iJ}'$ for the
projection of $\chi_i$ onto $S_J^\perp$, Chebyshev's inequality and
Theorem~\ref{thm:pairwise_variance} (applied to the i.i.d., hence pairwise
independent, mean-zero elements $\chi_{iJ}'$) give
\[
\mathbb{P}(\|\xi_{nJ}'\|>\epsilon) \leq \epsilon^{-2}\mathbb{E}\|\xi_{nJ}
'\|^2 = \epsilon^{-2}n^{-1}\sum_{i=1}^n\mathbb{E}\|\chi_{iJ}'\|^2 =
\epsilon^{-2}\mathbb{E}\|\chi_{1J}'\|^2,
\]
uniformly in $n$; since $\mathbb{E}\|\chi_1\|^2<\infty$, $\mathbb{E}
\|\chi_{1J}'\|^2\to 0$ as $J\to\infty$ (dominated convergence, as
$\chi_{1J}'\to 0$ pointwise on $\Omega$), so this can be made $<\delta$
for $J$ large, verifying the hypothesis of
Theorem~\ref{thm:weak_convergence_criterion} with $S=S_J$.
\end{proof}

Theorem~\ref{thm:hilbert_clt} is the tool that
Section~\ref{sec:mean_cov_estimation} applies to $\chi_i=X_i-m$ (giving
the asymptotic normality of the sample mean $\overline{X}_n$) and, via the
linearization $\Delta=\widetilde\Gamma-\Gamma$ of
Section~\ref{sec:perturbation}, to the sample covariance operator and, in
turn, to the estimated FPCA eigenvalues and eigenfunctions of
Section~\ref{sec:fpca_theory}.

\section{Mean and Covariance Operator Estimation}
\label{sec:mean_cov_estimation}

This section studies the large-sample behavior of estimators of the mean
element $m$ and covariance operator $\Gamma$, under three progressively
more realistic observation schemes: (i) completely observed functional
data (the sample mean and sample covariance operator);
(ii) discretely, noisily observed sample paths, estimated by local-linear
smoothing; (iii) the same discretely observed data, estimated instead by
penalized least squares (smoothing splines). We follow
\cite{hsingeubank2015}, Chapter~8. Case (i) is proved in full, since it
follows directly and briefly from the large-sample theory of
Section~\ref{subsec:large_sample}; cases (ii) and (iii) involve very long
uniform-consistency arguments (empirical-process bounds via Bernstein's
inequality over a discretized grid, for (ii); a Fr\'echet-derivative
linearization of a penalized criterion in a Sobolev basis, for (iii)) that
are not directly used by this thesis's estimation strategy --- which
individually smooths each country's trajectory with an adaptive B-spline
smoothing spline (Section~\ref{subsec:bspline}, justified in full by
Section~\ref{sec:smoothing}) and then applies the case-(i) sample
covariance operator to the resulting smoothed curves. For (ii) and (iii) we
therefore state the precise models and rate theorems, and give a proof
sketch conveying the strategy and the source of each rate, rather than the
complete technical derivation; full detail is in
\cite{hsingeubank2015}, Sections~8.2--8.3.

\subsection{Sample Mean and Covariance Operator}
\label{subsec:sample_mean_cov}

Let $\chi$ be a random element of a separable Hilbert space $\mathbb{H}$
with $\mathbb{E}\|\chi\|^2<\infty$, mean $m$, and covariance operator
$\Gamma$ (Section~\ref{subsec:mean_cov_operator}), and let $\chi_1,\dots,
\chi_n$ be an i.i.d.\ sample from $\chi$. Define the \emph{sample mean}
$m_n:=n^{-1}\sum_{i=1}^n\chi_i$ and \emph{sample covariance operator}
\begin{equation}
\Gamma_n := \frac{1}{n-1}\sum_{i=1}^n(\chi_i-m_n)\otimes(\chi_i-m_n).
\label{eq:sample_covariance}
\end{equation}

\begin{theorem}
\label{thm:sample_mean_asymptotics}
If $\mathbb{E}\|\chi_1\|<\infty$, $m_n\to m$ almost surely. If $\mathbb{E}
\|\chi_1\|^2<\infty$, $\sqrt{n}(m_n-m)\xrightarrow{d}\xi$ in $\mathbb{H}$,
where $\xi$ is Gaussian with mean $0$ and covariance operator $\Gamma$.
\end{theorem}

\begin{proof}
Immediate from Theorem~\ref{thm:slln} (with $\chi_i$ in place of $\chi_i-m$,
noting $n^{-1}S_n-m = n^{-1}\sum(\chi_i-m)$) and Theorem
\ref{thm:hilbert_clt} (applied to $\chi_i-m$, whose covariance operator is
$\Gamma$ by definition).
\end{proof}

\begin{theorem}
\label{thm:sample_covariance_asymptotics}
If $\mathbb{E}\|\chi_1\|^2<\infty$, $\Gamma_n\to\Gamma$ almost surely. If
$\mathbb{E}\|\chi_1\|^4<\infty$, $\sqrt{n}(\Gamma_n-\Gamma)\xrightarrow{d}
\mathfrak{Z}$ in $\mathcal{B}_{HS}(\mathbb{H})$, where $\mathfrak{Z}$ is
Gaussian with mean $0$ and covariance operator $\mathbb{E}\big[\{(\chi_1-m)
\otimes(\chi_1-m)-\Gamma\}\otimes_{HS}\{(\chi_1-m)\otimes(\chi_1-m)-\Gamma\}
\big]$ (an element of $\mathcal{B}(\mathcal{B}_{HS}(\mathbb{H}))$, via the
HS-space tensor product of Definition~\ref{def:hs_operator}, applied one
level up to the Hilbert space $\mathcal{B}_{HS}(\mathbb{H})$ itself).
\end{theorem}

\begin{proof}
Write $\Gamma_n = \frac{1}{n-1}\sum_{i=1}^n(\chi_i-m)\otimes(\chi_i-m) -
\frac{n}{n-1}(m_n-m)\otimes(m_n-m)$ (direct algebraic expansion). The
second term is $o(1)$ a.s.\ by Theorem~\ref{thm:sample_mean_asymptotics}
(so $m_n-m\to0$ a.s., hence $(m_n-m)\otimes(m_n-m)\to0$ in HS norm), and
the first term $\to\Gamma$ a.s.\ by Theorem~\ref{thm:slln} applied to the
i.i.d.\ $\mathcal{B}_{HS}(\mathbb{H})$-valued random elements $(\chi_i-m)
\otimes(\chi_i-m)$, giving the strong consistency of $\Gamma_n$.

For the CLT, Theorem~\ref{thm:hilbert_clt} applies to $(\chi_i-m)\otimes
(\chi_i-m)-\Gamma$ (i.i.d., mean $0$, valued in the Hilbert space
$\mathcal{B}_{HS}(\mathbb{H})$) provided $\mathbb{E}\|(\chi_1-m)\otimes
(\chi_1-m)-\Gamma\|_{HS}^2<\infty$; by Theorem~\ref{thm:mean_variance_identity}
applied in $\mathcal{B}_{HS}(\mathbb{H})$, this is bounded by $\mathbb{E}
\|(\chi_1-m)\otimes(\chi_1-m)\|_{HS}^2 = \mathbb{E}\|\chi_1-m\|^4<\infty$.
The $n/(n-1)$-weighted second term of $\Gamma_n$ is $O_p(n^{-1})$ (since
$\sqrt{n}(m_n-m)$ is asymptotically Gaussian by
Theorem~\ref{thm:sample_mean_asymptotics}, so $n(m_n-m)\otimes(m_n-m) =
O_p(1)$), hence contributes nothing to the limit of $\sqrt{n}(\Gamma_n-
\Gamma)$, which therefore coincides with that of $\sqrt{n}\big(n^{-1}
\sum_i[(\chi_i-m)\otimes(\chi_i-m)-\Gamma]\big)$.
\end{proof}

To interpret Theorem~\ref{thm:sample_covariance_asymptotics}: for any
$\mathcal{G}\in\mathcal{B}_{HS}(\mathbb{H})$, $\langle\sqrt{n}(\Gamma_n-
\Gamma),\mathcal{G}\rangle_{HS}\xrightarrow{d}\langle\mathfrak{Z},
\mathcal{G}\rangle_{HS}\sim N(0,\mathrm{Var}(\langle(\chi_1-m)\otimes(\chi_1
-m),\mathcal{G}\rangle_{HS}))$, and choosing a CONS $\{e_j\}$ with $e_1=
(\chi_1-m)/\|\chi_1-m\|$ shows $\langle(\chi_1-m)\otimes(\chi_1-m),
\mathcal{G}\rangle_{HS} = \langle\chi_1-m,\mathcal{G}(\chi_1-m)\rangle$, so
$\mathrm{Var}(\langle\sqrt{n}(\Gamma_n-\Gamma),\mathcal{G}\rangle_{HS})\to
\mathrm{Var}(\langle\chi_1-m,\mathcal{G}(\chi_1-m)\rangle)$ -- an explicit,
directly interpretable variance for any linear functional of the estimated
covariance operator. This is the estimator, and the exact asymptotic theory,
that Chapter~\ref{chap:methods} applies (with $\mathbb{H}=\mathbb{L}^2
(\mathcal{T})$ and $\chi_i$ the B-spline-smoothed country trajectory) to
obtain the empirical covariance operator whose spectral decomposition is
the FPCA of Chapter~\ref{chap:results}.

\subsection{Local Linear Estimation}
\label{subsec:local_linear}

Suppose instead that $X_1,\dots,X_n$ are i.i.d.\ copies of a mean-square
continuous process $X$ on $[0,1]$ with mean $m$ and covariance $\gamma$,
observed only at discrete, noisy points:
\begin{equation}
Y_{ij} = X_i(T_{ij})+\varepsilon_{ij}, \qquad j=1,\dots,r,\ i=1,\dots,n,
\label{eq:local_linear_model}
\end{equation}
with $T_{ij}$ i.i.d.\ sampling points (density $f$, $f'$ continuous) and
$\varepsilon_{ij}$ i.i.d.\ mean-zero noise, all mutually independent
collections. A \emph{local linear} estimator of $m(t)$ minimizes, in
$(a_0,a_1)$, the kernel-weighted criterion $\sum_{i,j}W_h(T_{ij}-t)(Y_{ij}
-a_0-a_1(T_{ij}-t))^2$ for a symmetric density $W$ on $[-1,1]$ (of bounded
variation, with $\int u^jW(u)\,du$ equal to $1,0,C\neq0$ for $j=0,1,2$) and
bandwidth $h$, giving the closed-form estimator
\begin{equation}
m_h(t) = \frac{S_0(t)M_2(t)-S_1(t)M_1(t)}{M_0(t)M_2(t)-M_1(t)^2}, \qquad
M_p(t):=\frac{1}{nr}\sum_{i,j}W_{h^p}(T_{ij}-t),\ \ S_p(t):=\frac{1}{nr}
\sum_{i,j}W_{h^p}(T_{ij}-t)Y_{ij},
\label{eq:local_linear_estimator}
\end{equation}
where $W_{h^p}(u):=(u/h)^pW_h(u)$.

\begin{theorem}
\label{thm:local_linear_mean_rate}
Suppose $m''$ is uniformly bounded, $\mathbb{E}|\varepsilon_{11}|^q<\infty$
and $\mathbb{E}\sup_t|X(t)|^q<\infty$ for some $q\in(2,\infty)$, and $h\to0$
with $(h^2+h/r)^{-1}(\log n/n)^{1-2/q}\to0$. Then, with
$\delta_{n1}(h):=(\{1+(hr)^{-1}\}\log n/n)^{1/2}$,
\begin{equation}
\sup_{t\in[0,1]}|m_h(t)-m(t)| = O(h^2+\delta_{n1}(h)) \quad\text{a.s.}
\label{eq:local_linear_mean_rate}
\end{equation}
\end{theorem}

\begin{proof}[Proof sketch]
Write $m_h(t)-m(t) = \{\widetilde S_0(t)M_2(t)-\widetilde S_1(t)M_1(t)\}/
\{M_0(t)M_2(t)-M_1(t)^2\}$ for $\widetilde S_p(t):=S_p(t)-m(t)M_p(t)-
hm'(t)M_{p+1}(t)$. The denominator and the $M_p$ converge a.s.\ uniformly
to nonzero limits by a general uniform strong law for kernel-weighted
sums of bounded triangular arrays (stated below as
Lemma~\ref{lem:uniform_kernel_slln}), so the rate is governed by
$\widetilde{S}_p$. Splitting $\widetilde{S}_p(t) = \{$a deterministic
Taylor-remainder term, $O(h^2)$ a.s.\ by boundedness of $m''\}$ $+\, U_p(t)$
for $U_p(t):=(nr)^{-1}\sum_{i,j}W_{h^p}(T_{ij}-t)[\varepsilon_{ij}+X_i
(T_{ij})-m(T_{ij})]$, a further application of
Lemma~\ref{lem:uniform_kernel_slln} (to the two zero-mean triangular arrays
comprising $U_p$) gives $\sup_t|U_p(t)| = O(\delta_{n1}(h))$ a.s.
\end{proof}

\begin{lemma}
\label{lem:uniform_kernel_slln}
Let $\{Z_{ij}\}$ be a triangular array with $\sup_{i,j}\mathbb{E}|Z_{ij}|^q
<\infty$, $(nr)^{-1}\sum_{i,j}|Z_{ij}|^q=O(1)$ a.s., and $\sup_{i,j,k}
\mathbb{E}[|Z_{ij}Z_{ik}|\mid T_{ij},T_{ik}]<C$ a.s., for some $q\in(2,
\infty)$; let $\overline{Z}_p(t):=(nr)^{-1}\sum_{i,j}W_{h^p}(T_{ij}-t)
Z_{ij}$ and $\beta_n:=h^2+h/r$ with $\beta_n^{-1}(\log n/n)^{1-2/q}\to0$.
Then $\sup_t|\overline{Z}_p(t)-\mathbb{E}\overline{Z}_p(t)| = O(\{\beta_n
\log n/n\}^{1/2})$ a.s.
\end{lemma}

This is \cite{hsingeubank2015}, Lemma 8.2.2; its proof is a uniform law of
large numbers for kernel-weighted triangular arrays, obtained by
discretizing $[0,1]$ into a grid of $O(1/h)$ points, truncating $Z_{ij}$ at
a slowly growing threshold $Q_n$, and applying Bernstein's inequality plus
Borel--Cantelli at each grid point (a two-page empirical-process argument
we do not reproduce, as it is generic nonparametric-regression technology
unrelated to the FDA-specific content of this thesis).

\begin{corollary}
\label{cor:local_linear_mean_rates}
Under the conditions of Theorem~\ref{thm:local_linear_mean_rate}: if $r$ is
bounded, $\sup_t|m_h(t)-m(t)| = O(h^2+[\log n/(nh)]^{1/2})$ a.s.\ (the
classical univariate nonparametric rate); if $r=r_n$ with $r_n^{-1}
\lesssim h\lesssim(\log n/n)^{1/4}$, $\sup_t|m_h(t)-m(t)| = O([\log n/n]
^{1/2})$ a.s.\ (a parametric rate, attained once $r\gtrsim n^{1/4}$).
\end{corollary}

An analogous, more involved construction (weighting products $Y_{ij}Y_{ik}$,
$j\neq k$, by $W_h(T_{ij}-s)W_h(T_{ik}-t)$ to estimate $R(s,t):=\mathbb{E}
[X(s)X(t)]$, then setting $\gamma_h(s,t):=R_{h_R}(s,t)-m_{h_m}(s)m_{h_m}
(t)$) yields a local-linear covariance estimator $\gamma_h$ with rate
$\sup_{s,t}|\gamma_h(s,t)-\gamma(s,t)| = O(h_m^2+h_R^2+\delta_{n1}(h_m)+
\delta_{n2}(h_R))$ a.s.\ (with $\delta_{n2}(h):=(\{1+(hr)^{-1}+(hr)^{-2}\}
\log n/n)^{1/2}$), under fourth-moment analogs of the hypotheses above
(\cite{hsingeubank2015}, Theorem 8.2.4); the sparse and dense rates parallel
Corollary~\ref{cor:local_linear_mean_rates}, with the sparse rate
$O(n^{-1/3})$ optimized at bandwidth $\widehat{h}_R\asymp n^{-1/6}$.

The rate that will actually matter for FPCA is not the pointwise kernel
rate for $\gamma_h(s,t)-\gamma(s,t)$, but the rate at which the induced
\emph{integral operator} converges -- and these turn out to differ.

\begin{theorem}
\label{thm:local_linear_operator_rate}
Under the conditions of the local-linear covariance rate above, for any
bounded measurable $e$ on $[0,1]$, with $(\Delta_he)(t):=\int_0^1\{\gamma_h
(s,t)-\gamma(s,t)\}e(s)\,ds$,
\begin{equation}
\sup_{t\in[0,1]}|(\Delta_he)(t)| = O(h_m^2+h_R^2+\delta_{n1}(h_m)+
\delta_{n1}(h_R)) \quad\text{a.s.}
\label{eq:local_linear_operator_rate}
\end{equation}
\end{theorem}

Note $\delta_{n1}(h_R)$, not the (slower) $\delta_{n2}(h_R)$ of the
pointwise rate: integrating against a fixed $e$ before taking the supremum
in $t$ averages out one order of the kernel's stochastic fluctuation. This
gain is exactly the mechanism Section~\ref{sec:fpca_theory} needs to
transfer local-linear covariance estimation into a rate on the estimated
covariance \emph{operator} in $\mathcal{B}(\mathbb{L}^2[0,1])$, hence (via
Section~\ref{sec:perturbation}) into a rate for the estimated eigenvalues
and eigenfunctions; we omit the proof (\cite{hsingeubank2015}, Theorem
8.2.7), which decomposes $\Delta_he$ analogously to
Theorem~\ref{thm:local_linear_mean_rate} and reapplies
Lemma~\ref{lem:uniform_kernel_slln}-type bounds to the resulting
integrated stochastic terms.

\subsection{Penalized Least-Squares Estimation}
\label{subsec:penalized_ls_estimation}

We return to model~\eqref{eq:local_linear_model} (with $T_{ij}$ now taken
uniform on $[0,1]$, for simplicity), but assume instead that $X$ is a
random element of the Sobolev space $\mathbb{W}^q[0,1]$ of Definition
\ref{def:sobolev}, $q>1$, with CONS $\{e_j\}$ (the eigenfunctions of
$(-1)^qD^{2q}$ from Theorem~2.8.3 of \cite{hsingeubank2015}, satisfying
$\langle e_j,e_k\rangle_2=\delta_{jk}$, $\langle e_j^{(q)},e_k^{(q)}
\rangle_2=\gamma_j\delta_{jk}$ with $\gamma_1=\cdots=\gamma_q=0$ and
$C_1j^{2q}\leq\gamma_{j+q}\leq C_2j^{2q}$), so that $\|v\|_{\mathbb{W}^q
[0,1]}^2=\sum_j(1+\gamma_j)v_j^2$ for $v=\sum_jv_je_j$. This is a more
restrictive but sharper setting than Section~\ref{subsec:local_linear}'s,
and the resulting smoothing-spline estimators are rate-optimal.

\begin{theorem}[Minimax lower bound; Cai and Yuan, 2011]
\label{thm:cai_yuan_lower_bound}
Let $\mathbb{P}(q,C_1)$ be the probability laws of $\mathbb{W}^q[0,1]$-valued
processes $X$ with $\mathbb{E}\|X^{(q)}\|_2^2\leq C_1$. There is $C_2>0$,
depending only on $C_1,\sigma^2=\mathbb{E}\varepsilon_{11}^2$, such that
every estimator $\widehat{m}$ of $m$ satisfies
\begin{equation}
\limsup_{n\to\infty}\ \sup_{\mathbb{P}_X\in\mathbb{P}(q,C_1)}\mathbb{P}
\big(\|\widehat{m}-m\|_2^2 > C_2\{(nr)^{-2q/(2q+1)}+n^{-1}\}\big) > 0.
\label{eq:minimax_lower_bound}
\end{equation}
\end{theorem}

The smoothing-spline estimator $m_\eta:=\mathrm{argmin}_{v\in\mathbb{W}^q
[0,1]}f_{rn,\eta}(v)$ for $f_{rn,\eta}(v):=(nr)^{-1}\sum_{i,j}(Y_{ij}-
v(T_{ij}))^2+\eta\|v^{(q)}\|_2^2$ (the case of Theorem~\ref{thm:representer}
with $\mathcal{T}$ the discrete evaluation map on $[0,1]$ into $\mathbb{R}^{
nr}$) attains this lower bound.

\begin{theorem}
\label{thm:penalized_mean_rate}
If $\eta\asymp(rn)^{-2q/(2q+1)}$, $\|m_\eta-m\|_2^2 = O_p((nr)^{-2q/(2q+1)}
+n^{-1})$.
\end{theorem}

\begin{proof}[Proof sketch]
Following Cai and Yuan (2011), let $f_{\infty,\eta}(v):=\mathbb{E}
f_{rn,\eta}(v) = \mathbb{E}[(Y_{11}-m(T_{11}))^2]+\|m-v\|_2^2+\eta\|v^{(q)}
\|_2^2$, $\overline{m}_\eta:=\mathrm{argmin}_vf_{\infty,\eta}(v)$, and the
one-step Newton-type intermediate $\widetilde{m}_\eta := m_\eta - (f_{
\infty,\eta}'')^{-1}f_{rn,\eta}'(m_\eta)$ (Fr\'echet derivatives, in the
sense of Section~3.6 of \cite{hsingeubank2015}), so that $m_\eta-m =
(m_\eta-\widetilde{m}_\eta)+(\widetilde{m}_\eta-\overline{m}_\eta)+
(\overline{m}_\eta-m)$.

The last term is the easiest: since $m_\eta$ minimizes $f_{\infty,\eta}$
with $m$ itself a competitor, $\|\overline{m}_\eta-m\|_2^2 \leq
f_{\infty,\eta}(\overline{m}_\eta)-\mathbb{E}[(Y_{11}-m(T_{11}))^2] \leq
f_{\infty,\eta}(m)-\mathbb{E}[(Y_{11}-m(T_{11}))^2] = \eta\|m^{(q)}\|_2^2 =
O((rn)^{-2q/(2q+1)})$ under the assumed rate on $\eta$ -- an
instance of the bias bound Theorem~\ref{thm:bias_variance} specialized to
this criterion.

The remaining two terms, $\|m_\eta-\widetilde{m}_\eta\|_2^2$ and $\|
\widetilde{m}_\eta-\overline{m}_\eta\|_2^2$, are shown to be $O_p((nr)^{
-2q/(2q+1)}+n^{-1})$ by expressing both in the interpolation-scale norms
$\|v\|_\alpha^2:=\sum_j(1+\gamma_j)^\alpha v_j^2$, $\alpha\in[0,1]$ (which
satisfy $\|v\|_0=\|v\|_2\leq\|v\|_\alpha\leq\|v\|_1=\|v\|_{\mathbb{W}^q
[0,1]}$), using the explicit representation $(\widetilde{f}_{\infty,\eta}
'')^{-1}v = \frac{1}{2}\sum_j\frac{1+\gamma_j}{1+\eta\gamma_j}v_je_j$ of
the (Riesz-representer form of the) second Fr\'echet derivative's inverse.
This reduces $\|m_\eta-\widetilde{m}_\eta\|_\alpha^2$ and $\|\widetilde{m}
_\eta-\overline{m}_\eta\|_\alpha^2$ to explicit weighted sums,
$\frac{1}{4}\sum_k\frac{(1+\gamma_k)^\alpha}{(1+\eta\gamma_k)^2}(f_{rn,
\eta}'(\overline{m}_\eta)e_k)^2$ and an analogous expression, whose
expectations are bounded via a variance computation on $f_{rn,\eta}'(
\overline{m}_\eta)e_k$ (using $\mathbb{E}[f_{rn,\eta}'(\overline{m}_\eta)
e_k]=0$, since $f_{\infty,\eta}'(\overline{m}_\eta)=0$) together with the
polynomial growth rate $\gamma_j\asymp j^{2q}$, giving $O_p(n^{-1}+
(nr)^{-1}\eta^{-\alpha-1/(2q)})$-type bounds that, at $\alpha=0$ and the
stated rate for $\eta$, both reduce to $O_p((nr)^{-2q/(2q+1)}+n^{-1})$. We
omit the (lengthy, but routine) algebraic details; full derivation in
\cite{hsingeubank2015}, Theorem 8.3.2.
\end{proof}

An entirely parallel construction --- penalizing a bivariate smoothing
criterion $f_{rn,\eta}(K):=\{r(r-1)n\}^{-1}\sum_{i}\sum_{j\neq k}(Y_{ij}
Y_{ik}-K(T_{ij},T_{ik}))^2+\eta\|K\|_{\mathbb{H}}^2$ over the tensor-product
Sobolev space $\mathbb{H}=\mathbb{W}^q[0,1]\otimes\mathbb{W}^q[0,1]$
(Definition~\ref{def:rkhs}'s product-kernel construction) --- gives a
smoothing-spline covariance estimator $K_\eta$.

\begin{theorem}
\label{thm:penalized_cov_rate}
At the rate $\eta\asymp(rn)^{-2q/(2q+1)}$, with $\delta_n^2:=(nr)^{-2q/(2q+1)}
+n^{-1}$,
\begin{equation}
\|K_\eta-\gamma\|_{\mathbb{L}^2([0,1]^2)}^2 = O_p(\delta_n^2).
\label{eq:penalized_cov_rate}
\end{equation}
Consequently, if $\Delta_\eta\in\mathcal{B}(\mathbb{L}^2[0,1])$ is the
integral operator with kernel $K_\eta-\gamma$, $\|\Delta_\eta\|^2 =
O_p(\delta_n^2)$.
\end{theorem}

\begin{proof}[Proof sketch]
By the same three-term decomposition and interpolation-norm argument as
Theorem~\ref{thm:penalized_mean_rate}, now applied to the bivariate Fourier
coefficients $\{h_{\nu\tau}\}$ of $K_\eta-\overline{K}_\eta$ in the tensor
basis $\{e_\nu\otimes e_\tau\}$ (\cite{hsingeubank2015}, culminating result
of Section~8.3); the operator-norm bound follows from
\eqref{eq:penalized_cov_rate} via the Cauchy--Schwarz bound of Definition
\ref{def:integral_operator}, $\|\Delta_\eta\| \leq \big(\iint(K_\eta-
\gamma)^2\big)^{1/2} = \|K_\eta-\gamma\|_{\mathbb{L}^2([0,1]^2)}$.
\end{proof}

Comparing Sections~\ref{subsec:local_linear} and
\ref{subsec:penalized_ls_estimation}: at $q=2$ the local-linear and
penalized-least-squares rates essentially coincide, but the penalized
least-squares (smoothing-spline) approach requires the stronger assumption
that sample paths be differentiable (as opposed to merely bounded), while
offering better computational efficiency -- the trade-off underlying the
choice, in Section~\ref{subsec:bspline}, of an adaptive smoothing-spline
(rather than local-linear kernel) estimator for each country's GDP growth
trajectory.

\section{Functional Principal Component Analysis}
\label{sec:fpca_theory}

This closing section is the capstone of the chapter: it combines the
spectral decomposition of Section~\ref{sec:spectral_theory} (Mercer's
Theorem identifying a covariance kernel with its eigen-expansion), the
Karhunen--Lo\`eve Theorem of Section~\ref{sec:random_elements} (FPCA at the
population level), the covariance operator estimators of
Section~\ref{sec:mean_cov_estimation}, and the eigenprojection/eigenvalue/
eigenvector perturbation theory of Section~\ref{sec:perturbation}, to
establish the large-sample distribution of the \emph{estimated} FPCA
eigenvalues and eigenfunctions -- exactly the objects Chapter
\ref{chap:results} reports (Table~\ref{tab:fpca_summary}, the cumulative
variance curves of Figure~\ref{fig:cumvar}, and the FPC score plots of
Figure~\ref{fig:score_plots}). We follow \cite{hsingeubank2015},
Chapter~9, again treating the two discretely-observed-data variants
(Sections~\ref{subsec:fpca_local_linear}, \ref{subsec:fpca_pls}) with proof
sketches, since Section~\ref{subsec:fpca_sample_covariance}'s
completely-observed-data theory is what applies once each country's
trajectory has already been individually smoothed
(Section~\ref{subsec:bspline}).

Throughout, $\chi$ (or $X$) denotes the (mean-zero, WLOG) functional
observation with covariance operator $\Gamma=\sum_j\lambda_j\mathcal{P}_j$,
nonrepeated eigenvalues $\lambda_1>\lambda_2>\cdots$ of multiplicity
$r_j\geq1$ and eigenprojections $\mathcal{P}_j$ (Theorem
\ref{thm:spectral_theorem}); by the Karhunen--Lo\`eve Theorem (Corollary
\ref{cor:population_kl}/Theorem~\ref{thm:karhunen_loeve}), $\chi=\sum_j
Z_je_j$ a.s.\ for uncorrelated scores $Z_j$ with $\mathrm{Var}(Z_j)=
\lambda_j$ -- functional PCA \emph{is} the empirical version of this
decomposition.

\subsection{Estimation via the Sample Covariance Operator}
\label{subsec:fpca_sample_covariance}

Let $\chi_1,\dots,\chi_n$ be i.i.d.\ copies of $\chi$ with sample
covariance operator $\Gamma_n$ (Definition~\eqref{eq:sample_covariance}),
(possibly repeated) eigenvalues $\{\lambda_{jn}\}$, and matched
eigenprojections $\mathcal{P}_{jn}$ (Equation~\eqref{eq:matched_projection}).
Set $\Delta_n:=\Gamma_n-\Gamma$.

\begin{theorem}
\label{thm:fpca_projection_clt}
If $\mathbb{E}\|\chi\|^4<\infty$,
\begin{equation}
n^{1/2}(\mathcal{P}_{jn}-\mathcal{P}_j) \xrightarrow{d} \mathcal{S}_j
\mathfrak{Z}\mathcal{P}_j+\mathcal{P}_j\mathfrak{Z}\mathcal{S}_j,
\label{eq:fpca_projection_clt}
\end{equation}
where $\mathfrak{Z}$ is the distributional limit of $n^{1/2}(\Gamma_n-
\Gamma)$ from Theorem~\ref{thm:sample_covariance_asymptotics} and
$\mathcal{S}_j=\sum_{k\neq j}(\lambda_j-\lambda_k)^{-1}\mathcal{P}_k$ as in
Theorem~\ref{thm:projection_perturbation}.
\end{theorem}

\begin{proof}
By Theorem~\ref{thm:sample_covariance_asymptotics}, $\Delta_n\to0$ a.s., so
$\mathbb{P}(\|\Delta_n\|<\eta_j\text{ eventually})=1$ for $\eta_j=\frac12
\min_{k\neq j}|\lambda_k-\lambda_j|$; Theorem~\ref{thm:projection_perturbation}
then applies eventually a.s., giving $\mathcal{P}_{jn}-\mathcal{P}_j =
\mathcal{S}_j\Delta_n\mathcal{P}_j+\mathcal{P}_j\Delta_n\mathcal{S}_j +
O(\|\Delta_n\|^2)$. Since $\|\Delta_n\|=O_p(n^{-1/2})$ (Theorem
\ref{thm:sample_covariance_asymptotics}), $n^{1/2}\cdot O(\|\Delta_n\|^2) =
O_p(n^{-1/2}) = o_p(1)$, so $n^{1/2}(\mathcal{P}_{jn}-\mathcal{P}_j)$ is
asymptotically equivalent to $\mathcal{S}_j(n^{1/2}\Delta_n)\mathcal{P}_j+
\mathcal{P}_j(n^{1/2}\Delta_n)\mathcal{S}_j$, a continuous (bounded
bilinear) functional of $n^{1/2}\Delta_n$; the continuous mapping theorem
and $n^{1/2}\Delta_n\xrightarrow{d}\mathfrak{Z}$
(Theorem~\ref{thm:sample_covariance_asymptotics}) give
\eqref{eq:fpca_projection_clt}.
\end{proof}

\begin{theorem}
\label{thm:fpca_eigenvalue_clt}
If $\mathbb{E}\|\chi\|^4<\infty$, for any fixed $j$,
\begin{equation}
\big(n^{1/2}(\lambda_{kn}-\lambda_j)\big)_{k=N_{j-1}+1,\dots,N_j}
\xrightarrow{d} \mathrm{eig}(\mathcal{P}_j\mathfrak{Z}\mathcal{P}_j),
\label{eq:fpca_eigenvalue_clt}
\end{equation}
the (random) eigenvalues of the $r_j\times r_j$ self-adjoint operator
$\mathcal{P}_j\mathfrak{Z}\mathcal{P}_j$ (restricted to the eigenspace of
$\mathcal{P}_j$).
\end{theorem}

\begin{proof}
The $n^{1/2}(\lambda_{kn}-\lambda_j)$, $k=N_{j-1}+1,\dots,N_j$, are the
eigenvalues of $n^{1/2}(\mathcal{P}_{jn}\Gamma_n\mathcal{P}_{jn}-\lambda_j
\mathcal{P}_{jn})$. By Theorem~\ref{thm:eigenvalue_perturbation_exact}
(with $\widetilde{\mathcal{T}}=\Gamma_n$),
\[
\mathcal{P}_{jn}\Gamma_n\mathcal{P}_{jn}-\lambda_j\mathcal{P}_{jn} =
\mathcal{P}_j\Delta_n\mathcal{P}_j + (\mathcal{P}_{jn}-\mathcal{P}_j)
\Delta_n\mathcal{P}_{jn} + \mathcal{P}_j\Delta_n(\mathcal{P}_{jn}-
\mathcal{P}_j) + (\mathcal{P}_{jn}-\mathcal{P}_j)(\Gamma-\lambda_jI)
(\mathcal{P}_{jn}-\mathcal{P}_j),
\]
using $\mathcal{P}_{jn}\Delta_n\mathcal{P}_{jn}=\mathcal{P}_j\Delta_n
\mathcal{P}_j+(\mathcal{P}_{jn}-\mathcal{P}_j)\Delta_n\mathcal{P}_{jn}+
\mathcal{P}_j\Delta_n(\mathcal{P}_{jn}-\mathcal{P}_j)$ (algebraic
expansion). By Theorem~\ref{thm:fpca_projection_clt}, $\mathcal{P}_{jn}-
\mathcal{P}_j=O_p(n^{-1/2})$, and $\Delta_n=O_p(n^{-1/2})$, so each of the
last three terms is $O_p(n^{-1})$; multiplying by $n^{1/2}$, these vanish
in probability, leaving $n^{1/2}(\mathcal{P}_{jn}\Gamma_n\mathcal{P}_{jn}-
\lambda_j\mathcal{P}_{jn})$ asymptotically equivalent to $\mathcal{P}_j
(n^{1/2}\Delta_n)\mathcal{P}_j$. Since eigenvalues of a self-adjoint
operator restricted to a fixed finite-dimensional subspace are a continuous
function of the operator (Theorem~\ref{thm:eigenvalue_perturbation}, applied
within the $r_j$-dimensional range of $\mathcal{P}_j$), the continuous
mapping theorem and $n^{1/2}\Delta_n\xrightarrow{d}\mathfrak{Z}$ give
\eqref{eq:fpca_eigenvalue_clt}.
\end{proof}

The case $r_j=1$, of primary interest (and the generic case for a
covariance operator with simple spectrum), admits an explicit limit for
both the eigenvalue and the eigenfunction.

\begin{theorem}
\label{thm:fpca_simple_eigenvalue_clt}
Suppose $\mathbb{E}\|\chi\|^4<\infty$ and $\lambda_j$ has unit multiplicity
with eigenfunction $e_j$; let $(\lambda_{jn},e_{jn})$ be the matched
eigenpair of $\Gamma_n$ (with $\langle e_{jn},e_j\rangle\geq0$). Then
\begin{equation}
n^{1/2}(\lambda_{jn}-\lambda_j) \xrightarrow{d} \langle\mathfrak{Z}e_j,
e_j\rangle, \qquad
n^{1/2}(e_{jn}-e_j) \xrightarrow{d} \sum_{k\neq j}(\lambda_j-\lambda_k)^{-1}
\mathcal{P}_k\mathfrak{Z}e_j.
\label{eq:fpca_simple_clt}
\end{equation}
\end{theorem}

\begin{proof}
By Theorem~\ref{thm:eigenvalue_perturbation_exact}
(Equation~\eqref{eq:eigenvalue_perturbation_simple}) with $\Delta=\Delta_n$,
$\lambda_{jn}-\lambda_j = \langle\Delta_ne_j,e_j\rangle+O_p(\|\Delta_n\|^2)$,
and $n^{1/2}\cdot O_p(\|\Delta_n\|^2)=O_p(n^{-1/2})=o_p(1)$; the first
identity follows from $n^{1/2}\Delta_n\xrightarrow{d}\mathfrak{Z}$ and the
continuous mapping theorem applied to $\Delta\mapsto\langle\Delta e_j,
e_j\rangle$. Similarly, by Theorem~\ref{thm:eigenvector_perturbation}
(Equation~\eqref{eq:eigenvector_perturbation_bound}), $e_{jn}-e_j =
\mathcal{S}_j\Delta_ne_j+O_p(\|\Delta_n\|^2)$, and the same argument, with
continuous mapping applied to $\Delta\mapsto\mathcal{S}_j\Delta e_j$, gives
the second identity.
\end{proof}

Theorem~\ref{thm:fpca_simple_eigenvalue_clt} is the precise asymptotic
justification for treating the eigenvalues and eigenfunctions reported in
Table~\ref{tab:fpca_summary} as consistent, asymptotically Gaussian
estimates of the population FPCA of the (B-spline-smoothed) GDP growth
process: $\sqrt{n}$-consistency holds under only a fourth-moment condition
on the smoothed trajectories, with no smoothing-parameter-dependent bias
term, precisely because the input to $\Gamma_n$ is already-smoothed,
effectively fully observed functional data.

\subsection{Estimation via Local Linear Smoothing}
\label{subsec:fpca_local_linear}

We return to the discretely, noisily observed model
\eqref{eq:local_linear_model} and the local-linear covariance kernel
estimator $\gamma_h(s,t)=\sum_j\widehat\lambda_j\widehat{e}_j(s)\widehat{e}
_j(t)$ of Section~\ref{subsec:local_linear}, with $\Delta_h$ the integral
operator with kernel $\gamma_h-\gamma$ (rate given by
Theorem~\ref{thm:local_linear_operator_rate}). Assume the $\lambda_j$ are
distinct, and resolve the sign ambiguity of $\widehat{e}_j$ by choosing
$\langle\widehat{e}_j,e_j\rangle\geq0$.

Since $\|\Delta_h\|\leq\|\Delta_he\|$ is not generally true, but
$\|\mathcal{P}_j-\widehat{\mathcal{P}}_j\| = O_p(\|\Delta_h\|/\gamma_j)$
follows directly from \eqref{eq:projection_norm_bound} (with $\gamma_j:=
\frac12\min_{\lambda_k\neq\lambda_j}|\lambda_k-\lambda_j|$) whenever
$\|\Delta_h\|/\gamma_j\to0$, Theorem~\ref{thm:local_linear_operator_rate}
already gives the rate of $\widehat{\mathcal{P}}_j-\mathcal{P}_j$
(simultaneously over $j$). We now refine this to $\widehat\lambda_j$ and
$\widehat{e}_j$ individually.

\begin{theorem}
\label{thm:fpca_local_linear_eigenvalue}
Under the conditions of Theorem~\ref{thm:local_linear_operator_rate}, for
any $j$ with $\lambda_j>0$,
\begin{equation}
\widehat\lambda_j-\lambda_j = O\big((\log n/n)^{1/2}+h_m^2+h_R^2+
\delta_{n1}^2(h_m)+\delta_{n2}^2(h_R)\big) \quad\text{a.s.}
\label{eq:fpca_local_linear_eigenvalue}
\end{equation}
\end{theorem}

\begin{proof}[Proof sketch]
By Equations~\eqref{eq:eigenvalue_perturbation_simple} and
\eqref{eq:projection_norm_bound}, $\widehat\lambda_j-\lambda_j =
\iint(\gamma_h-\gamma)(s,t)e_j(s)e_j(t)\,ds\,dt + O(\|\Delta_h\|^2) =: A_{n1}
-A_{n2}+O(\|\Delta_h\|^2)$, splitting $\gamma_h=R_{h_R}-m_{h_m}\otimes
m_{h_m}$ into its two constituent local-linear estimators (Section
\ref{subsec:local_linear}). $A_{n2}$, the $m_{h_m}$-driven term, is
$O((\log n/n)^{1/2}+h_m^2)$ a.s.\ by Theorem
\ref{thm:local_linear_mean_rate} applied twice. For $A_{n1}$, substituting
the decomposition of $R_{h_R}-R$ from the (sketched) proof of the
local-linear covariance rate reduces $A_{n1}$ to an average of the form
$n^{-1}\sum_iZ_{ni}$ for bounded-moment, independent-across-$i$ $Z_{ni}$
(pairing sampling points within each curve $i$), to which the truncation
argument of Lemma~\ref{lem:fpca_truncated_slln} below applies, giving
$A_{n1}=O((\log n/n)^{1/2}+h_R^2)$ a.s. Combining the two bounds with
$\|\Delta_h\|^2=O(h_m^4+h_R^4+\delta_{n1}^2(h_m)+\delta_{n2}^2(h_R))$ a.s.\
(Theorems~\ref{thm:local_linear_operator_rate}--\ref{thm:local_linear_mean_rate}
combined, as in the source) and simplifying (since $(\log n/n)^{1/2}=
o(\delta_{n1}(h_m))$) gives \eqref{eq:fpca_local_linear_eigenvalue}.
\end{proof}

\begin{lemma}
\label{lem:fpca_truncated_slln}
Let $Z_{n1},\dots,Z_{nn}$ be independent, mean-zero, with $\sup_{i,n}
\mathbb{E}|Z_{ni}|^\delta<\infty$ and $n^{-1}\sum_i|Z_{ni}|^\delta=O(1)$
a.s., for some $\delta\in(2,\infty)$. Then $n^{-1}\sum_iZ_{ni} =
O((\log n/n)^{1/2})$ a.s.
\end{lemma}

\begin{proof}
Let $a_n=(\log n/n)^{1/2}$ and split $Z_{ni}=Z_{ni}^\succ+Z_{ni}^\prec$ for
$Z_{ni}^\succ:=Z_{ni}I(|Z_{ni}|>a_n^{-1})$, $Z_{ni}^\prec:=Z_{ni}I(|Z_{ni}|
\leq a_n^{-1})$. By H\"older's inequality, $a_n^{-1}n^{-1}\sum_i|Z_{ni}
^\succ| \leq a_n^{-1}n^{-1}\sum_i|Z_{ni}|^{1-\delta}|Z_{ni}|^\delta \leq
a_n^{1-\delta}n^{-1}\sum_i|Z_{ni}|^\delta \to0$ a.s.\ (since $a_n^{1-\delta}
\to0$ and the sum is $O(1)$ a.s.), and likewise in expectation, so
$n^{-1}\sum_i(Z_{ni}^\succ-\mathbb{E}Z_{ni}^\succ)=o(a_n)$ a.s. For the
truncated part, Bernstein's inequality gives $\mathbb{P}\big(|n^{-1}\sum_i
(Z_{ni}^\prec-\mathbb{E}Z_{ni}^\prec)|>Ba_n\big) \leq \exp\{-B^2\log n/
(2\sigma^2+(2/3)B)\}$, summable in $n$ for $B$ large; Borel--Cantelli
completes the proof.
\end{proof}

\begin{theorem}
\label{thm:fpca_local_linear_eigenfunction}
Under the conditions of Theorem~\ref{thm:local_linear_operator_rate}, for
$\lambda_j>0$,
\begin{equation}
\|\widehat{e}_j-e_j\| = O(h_m^2+h_R^2+\delta_{n1}(h_m)+\delta_{n1}(h_R))
\quad\text{a.s.}, \qquad
\sup_t|\widehat{e}_j(t)-e_j(t)| = O(h_m^2+h_R^2+\delta_{n1}(h_m)+
\delta_{n1}(h_R)+\delta_{n2}^2(h_R)) \quad\text{a.s.}
\label{eq:fpca_local_linear_eigenfunction}
\end{equation}
\end{theorem}

\begin{proof}
By Theorem~\ref{thm:eigenvector_perturbation} (Equation
\eqref{eq:eigenvector_expansion}), $\widehat{e}_j-e_j = \sum_{k\neq j}
(\lambda_j-\lambda_k)^{-1}\langle\Delta_he_j,e_k\rangle e_k + O(\|\Delta_h
\|^2)$, so by Bessel's inequality $\|\widehat{e}_j-e_j\|\leq C(\|\Delta_h
e_j\|+\|\Delta_h\|^2)$ for a finite $C$ depending on $\gamma_j$. By
Theorems~\ref{thm:local_linear_operator_rate} and (as invoked in the proof
of Theorem~\ref{thm:fpca_local_linear_eigenvalue}) the local-linear rates,
$\|\Delta_h\|^2 = O(h_m^4+h_R^4+\delta_{n1}^2(h_m)+\delta_{n2}^2(h_R))$ and
$\|\Delta_he_j\| = O(h_m^2+h_R^2+\delta_{n1}(h_m)+\delta_{n1}(h_R))$ a.s.,
giving the first bound.

For the uniform bound, for any $t$, $\widehat\lambda_j\widehat{e}_j(t)-
\lambda_je_j(t) = \int_0^1(\gamma_h(s,t)-\gamma(s,t))e_j(s)\,ds + \int_0^1
\gamma_h(s,t)(\widehat{e}_j(s)-e_j(s))\,ds =: (\Delta_he_j)(t) + R_n(t)$,
where by Cauchy--Schwarz, $|R_n(t)|\leq\sup_{s,t}|\gamma_h(s,t)|\cdot\|
\widehat{e}_j-e_j\| = O(\|\widehat{e}_j-e_j\|)$ a.s.\ (using uniform
boundedness of $\gamma_h$, itself $O(1)$ a.s.\ by consistency); so, using
the already-established rate for $\|\widehat e_j-e_j\|$,
$\widehat\lambda_j\widehat{e}_j(t)-\lambda_je_j(t) = O(h_m^2+h_R^2+
\delta_{n1}(h_m)+\delta_{n1}(h_R))$ a.s., uniformly in $t$. Writing
$\lambda_j(\widehat{e}_j(t)-e_j(t)) = \{\widehat\lambda_j\widehat{e}_j(t)-
\lambda_je_j(t)\} - (\widehat\lambda_j-\lambda_j)\widehat{e}_j(t)$ and
combining this bound with Theorem
\ref{thm:fpca_local_linear_eigenvalue}'s rate for $\widehat\lambda_j-
\lambda_j$ (and uniform boundedness of $\widehat{e}_j$) via the triangle
inequality gives the second bound (the $(\log n/n)^{1/2}$ term from
Theorem~\ref{thm:fpca_local_linear_eigenvalue} being dominated by
$\delta_{n1}(h_m)$).
\end{proof}

\begin{corollary}
\label{cor:fpca_local_linear_rates}
Under the conditions of Theorem~\ref{thm:local_linear_operator_rate}, with
$\lambda_j>0$: if $r\leq M$ (sparse) and bandwidths chosen optimally,
$\widehat\lambda_j-\lambda_j=O(\{\log n/n\}^{1/2})$ a.s., while
$\|\widehat{e}_j-e_j\|$ and $\sup_t|\widehat{e}_j(t)-e_j(t)|$ attain the
classical nonparametric rate $O(h_R^2+\{\log n/(nh_R)\}^{1/2})$ a.s.\ (at
$h_R\asymp n^{-1/5}$); if $r=r_n\to\infty$ with $r_n^{-1}\lesssim h_m,h_R
\lesssim(\log n/n)^{1/4}$ (dense), all three quantities attain the
parametric rate $O(\{\log n/n\}^{1/2})$ a.s.
\end{corollary}

So, unlike the mean and covariance \emph{kernel} estimators themselves
(Section~\ref{subsec:local_linear}), the eigenvalues of the local-linear
covariance operator are always $\sqrt{n}$-consistent, sparse or dense; only
the eigenfunctions pay the nonparametric-rate penalty in the sparse
regime -- a reflection of the operator-level rate gain already noted after
Theorem~\ref{thm:local_linear_operator_rate}.

\subsection{Estimation via Penalized Least Squares}
\label{subsec:fpca_pls}

Finally, we obtain FPCA from the smoothing-spline covariance estimator
$K_\eta$ of Theorem~\ref{thm:penalized_cov_rate}. Recall $K_\eta(s,t) =
\sum_{i=1}^nR_i(s)^T\mathcal{A}_iR_i(t)$ for symmetric, zero-diagonal
matrices $\mathcal{A}_i$ and $R_i(s):=(R_q(s,T_{i1}),\dots,R_q(s,T_{ir}))^T$,
$R_q$ the reproducing kernel of $\mathbb{W}^q[0,1]$ (Definition
\ref{def:sobolev}). Every eigenfunction of $K_\eta$ is thus of the form
$\widehat{e}(\cdot)=v^TR(\cdot)$ for $R(\cdot):=(R_1(\cdot)^T,\dots,R_n
(\cdot)^T)^T$ and some vector $v$.

\begin{theorem}
\label{thm:fpca_pls_matrix_form}
Let $\mathcal{A}:=\mathrm{diag}(\mathcal{A}_1,\dots,\mathcal{A}_n)$ and
$\mathcal{G}:=\{\int_0^1R_q(s,T_{i\nu})R_q(s,T_{j\tau})\,ds\}_{i,j;\,\nu,
\tau=1:r}$ (the Gram matrix of the $\{R_i\}$). If $u$ is an eigenvector of
$\mathcal{B}:=\mathcal{G}^{1/2}\mathcal{A}\mathcal{G}^{1/2}$ with eigenvalue
$\widehat\lambda$, then $\widehat{e}(\cdot):=u^T\mathcal{G}^{-1/2}R(\cdot)$
is an eigenfunction of $K_\eta$ with eigenvalue $\widehat\lambda$, and
distinct eigenvectors of $\mathcal{B}$ give $\mathbb{L}^2[0,1]$-orthonormal
eigenfunctions of $K_\eta$.
\end{theorem}

\begin{proof}
For eigenvectors $u_i,u_j$ of $\mathcal{B}$ with $v_i:=\mathcal{G}^{-1/2}
u_i$, $\widehat{e}_i(\cdot):=v_i^TR(\cdot)$, orthonormality holds since
$\int_0^1\widehat{e}_i(s)\widehat{e}_j(s)\,ds = v_i^T\mathcal{G}v_j =
u_i^Tu_j=\delta_{ij}$ (using the definition of $\mathcal{G}$ and
$v_i=\mathcal{G}^{-1/2}u_i$). And, writing $\{\widehat\lambda_j\}$ for the
eigenvalues of $\mathcal{B}$ with eigenvectors $\{u_j\}$, $\sum_j
\widehat\lambda_j\widehat{e}_j(s)\widehat{e}_j(t) = R(s)^T\mathcal{G}^{
-1/2}\big(\sum_j\widehat\lambda_ju_ju_j^T\big)\mathcal{G}^{-1/2}R(t) =
R(s)^T\mathcal{G}^{-1/2}\mathcal{B}\mathcal{G}^{-1/2}R(t) = R(s)^T
\mathcal{A}R(t) = K_\eta(s,t)$, identifying $\{(\widehat\lambda_j,
\widehat{e}_j)\}$ as the eigen-decomposition of $K_\eta$.
\end{proof}

We estimate $\lambda_j$ (assumed simple) by the $j$th largest eigenvalue
$\widehat\lambda_j$ of $\mathcal{B}$, and $e_j$ by $\widehat{e}_j(\cdot)=
u_j^T\mathcal{G}^{-1/2}R(\cdot)$ (sign resolved by $\langle\widehat{e}_j,
e_j\rangle\geq0$).

\begin{theorem}
\label{thm:fpca_pls_rate}
Suppose $\lambda_j$ has unit multiplicity and $\eta\asymp(\log n/(rn))^{
2q/(2q+1)}$. Under the conditions of Theorem~\ref{thm:penalized_cov_rate},
with $\delta_n^2 := (\log n/(rn))^{2q/(2q+1)}+n^{-1}$,
\begin{equation}
\widehat\lambda_j-\lambda_j = O_p(\delta_n), \qquad \|\widehat{e}_j-e_j\|^2
= O_p(\delta_n^2).
\label{eq:fpca_pls_rate}
\end{equation}
\end{theorem}

\begin{proof}
By Theorem~\ref{thm:penalized_cov_rate}, $\|\Delta_\eta\|^2=O_p(\delta_n^2)$
for $\Delta_\eta$ the operator with kernel $K_\eta-\gamma$; by
\eqref{eq:projection_norm_bound}, $\|\widehat{\mathcal{P}}_j-\mathcal{P}_j
\|=\|\widehat{e}_j\otimes\widehat{e}_j-e_j\otimes e_j\|=O_p(\delta_n)$.
By Lemma~\ref{lem:eigenvector_projection_norm}(2) and $(1-x)^{1/2}=1-x/2+
o(x)$ as $x\to0$, $\|\widehat{e}_j-e_j\|^2 = 2\big[1-(1-\|\widehat{
\mathcal{P}}_j-\mathcal{P}_j\|^2)^{1/2}\big] = \|\widehat{\mathcal{P}}_j-
\mathcal{P}_j\|^2+o_p(\delta_n^2) = O_p(\delta_n^2)$. The eigenvalue rate
follows from Equation~\eqref{eq:eigenvalue_perturbation_simple} exactly as
in the proof of Theorem~\ref{thm:fpca_simple_eigenvalue_clt}, using
$\|\Delta_\eta\|=O_p(\delta_n)$.
\end{proof}

Theorems~\ref{thm:fpca_local_linear_eigenvalue}--\ref{thm:fpca_pls_rate}
show that all three estimation schemes considered in this chapter --
sample covariance operator, local-linear smoothing, and penalized least
squares -- yield eigenvalue and eigenfunction estimators whose consistency
and (in the completely-observed case) asymptotic normality are direct,
mechanical consequences of two facts established earlier in this chapter:
operator-norm (or Hilbert--Schmidt-norm) consistency of the estimated
covariance operator for $\Gamma$ (Section~\ref{sec:mean_cov_estimation}),
and the linearization of eigenprojections, eigenvalues, and eigenfunctions
around this operator-norm perturbation (Section~\ref{sec:perturbation}).
This is the theoretical justification, in full, for treating the FPCA
eigenvalues and eigenfunctions estimated in Chapter~\ref{chap:results} --
via the sample covariance operator of the B-spline-smoothed country
trajectories -- as consistent, asymptotically normal estimates of the
population functional principal components of Latin American GDP growth.

\section{Chapter Summary and Scope for the Remaining Bases}
\label{sec:theory_summary}

This chapter has developed, from first principles, the complete rigorous
justification for the B-spline-based FPCA pipeline of
Chapter~\ref{chap:methods}: Sections~\ref{sec:preliminaries}--
\ref{sec:spectral_theory} built the Hilbert-space and operator-theoretic
machinery culminating in Mercer's Theorem, which identifies a covariance
kernel with the eigen-decomposition FPCA estimates; Section
\ref{sec:perturbation} showed that this eigen-decomposition depends
continuously, and to first order linearly, on the covariance operator
itself; Section~\ref{sec:smoothing} established that the smoothing spline
underlying the adaptive B-spline basis of Section~\ref{subsec:bspline} is
both well defined (Theorem~\ref{thm:representer}) and, via the
natural-spline optimality theorem, exactly computable in finite dimensions;
Section~\ref{sec:random_elements} gave the Karhunen--Lo\`eve Theorem, the
population-level statement of FPCA; and Sections~\ref{sec:mean_cov_estimation}
--\ref{sec:fpca_theory} combined all of the above into consistency and
asymptotic normality of the estimated FPCA eigenvalues and eigenfunctions
themselves (Theorem~\ref{thm:fpca_simple_eigenvalue_clt}), the precise
sense in which the empirical results of Chapter~\ref{chap:results} are
statistically meaningful rather than merely descriptive.

As flagged in Section~\ref{sec:theory_overview}, \cite{hsingeubank2015}
develops this pipeline for a general basis system and, concretely, for
smoothing splines (B-splines with a roughness penalty) -- matching
Section~\ref{subsec:bspline} of this thesis -- but does not treat the
Fourier or wavelet basis systems that Chapter~\ref{chap:methods} compares
against the B-spline baseline. The theory that would complete this
chapter's coverage of those two bases includes: for the \textbf{Fourier
basis}, classical approximation-theoretic results on the rate of trigonometric
approximation for functions of bounded variation, and a rigorous account of
the Gibbs phenomenon near jump discontinuities, explaining the artefacts
documented empirically in Section~\ref{subsec:gibbs_results}; for the
\textbf{wavelet basis}, the characterization of Besov spaces
$B_{p,q}^s$ as the natural regularity classes for spatially inhomogeneous
functions, the minimax theory of wavelet shrinkage of \cite{donoho1998}
already cited in Section~\ref{subsec:wavelets}, and nonlinear $n$-term
approximation theory establishing why a sparse wavelet representation can
outperform any linear (fixed-basis) method, including B-splines and
Fourier truncation, on functions with localized irregularities such as the
crisis episodes analyzed in Chapter~\ref{chap:results}. This material will
be incorporated from a dedicated approximation-theory source in a
subsequent revision of this chapter.

%---------------------------------------------------------
